\paragraph{局部化相位采样、solenoid 刚性与分层 prime-register 的极小性}
以下条目把有限素数局部化数系的相位采样闭式、基本扩张的不分裂性、solenoid 端同态环以及有限指数子群/有限商/profinite 完成/连通 Lie 商的刚性分类、dyadic 读出的奇素失明、与有限生成离散交换群的严格不可混同性，以及分层 prime-register 外置在最坏情形下的锐最优性，压缩为同一条可直接调用的终端链条。

\begin{theorem}[局部化数系相位采样函数的闭式与素数剥离律]\label{thm:xi-localized-phase-sampling-closed-form}
\leanverified{paper\_xi\_localized\_phase\_sampling\_closed\_form}
设 \(S\subset\mathbb P\) 为有限素数集，并记
\[
G_S:=\ZZ[S^{-1}]
\]
为加法群。定义相位采样函数
\[
\Sigma^{\mathrm{ph}}_{G_S}(N):=\card{\operatorname{Hom}(G_S,\ZZ/N\ZZ)}
\qquad(N\ge 1).
\]
则对一切 \(N\ge 1\) 都有闭式
\[
\boxed{\;
\Sigma^{\mathrm{ph}}_{G_S}(N)
=
\prod_{\substack{q^{v_q(N)}\parallel N\\ q\notin S}}q^{v_q(N)}
=
\frac{N}{\prod_{p\in S}p^{v_p(N)}}.\;}
\]
等价地，\(\Sigma^{\mathrm{ph}}_{G_S}(N)\) 恰由 \(N\) 去掉全部 \(S\)-主部分得到。特别地，
\[
p\in S \iff \Sigma^{\mathrm{ph}}_{G_S}(p)=1,
\qquad
p\notin S \iff \Sigma^{\mathrm{ph}}_{G_S}(p)=p.
\]
因此 \(\Sigma^{\mathrm{ph}}_{G_S}\) 唯一恢复素数集 \(S\)，并因定理 \ref{thm:killo-localization-circle-dimension-prime-ledger-splitting} 进一步唯一恢复 Pontryagin 对偶短正合列
\[
0\to \prod_{p\in S}\ZZ_p\to \widehat{G_S}\to \TT\to 0
\]
中的 profinite 核；与此同时 \(\cdim_{\mathrm{fr}}(G_S)=1\) 与 \(S\) 无关。
\end{theorem}
\begin{proof}
定理 \ref{thm:xi-cdim-localization-zeta-AS} 已证明：若把同一加法群记作
\(
A_S=\ZZ[S^{-1}]
\)，则
\[
\card{\operatorname{Hom}(A_S,\ZZ/N\ZZ)}
=
\prod_{\substack{q^{v_q(N)}\parallel N\\ q\notin S}}q^{v_q(N)}
=
\frac{N}{\prod_{p\in S}p^{v_p(N)}}
\qquad(N\ge 1).
\]
将 \(A_S\) 记为 \(G_S\) 即得所述闭式。
\par
对任意素数 \(p\)，代入 \(N=p\) 得
\[
\Sigma^{\mathrm{ph}}_{G_S}(p)=
\begin{cases}
1,& p\in S,\\[1mm]
p,& p\notin S.
\end{cases}
\]
故 \(S\) 由 \(\Sigma^{\mathrm{ph}}_{G_S}\) 唯一恢复。
再由定理 \ref{thm:killo-localization-circle-dimension-prime-ledger-splitting}，核
\(\prod_{p\in S}\ZZ_p\) 的非平凡 \(p\)-主因子与 \(p\in S\) 等价，而 \(\cdim_{\mathrm{fr}}(G_S)=1\) 始终成立。证毕。
\end{proof}

\begin{corollary}[dyadic 相位谱对奇素支撑的完全失明]\label{cor:xi-localized-dyadic-phase-blindness}
\leanverified{paper\_xi\_localized\_dyadic\_phase\_blindness}
沿用定理 \ref{thm:xi-localized-phase-sampling-closed-form} 的记号，定义 dyadic 相位谱
\[
\Sigma^{(2)}_{G_S}(a):=\Sigma^{\mathrm{ph}}_{G_S}(2^a)
\qquad(a\ge 1).
\]
则有严格二分
\[
\boxed{\;
\Sigma^{(2)}_{G_S}(a)=
\begin{cases}
2^a,& 2\notin S,\\[1mm]
1,& 2\in S.
\end{cases}\;}
\]
因此若有限素数集 \(S,T\subset\mathbb P\) 满足
\[
2\in S \iff 2\in T,
\]
则
\[
\Sigma^{(2)}_{G_S}(a)=\Sigma^{(2)}_{G_T}(a)
\qquad(a\ge 1).
\]
换言之，dyadic 相位谱只能探测 \(2\) 是否进入素数账本，而对全部奇素支撑严格失明。
\end{corollary}
\begin{proof}
把 \(N=2^a\) 代入定理 \ref{thm:xi-localized-phase-sampling-closed-form} 即得
\[
\Sigma^{\mathrm{ph}}_{G_S}(2^a)
=
\frac{2^a}{\prod_{p\in S}p^{v_p(2^a)}}
=
\begin{cases}
2^a,& 2\notin S,\\[1mm]
1,& 2\in S.
\end{cases}
\]
而对任意奇素 \(p\) 恒有 \(v_p(2^a)=0\)，故其是否属于 \(S\) 不会改变数值。证毕。
\end{proof}

\begin{theorem}[局部化数系对偶短正合列的分裂判据]\label{thm:xi-localized-basic-extension-strictly-nonsplit}
设 \(S\subset\mathbb P\) 为有限素数集，并记
\[
G_S:=\ZZ[S^{-1}].
\]
再记
\[
\Sigma_S:=\widehat{G_S}.
\]
前述结果给出互为对偶的短正合列
\[
0\longrightarrow \ZZ\longrightarrow G_S\longrightarrow \bigoplus_{p\in S}\ZZ(p^\infty)\longrightarrow 0
\]
\[
0\longrightarrow \prod_{p\in S}\ZZ_p\longrightarrow \widehat{G_S}\longrightarrow \TT\longrightarrow 0.
\]
则下列两条等价：
\begin{enumerate}
  \item \(S=\varnothing\)；
  \item 上述两条短正合列在各自范畴中都分裂。
\end{enumerate}
等价地，
\[
S\neq\varnothing
\quad\Longrightarrow\quad
\Sigma_S\not\cong \TT\times \prod_{p\in S}\ZZ_p.
\]
特别地，当 \(S\neq\varnothing\) 时，不存在群同态截面
\[
\bigoplus_{p\in S}\ZZ(p^\infty)\longrightarrow G_S,
\]
也不存在连续群同态截面
\[
\TT\longrightarrow \widehat{G_S}.
\]
\end{theorem}
\begin{proof}
若 \(S=\varnothing\)，则 \(G_S=\ZZ\) 且 \(\Sigma_S=\widehat{\ZZ}\cong \TT\)，两条短正合列都退化为平凡扩张，故显然分裂。

现设 \(S\neq\varnothing\)。若第一条短正合列分裂，则
\[
G_S\cong \ZZ\oplus \bigoplus_{p\in S}\ZZ(p^\infty).
\]
右端含有非零挠子群，而
\[
G_S=\ZZ[S^{-1}]\subset \QQ
\]
作为有理数加法群的子群无挠，矛盾。故离散侧不分裂。

Pontryagin 对偶在局部紧交换群范畴中给出精确反变等价。若第二条短正合列分裂，则再对偶化便得到第一条离散短正合列分裂，与上一步矛盾。故紧侧亦不分裂。证毕。
\end{proof}

\begin{theorem}[局部化数系端同态环的完全分类与 solenoid 的直积不可分解性]\label{thm:xi-localized-endomorphism-ring-solenoid-indecomposable}
沿用定理 \ref{thm:xi-localized-basic-extension-strictly-nonsplit} 的记号。则：
\begin{enumerate}
  \item 任意群同态 \(\phi:G_S\to G_S\) 都唯一地写成
  \[
  \phi(x)=qx
  \qquad (x\in G_S)
  \]
  的形式，其中 \(q\in G_S\)。
  \item 因而
  \[
  \operatorname{End}(G_S)\cong G_S
  \]
  作为交换环同构成立，其中加法为点态相加，乘法为复合。
  \item \(\operatorname{End}(G_S)\) 的幂等元只有 \(0,1\)。等价地，
  \[
  \operatorname{End}_{\mathrm{cts}}(\widehat{G_S})
  \cong
  \operatorname{End}(G_S)^{\mathrm{op}}
  \cong
  G_S
  \]
  亦只含平凡幂等元；特别地，\(\widehat{G_S}\) 不存在非平凡连续幂等自同态，因而不能分解为两个非平凡紧交换群的直积。
\end{enumerate}
\end{theorem}
\begin{proof}
取任意 \(\phi\in \operatorname{End}(G_S)\)，记
\[
q:=\phi(1)\in G_S.
\]
对任意 \(x\in G_S\)，可写成
\[
x=\frac{a}{D},
\qquad
a\in \ZZ,
\qquad
D=\prod_{p\in S}p^{k_p}
\]
的形式。由 \(Dx=a\) 可得
\[
D\,\phi(x)=\phi(a)=aq.
\]
由于 \(G_S\subset \QQ\) 无挠，乘以 \(D\) 的映射单射，故
\[
\phi(x)=\frac{aq}{D}=qx.
\]
这证明了第 \(1\) 条；唯一性由在 \(x=1\) 处取值立即得到。反过来，任取 \(q\in G_S\)，映射 \(x\mapsto qx\) 因 \(G_S\) 为 \(\QQ\) 的子环而确为 \(G_S\) 的自同态。于是
\[
q\longmapsto m_q,\qquad m_q(x):=qx
\]
给出 \(\operatorname{End}(G_S)\) 与 \(G_S\) 之间的双射。

该双射保持加法，且满足
\[
m_q\circ m_r=m_{qr},
\]
故第 \(2\) 条成立。

由于 \(G_S=\ZZ[S^{-1}]\subset \QQ\) 是整环的子环，幂等方程
\[
e^2=e
\]
只可能有解 \(e=0\) 或 \(e=1\)。故 \(\operatorname{End}(G_S)\) 仅有平凡幂等元。Pontryagin 对偶把连续端同态环识别为离散端端同态环的反环；又因 \(G_S\) 交换，反环与自身一致，故 \(\operatorname{End}_{\mathrm{cts}}(\widehat{G_S})\) 同样只含平凡幂等元。

若 \(\widehat{G_S}\cong K_1\times K_2\) 为两个非平凡紧交换群的直积，则合成
\[
\widehat{G_S}\xrightarrow{\ \mathrm{pr}_1\ }K_1
\hookrightarrow
K_1\times K_2\cong \widehat{G_S}
\]
给出一个既非零也非恒等的连续幂等自同态，矛盾。故 \(\widehat{G_S}\) 不能作非平凡直积分解。证毕。
\end{proof}

\begin{theorem}[有限商与有限相位层的完全分类]\label{thm:xi-localized-finite-quotient-phase-layer-classification}
\leanverified{paper\_xi\_localized\_finite\_quotient\_phase\_layer\_classification}
沿用定理 \ref{thm:xi-localized-basic-extension-strictly-nonsplit} 的记号。对任意正整数 \(n\)，定义
\[
n_{S^\perp}:=\frac{n}{\prod_{p\in S}p^{v_p(n)}},
\]
即从 \(n\) 中剥离全部 \(S\)-主因子后得到的 \(S\)-自由部分。则：
\begin{enumerate}
  \item 商群满足
  \[
  G_S/nG_S\cong \ZZ/n_{S^\perp}\ZZ.
  \]
  \item 因而 \(G_S\) 的全部有限商恰为循环群 \(\ZZ/d\ZZ\)，其中 \(d\) 的全部素因子都不属于 \(S\)。
  \item 对偶地，\(\widehat{G_S}\) 的全部有限闭子群也恰为这些循环群。
\end{enumerate}
\end{theorem}
\begin{proof}
把 \(n\) 分解为
\[
n=n_S\,n_{S^\perp},
\qquad
n_S:=\prod_{p\in S}p^{v_p(n)}.
\]
对每个 \(p\in S\)，乘法映射 \(x\mapsto px\) 在 \(G_S=\ZZ[S^{-1}]\) 上是自同构，因为 \(x\mapsto x/p\) 仍取值于 \(G_S\)。故
\[
nG_S=n_{S^\perp}G_S.
\]
因此只需证明当 \(n\) 与 \(\prod_{p\in S}p\) 互素时，
\[
G_S/nG_S\cong \ZZ/n\ZZ.
\]

考虑自然映射
\[
\rho:\ZZ\longrightarrow G_S/nG_S.
\]
任取 \(x=a/b\in G_S\)，其中 \(a\in\ZZ\)，且 \(b\) 的全部素因子属于 \(S\)。由于 \(\gcd(b,n)=1\)，可取 \(c,k\in\ZZ\) 使
\[
bc=1+nk.
\]
于是
\[
x-ac
=
\frac{a}{b}-ac
=
\frac{a(1-bc)}{b}
=
-n\,\frac{ak}{b}\in nG_S,
\]
故 \(x\equiv ac\pmod{nG_S}\)。这说明 \(\rho\) 满射。

再看核。若 \(m\in\ker\rho\)，则
\[
m=n\frac{a}{b}
\]
对某个 \(\frac{a}{b}\in G_S\) 成立，且 \(\gcd(b,n)=1\)。于是
\[
bm=na.
\]
由于 \(\gcd(b,n)=1\)，必有 \(b\mid a\)，从而 \(a/b\in \ZZ\)，进而 \(m\in n\ZZ\)。故
\[
\ker(\rho)=n\ZZ,
\]
于是
\[
G_S/nG_S\cong \ZZ/n\ZZ.
\]
把 \(n\) 替换回 \(n_{S^\perp}\) 即得第 \(1\) 条。

若 \(F\) 为 \(G_S\) 的任一有限商，令 \(n\) 为 \(F\) 的指数，则满射 \(G_S\twoheadrightarrow F\) 经由 \(G_S/nG_S\) 因子化。由第 \(1\) 条，\(G_S/nG_S\) 为循环群，故 \(F\) 亦为循环群，且其阶的全部素因子都不属于 \(S\)。反之，若 \(d\) 的全部素因子都不属于 \(S\)，则 \(d=d_{S^\perp}\)，从而
\[
G_S/dG_S\cong \ZZ/d\ZZ.
\]
第 \(2\) 条得证。

第 \(3\) 条由 Pontryagin 对偶把有限离散商与有限闭子群一一对应，并且有限循环群的对偶仍为同阶循环群，立即得到。证毕。
\end{proof}

\paragraph{solenoid 覆盖度谱、乘法周期律与有限闭子群的循环实现}
在局部化数系的端同态环、有限商与有限闭子群已经完成分类之后，还可把整数乘法在 \(\widehat{G_S}\) 上进一步压缩为一组完全显式的覆盖与动力学结论。其核心现象是：所有有限层数据都只由 \(S\)-自由部分
\[
n_{S^\perp}:=\frac{n}{\prod_{p\in S}p^{v_p(n)}}
\]
控制。

\begin{theorem}[整数乘法在 \texorpdfstring{$\widehat{G_S}$}{G-hat-S} 上的核度公式]\label{thm:xi-localized-solenoid-multiplication-kernel-degree}
设 \(S\subset\mathbb P\) 为有限非空素数集，记
\[
G_S:=\ZZ[S^{-1}],
\qquad
\Sigma_S:=\widehat{G_S}.
\]
对任意正整数 \(n\)，考虑整数乘法自同态
\[
[n]:\Sigma_S\longrightarrow \Sigma_S,
\qquad
x\longmapsto nx.
\]
则 \([n]\) 必为满射，其核为有限循环群，并且
\[
\boxed{\;
\ker[n]\cong \ZZ/n_{S^\perp}\ZZ,
\qquad
\deg([n])=n_{S^\perp}.\;}
\]
\end{theorem}
\begin{proof}
令
\[
m_n:G_S\to G_S,
\qquad
x\mapsto nx
\]
为 Pontryagin 对偶端映射。由于
\[
G_S=\ZZ[S^{-1}]\subset \QQ
\]
无挠，\(m_n\) 为单射。于是对偶映射 \([n]\) 为满射。

再看短正合列
\[
0\longrightarrow G_S\xrightarrow{\,m_n\,}G_S\longrightarrow G_S/nG_S\longrightarrow 0.
\]
Pontryagin 对偶的精确性给出
\[
0\longrightarrow \widehat{G_S/nG_S}\longrightarrow \Sigma_S\xrightarrow{\,[n]\,}\Sigma_S\longrightarrow 0,
\]
故
\[
\ker[n]\cong \widehat{G_S/nG_S}.
\]
而定理 \ref{thm:xi-localized-finite-quotient-phase-layer-classification} 已证明
\[
G_S/nG_S\cong \ZZ/n_{S^\perp}\ZZ.
\]
有限循环群的 Pontryagin 对偶仍为同阶循环群，因此
\[
\ker[n]\cong \ZZ/n_{S^\perp}\ZZ.
\]

最后求覆盖度。记 \(F:=\ker[n]\)。由于 \([n]\) 为满射群同态，诱导同构
\[
\Sigma_S/F\cong \Sigma_S.
\]
有限群 \(F\) 通过平移自由作用于 \(\Sigma_S\)。取单位元的一个开邻域 \(U\) 使
\[
(U-U)\cap F=\{0\},
\]
则不同 \(F\)-平移 \(U+\eta\) 两两不交。故商映射
\[
q:\Sigma_S\to \Sigma_S/F
\]
是度数为 \(|F|\) 的有限正则覆盖，而 \([n]\) 与该商映射只差一个群同构。于是
\[
\deg([n])=|F|=n_{S^\perp}.
\]
证毕。
\end{proof}

\begin{theorem}[局部化标量乘法在 \texorpdfstring{$\widehat{G_S}$}{G-hat-S} 上的核度公式]\label{thm:xi-localized-solenoid-localized-scalar-kernel-degree}
\leanverified{paper\_xi\_localized\_solenoid\_localized\_scalar\_kernel\_degree}
设 \(S\subset\mathbb P\) 为有限素数集，记
\[
G_S:=\ZZ[S^{-1}],
\qquad
\Sigma_S:=\widehat{G_S}.
\]
任取非零标量
\[
q\in \ZZ[S^{-1}],
\qquad
q=\varepsilon\,a\prod_{p\in S}p^{-k_p},
\qquad
\varepsilon\in\{\pm1\},
\quad
a\in\NN,
\]
并定义其 \(S\)-自由部分
\[
a_{S^\perp}:=\frac{a}{\prod_{p\in S}p^{v_p(a)}}.
\]
令
\[
[q]:\Sigma_S\longrightarrow \Sigma_S
\]
表示与离散侧“乘以 \(q\)”对偶的连续自同态。则 \([q]\) 为满射，其核为有限循环群，并且
\[
\boxed{\;
\ker[q]\cong \ZZ/a_{S^\perp}\ZZ,
\qquad
\deg([q])=a_{S^\perp}.\;}
\]
特别地，
\[
[q]\ \text{是同构}
\iff
a_{S^\perp}=1
\iff
q\in \ZZ[S^{-1}]^\times.
\]
\end{theorem}
\begin{proof}
令
\[
m_q:G_S\to G_S,
\qquad
x\mapsto qx
\]
为离散侧端映射。由于 \(G_S\subset \QQ\) 无挠且 \(q\neq 0\)，故 \(m_q\) 为单射，从而其对偶 \([q]\) 为满射。

把 \(a\) 分解为
\[
a=a_S\,a_{S^\perp},
\qquad
a_S:=\prod_{p\in S}p^{v_p(a)}.
\]
则
\[
u:=\varepsilon\,a_S\prod_{p\in S}p^{-k_p}
\]
是局部化环 \(\ZZ[S^{-1}]\) 的单位，且
\[
q=u\,a_{S^\perp}.
\]
于是
\[
qG_S=a_{S^\perp}G_S.
\]
对短正合列
\[
0\longrightarrow G_S\xrightarrow{\,m_q\,}G_S\longrightarrow G_S/qG_S\longrightarrow 0
\]
作 Pontryagin 对偶，得到
\[
0\longrightarrow \widehat{G_S/qG_S}\longrightarrow \Sigma_S\xrightarrow{\,[q]\,}\Sigma_S\longrightarrow 0,
\]
故
\[
\ker[q]\cong \widehat{G_S/qG_S}\cong \widehat{G_S/a_{S^\perp}G_S}.
\]
而定理 \ref{thm:xi-localized-finite-quotient-phase-layer-classification} 已给出
\[
G_S/a_{S^\perp}G_S\cong \ZZ/a_{S^\perp}\ZZ.
\]
有限循环群的 Pontryagin 对偶仍为同阶循环群，遂得
\[
\ker[q]\cong \ZZ/a_{S^\perp}\ZZ.
\]

最后，\([q]\) 是紧交换群上的满同态且核有限，故如定理 \ref{thm:xi-localized-solenoid-multiplication-kernel-degree} 的证明所示，它是度数等于核大小的有限正则覆盖。因此
\[
\deg([q])=\#\ker[q]=a_{S^\perp}.
\]
核平凡当且仅当 \(a_{S^\perp}=1\)，这又等价于 \(q\) 没有任何逃离 \(S\) 的素因子，即 \(q\in \ZZ[S^{-1}]^\times\)。证毕。
\leanverified{paper\_xi\_localized\_solenoid\_multiplication\_kernel\_degree}
\end{proof}

\begin{corollary}[覆盖度谱完全由被剥离后的素数支撑决定]\label{cor:xi-localized-solenoid-cover-degree-spectrum}
\leanverified{paper\_xi\_localized\_solenoid\_cover\_degree\_spectrum}
沿用定理 \ref{thm:xi-localized-solenoid-multiplication-kernel-degree} 的记号。则 \(\Sigma_S\) 上全部整数乘法覆盖的度数集合恰为
\[
\boxed{\;
\operatorname{Deg}(\Sigma_S)
:=
\{\deg([n]): n\ge 1\}
=
\{d\in\NN:\ \text{\(d\) 的全部素因子都不属于 }S\}.\;}
\]
并且对任意 \(n\ge 1\)，以下三条等价：
\[
[n]\ \text{是同构}
\iff
n_{S^\perp}=1
\iff
\text{\(n\) 的全部素因子都属于 }S.
\]
\end{corollary}
\begin{proof}
由定理 \ref{thm:xi-localized-solenoid-multiplication-kernel-degree}，
\[
\deg([n])=n_{S^\perp}.
\]
故任意整数乘法所产生的覆盖度都自动剥离了全部 \(S\)-主因子。反之，若 \(d\) 的全部素因子都不属于 \(S\)，则
\[
d_{S^\perp}=d,
\]
取 \(n=d\) 即得 \(\deg([d])=d\)。这就给出度数谱的精确描述。

又因为 \([n]\) 已知为满射，所以它为同构当且仅当核平凡；而定理
\ref{thm:xi-localized-solenoid-multiplication-kernel-degree}
给出
\[
|\ker[n]|=n_{S^\perp}.
\]
故 \([n]\) 为同构当且仅当 \(n_{S^\perp}=1\)，也即 \(n\) 的全部素因子都落在 \(S\) 中。证毕。
\end{proof}

\begin{theorem}[整数乘法动力学的周期点计数与 primitive 轨道公式]\label{thm:xi-localized-solenoid-periodic-point-formula}
固定整数 \(a\ge 2\)，考虑连续群自映射
\[
T_a:\Sigma_S\longrightarrow \Sigma_S,
\qquad
x\longmapsto ax.
\]
则对任意 \(m\ge 1\)，其 \(m\)-步不动点群满足
\[
\boxed{\;
\mathrm{Fix}(T_a^m)\cong \ZZ/(a^m-1)_{S^\perp}\ZZ,
\qquad
\#\mathrm{Fix}(T_a^m)=(a^m-1)_{S^\perp}.\;}
\]
于是 exact period 为 \(m\) 的 primitive 轨道数满足
\[
\boxed{\;
\mathcal O_m(T_a)
=
\frac1m\sum_{d\mid m}\mu(d)\,(a^{m/d}-1)_{S^\perp}.\;}
\]
\leanverified{paper\_xi\_localized\_solenoid\_periodic\_point\_formula}
\end{theorem}
\begin{proof}
由定义，
\[
T_a^m=[a^m].
\]
故
\[
x\in \mathrm{Fix}(T_a^m)
\iff
a^m x=x
\iff
(a^m-1)x=0
\iff
x\in \ker[a^m-1].
\]
因此
\[
\mathrm{Fix}(T_a^m)=\ker[a^m-1].
\]
将定理 \ref{thm:xi-localized-solenoid-multiplication-kernel-degree} 应用于整数
\[
n=a^m-1
\]
即得
\[
\ker[a^m-1]\cong \ZZ/(a^m-1)_{S^\perp}\ZZ,
\]
并且
\[
\#\mathrm{Fix}(T_a^m)=|\ker[a^m-1]|=(a^m-1)_{S^\perp}.
\]

对第二条，设 \(\mathcal O_d(T_a)\) 为 exact period 为 \(d\) 的 primitive 轨道数。则标准周期分解给出
\[
\#\mathrm{Fix}(T_a^m)=\sum_{d\mid m} d\,\mathcal O_d(T_a).
\]
对该等式施行 Möbius 反演，得到
\[
\mathcal O_m(T_a)
=
\frac1m\sum_{d\mid m}\mu(d)\,\#\mathrm{Fix}(T_a^{m/d})
=
\frac1m\sum_{d\mid m}\mu(d)\,(a^{m/d}-1)_{S^\perp}.
\]
证毕。
\end{proof}

\begin{corollary}[单质数核实验即可完全恢复 prime-register]\label{cor:xi-localized-solenoid-prime-kernel-tomography}
\leanverified{paper\_xi\_localized\_solenoid\_prime\_kernel\_tomography}
对任意素数 \(q\)，有
\[
\#\ker[q]=
\begin{cases}
1,& q\in S,\\[1mm]
q,& q\notin S.
\end{cases}
\]
等价地，
\[
\boxed{\;
q\in S
\iff
\#\ker[q]=1.\;}
\]
因此，\(\Sigma_S\) 的一阶覆盖核谱
\[
q\longmapsto \#\ker[q]
\qquad (q\ \text{为素数})
\]
已经足以唯一恢复素数账本 \(S\)。
\end{corollary}
\begin{proof}
把定理 \ref{thm:xi-localized-solenoid-multiplication-kernel-degree} 中的 \(n\) 取为素数 \(q\)。若 \(q\in S\)，则
\[
q_{S^\perp}=1,
\]
故 \(\ker[q]\) 平凡；若 \(q\notin S\)，则
\[
q_{S^\perp}=q,
\]
故
\[
\ker[q]\cong \ZZ/q\ZZ.
\]
于是所述判别式立即成立。证毕。
\end{proof}

\begin{theorem}[\texorpdfstring{$n$}{n}-挠子群与全部有限闭子群由整数乘法核唯一实现]\label{thm:xi-localized-solenoid-finite-subgroups-are-kernels}
\leanverified{paper\_xi\_localized\_solenoid\_finite\_subgroups\_are\_kernels}
对任意 \(n\ge 1\)，\(\Sigma_S=\widehat{G_S}\) 的 \(n\)-挠子群满足
\[
\Sigma_S[n]=\ker[n]\cong \ZZ/n_{S^\perp}\ZZ.
\]
更一般地，\(\Sigma_S\) 的每一个有限闭子群 \(F\le \Sigma_S\) 都必为循环群；更精确地，存在唯一满足“其全部素因子都不属于 \(S\)”的正整数 \(d\)，使得
\[
\boxed{\;
F=\ker[d]\cong \ZZ/d\ZZ.\;}
\]
反之，对每个这样的 \(d\)，\(\ker[d]\) 都给出一个有限闭子群。特别地，\(\Sigma_S\) 中不存在任何非循环有限闭子群。
\end{theorem}
\begin{proof}
按定义，
\[
\Sigma_S[n]=\{x\in \Sigma_S:\ nx=0\}=\ker[n].
\]
故第一条由定理 \ref{thm:xi-localized-solenoid-multiplication-kernel-degree} 立即得到。

任取有限闭子群 \(F\le \Sigma_S\)。由定理
\ref{thm:xi-localized-finite-quotient-phase-layer-classification}
第 \(3\) 条，\(F\) 同构于某个循环群 \(\ZZ/d\ZZ\)，其中 \(d\) 的全部素因子都不属于 \(S\)。因此 \(F\) 为阶 \(d\) 的循环群，尤其其每个元素都被 \(d\) 杀掉，故
\[
F\subseteq \ker[d].
\]
另一方面，由定理 \ref{thm:xi-localized-solenoid-multiplication-kernel-degree}，
\[
\ker[d]\cong \ZZ/d_{S^\perp}\ZZ.
\]
由于 \(d\) 本身已是 \(S\)-自由整数，故 \(d_{S^\perp}=d\)。于是
\[
|\ker[d]|=d=|F|.
\]
结合 \(F\subseteq \ker[d]\)，可得
\[
F=\ker[d].
\]
这就证明了存在性。

唯一性也随之成立：若
\[
\ker[d_1]=\ker[d_2],
\]
则两边阶数相同。再由定理 \ref{thm:xi-localized-solenoid-multiplication-kernel-degree}，
\[
d_1=|\ker[d_1]|=|\ker[d_2]|=d_2.
\]

反过来，若 \(d\) 的全部素因子都不属于 \(S\)，则 \(d_{S^\perp}=d\)，故定理
\ref{thm:xi-localized-solenoid-multiplication-kernel-degree}
给出
\[
\ker[d]\cong \ZZ/d\ZZ,
\]
从而它确为一个有限闭子群。

最后，由所有有限闭子群都等于某个 \(\ker[d]\) 且 \(\ker[d]\) 总为循环群，立即推出 \(\Sigma_S\) 中不存在任何非循环有限闭子群。证毕。
\end{proof}

若将整数乘法进一步推广为由任意 \(x\in G_S\setminus\{0\}\) 所诱导的连续自同态，则同样的 \(S\)-自由剥离律仍然保持，只是控制有限核与周期点计数的整数由 \(n\) 变为 \(x\) 的 \(S\)-外分子以及差分 \(a^n-b^n\) 的 \(S\)-外部分。

\begin{theorem}[连续自同态的核与余核由 \texorpdfstring{$S$}{S}-外分子唯一决定]\label{thm:xi-localized-solenoid-endomorphism-kernel-cokernel-sfree-numerator}
设
\[
x=\frac{a}{b}\in G_S=\ZZ[S^{-1}]\setminus\{0\},
\qquad
a\in\ZZ\setminus\{0\},
\qquad
b\in\langle S\rangle,
\qquad
\gcd(a,b)=1.
\]
定义
\[
v:=|a|_{S^\perp}
=
\frac{|a|}{\prod_{p\in S}p^{v_p(a)}}.
\]
记
\[
m_x:G_S\to G_S,
\qquad
y\mapsto xy,
\]
并令
\[
\alpha_x:=\widehat{m_x}:\Sigma_S\to \Sigma_S
\]
为其 Pontryagin 对偶连续自同态。则：
\begin{enumerate}
  \item
  \[
  \operatorname{coker}(m_x)\cong \ZZ/v\ZZ;
  \]
  \item
  \[
  \ker(\alpha_x)\cong \ZZ/v\ZZ;
  \]
  \item \(\alpha_x\) 为同构当且仅当 \(v=1\)，亦即 \(x\in G_S^\times\)。
\end{enumerate}
\end{theorem}
\begin{proof}
写
\[
|a|=uv,
\]
其中 \(u\) 的全部素因子都属于 \(S\)，而 \(v\) 与 \(\prod_{p\in S}p\) 互素。再记 \(\varepsilon:=\operatorname{sgn}(a)\in\{\pm1\}\)。于是
\[
x=\left(\varepsilon\frac{u}{b}\right)v.
\]
由于 \(u/b\in \ZZ[S^{-1}]^\times\) 且 \(\varepsilon=\pm1\) 亦为单位，故
\[
\varepsilon\frac{u}{b}\in G_S^\times.
\]
从而 \(m_x\) 与 \(m_v\) 只差一个群自同构，故
\[
\operatorname{coker}(m_x)\cong \operatorname{coker}(m_v)\cong G_S/vG_S.
\]
又因为 \(v\) 与 \(\prod_{p\in S}p\) 互素，定理 \ref{thm:xi-localized-finite-quotient-phase-layer-classification} 给出
\[
G_S/vG_S\cong \ZZ/v\ZZ,
\]
从而得到第 \(1\) 条。

另一方面，\(x\neq 0\) 且 \(G_S\subset \QQ\) 无挠，故 \(m_x\) 为单射。对短正合列
\[
0\longrightarrow G_S\xrightarrow{\,m_x\,}G_S\longrightarrow \operatorname{coker}(m_x)\longrightarrow 0
\]
施行 Pontryagin 对偶，得到
\[
0\longrightarrow \widehat{\operatorname{coker}(m_x)}\longrightarrow \Sigma_S\xrightarrow{\,\alpha_x\,}\Sigma_S\longrightarrow 0.
\]
因此
\[
\ker(\alpha_x)\cong \widehat{\operatorname{coker}(m_x)}.
\]
而有限循环群的 Pontryagin 对偶仍为同阶循环群，故第 \(2\) 条成立。

最后，由上式可知 \(\alpha_x\) 总为满射；因此 \(\alpha_x\) 为同构当且仅当 \(\ker(\alpha_x)=0\)，等价于 \(v=1\)。而 \(v=1\) 恰当且仅当 \(a\) 的全部素因子都属于 \(S\)，也即 \(x\) 在局部化环 \(G_S\) 中可逆。证毕。
\end{proof}

\leanverified{paper\_xi\_localized\_solenoid\_endomorphism\_kernel\_cokernel\_sfree\_numerator}

\paragraph{一般连续自同态的有限覆盖度与覆盖半群的算术谱}
在一般连续自同态的核与余核已经由 \(S\)-外分子完全刻画之后，还可把其拓扑非可逆性进一步压缩为覆盖度层面的两条闭式结论。

\leanverified{paper\_xi\_localized\_solenoid\_nonzero\_endomorphism\_finite\_cover\_degree}
\begin{corollary}[非零连续自同态必为有限覆盖，且覆盖度等于 \texorpdfstring{$S$}{S}-外分子]\label{cor:xi-localized-solenoid-nonzero-endomorphism-finite-cover-degree}
沿用定理 \ref{thm:xi-localized-solenoid-endomorphism-kernel-cokernel-sfree-numerator} 的记号。则 \(\alpha_x\) 是满射有限覆盖，并且
\[
\boxed{\;
\ker(\alpha_x)\cong \ZZ/v\ZZ,
\qquad
\deg(\alpha_x)=v.\;}
\]
特别地，\(\alpha_x\) 为同构当且仅当 \(v=1\)，亦即 \(x\in G_S^\times\)。
\end{corollary}
\begin{proof}
定理 \ref{thm:xi-localized-solenoid-endomorphism-kernel-cokernel-sfree-numerator} 已给出
\[
\ker(\alpha_x)\cong \ZZ/v\ZZ
\]
以及 \(\alpha_x\) 为同构当且仅当 \(v=1\)。
另一方面，由 \(x\neq 0\) 与 \(G_S\subset \QQ\) 无挠，离散端乘法映射
\[
m_x:G_S\to G_S
\]
是单射，故其对偶 \(\alpha_x\) 为满射。

再证有限覆盖性。由于 \(\ker(\alpha_x)\) 有限，可取单位元的开邻域 \(U\subset \Sigma_S\) 满足
\[
(U-U)\cap \ker(\alpha_x)=\{0\}.
\]
于是对不同的 \(\eta,\eta'\in \ker(\alpha_x)\)，平移集 \(U+\eta\) 与 \(U+\eta'\) 两两不交。又因 \(\alpha_x\) 为群同态，
\[
\alpha_x(U+\eta)=\alpha_x(U)
\qquad
(\eta\in \ker(\alpha_x)),
\]
并且若 \(u_1,u_2\in U\) 满足 \(\alpha_x(u_1)=\alpha_x(u_2)\)，则
\[
u_1-u_2\in (U-U)\cap \ker(\alpha_x)=\{0\},
\]
故 \(\alpha_x|_U\) 单射。由于连续满射群同态是开映射，\(\alpha_x(U)\) 为开集，且
\[
\alpha_x|_U:U\longrightarrow \alpha_x(U)
\]
为同胚。再利用平移共轭，可知每个 \(\alpha_x|_{U+\eta}\) 都是到 \(\alpha_x(U)\) 的同胚，故在该开邻域上给出恰好 \(|\ker(\alpha_x)|\) 个局部分支。于是 \(\alpha_x\) 是度数为
\[
\card{\ker(\alpha_x)}=v
\]
的有限覆盖。证毕。
\end{proof}

\begin{corollary}[非零连续自同态覆盖度半群的严格算术分类]\label{cor:xi-localized-solenoid-nonzero-endomorphism-degree-semigroup}
\leanverified{paper\_xi\_localized\_solenoid\_nonzero\_endomorphism\_degree\_semigroup}
沿用定理 \ref{thm:xi-localized-solenoid-endomorphism-kernel-cokernel-sfree-numerator} 的记号。定义
\[
\deg:\operatorname{End}_{\mathrm{cts}}(\Sigma_S)\setminus\{0\}\to \NN_{\ge 1},
\qquad
\deg(\alpha):=\card{\ker(\alpha)}.
\]
则其像恰为
\[
\boxed{\;
\operatorname{Im}(\deg)
=
\left\{
n\in \NN_{\ge 1}:\ \gcd\!\left(n,\prod_{p\in S}p\right)=1
\right\}.\;}
\]
并且对任意非零连续自同态 \(\alpha_x,\alpha_y\) 都有
\[
\boxed{\;
\deg(\alpha_x\circ \alpha_y)=\deg(\alpha_x)\deg(\alpha_y).\;}
\]
因此，\(\Sigma_S\) 的全部非平凡连续覆盖半群与去掉 \(S\)-素数后的正整数乘法半群完全同构。
\end{corollary}
\begin{proof}
由定理 \ref{thm:xi-localized-endomorphism-ring-solenoid-indecomposable}，
\[
\operatorname{End}_{\mathrm{cts}}(\Sigma_S)\cong G_S,
\]
故每个非零连续自同态都唯一写成某个 \(\alpha_x\)。

任取
\[
x=\frac{a}{b}\in G_S\setminus\{0\},
\qquad
\gcd(a,b)=1,
\qquad
b\in \langle S\rangle.
\]
把
\[
|a|=uv
\]
写成 \(u\) 的全部素因子都属于 \(S\)、而 \(v\) 与 \(\prod_{p\in S}p\) 互素的分解。由推论 \ref{cor:xi-localized-solenoid-nonzero-endomorphism-finite-cover-degree}，
\[
\deg(\alpha_x)=v.
\]
因此每个覆盖度都与 \(S\) 互素。

反之，若 \(n\in\NN_{\ge1}\) 与 \(\prod_{p\in S}p\) 互素，则 \(n\in G_S\)，并且它的 \(S\)-外分子就是 \(n\) 本身。故
\[
\deg(\alpha_n)=n.
\]
从而上述集合中的每个正整数都可实现，这就给出 \(\operatorname{Im}(\deg)\) 的精确刻画。

最后证明乘法性。设
\[
x=\varepsilon_x\frac{u_xv_x}{b_x},
\qquad
y=\varepsilon_y\frac{u_yv_y}{b_y},
\]
其中 \(\varepsilon_x,\varepsilon_y\in\{\pm1\}\)，\(u_x,u_y,b_x,b_y\) 的全部素因子都属于 \(S\)，而 \(v_x,v_y\) 都与 \(\prod_{p\in S}p\) 互素。则
\[
xy=\varepsilon_x\varepsilon_y\frac{u_xu_y}{b_xb_y}\,v_xv_y.
\]
由于 \(\frac{u_xu_y}{b_xb_y}\in G_S^\times\)，故 \(xy\) 的 \(S\)-外分子恰为 \(v_xv_y\)。再由交换性，
\[
\alpha_x\circ \alpha_y=\alpha_{xy},
\]
于是
\[
\deg(\alpha_x\circ \alpha_y)
=
\deg(\alpha_{xy})
=
v_xv_y
=
\deg(\alpha_x)\deg(\alpha_y).
\]
证毕。
\end{proof}

\endinput


\begin{theorem}[一般连续自同态的周期点计数与 Artin--Mazur \texorpdfstring{$\zeta$}{zeta} 的算术闭式]\label{thm:xi-localized-solenoid-general-endomorphism-periodic-zeta}
\leanverified{paper\_xi\_localized\_solenoid\_general\_endomorphism\_periodic\_zeta}
沿用定理 \ref{thm:xi-localized-solenoid-endomorphism-kernel-cokernel-sfree-numerator} 的记号，并额外假设
\[
x\notin\{\pm 1\}.
\]
对任意非零整数 \(m\)，记
\[
|m|_{S^\perp}:=\frac{|m|}{\prod_{p\in S}p^{v_p(m)}}.
\]
则对每个 \(n\ge 1\)，不动点群满足
\[
\boxed{\;
\mathrm{Fix}(\alpha_x^n)
\cong
\ZZ/|a^n-b^n|_{S^\perp}\ZZ,
\qquad
\#\mathrm{Fix}(\alpha_x^n)=|a^n-b^n|_{S^\perp}.\;}
\]
因此其 Artin--Mazur \(\zeta\)-函数作为形式幂级数满足
\[
\boxed{\;
\zeta_{\alpha_x}(z)
:=
\exp\!\left(
\sum_{n\ge 1}\frac{\#\mathrm{Fix}(\alpha_x^n)}{n}z^n
\right)
=
\exp\!\left(
\sum_{n\ge 1}\frac{|a^n-b^n|_{S^\perp}}{n}z^n
\right).\;}
\]
\end{theorem}
\begin{proof}
局部化环 \(G_S=\ZZ[S^{-1}]\subset \QQ\) 的有限阶单位只有 \(\pm1\)。故 \(x\notin\{\pm1\}\) 蕴含 \(x^n\neq 1\) 对全部 \(n\ge 1\) 成立，从而下面出现的差分同态都对应非零有理数乘法。

由定义，
\[
\alpha_x^n=\widehat{m_x}^n=\widehat{m_{x^n}}.
\]
因此
\[
\alpha_x^n-\mathrm{id}_{\Sigma_S}
=
\widehat{m_{x^n}-\mathrm{id}_{G_S}}
=
\widehat{m_{x^n-1}}.
\]
于是
\[
\mathrm{Fix}(\alpha_x^n)
=
\ker(\alpha_x^n-\mathrm{id}_{\Sigma_S})
=
\ker(\widehat{m_{x^n-1}}).
\]

再写
\[
x^n-1=\frac{a^n-b^n}{b^n}.
\]
由于 \(b^n\in\langle S\rangle\)，它在 \(G_S\) 中为单位，故 \(x^n-1\) 的 \(S\)-外分子恰为 \(|a^n-b^n|_{S^\perp}\)。将定理 \ref{thm:xi-localized-solenoid-endomorphism-kernel-cokernel-sfree-numerator} 应用于非零元素 \(x^n-1\)，得到
\[
\ker(\widehat{m_{x^n-1}})
\cong
\ZZ/|a^n-b^n|_{S^\perp}\ZZ.
\]
从而所示不动点计数公式成立。

最后把
\[
\#\mathrm{Fix}(\alpha_x^n)=|a^n-b^n|_{S^\perp}
\]
代入 Artin--Mazur \(\zeta\)-函数的定义即可得到所示形式幂级数闭式。证毕。
\end{proof}

\noindent
上述结果把 prime-register 机制推进为一条完整的覆盖--动力学--群论链：\(\Sigma_S=\widehat{G_S}\) 的整数乘法覆盖度、\(n\)-挠子群、全部有限闭子群、整数乘法的周期点计数与 primitive 轨道数，以及任意连续自同态的核、余核与 Artin--Mazur \(\zeta\)-函数，都由 \(S\)-自由剥离操作严格控制。由此，\(\Sigma_S\) 不再只是“带有 \(p\)-进核的 solenoid”这一拓扑模型，而成为一个其有限层对称性、非可逆性与周期动力学都可被 prime-register 完全算术化的对象。

\endinput


\begin{theorem}[局部化数系有限指数子群与 profinite 完成的完全分类]\label{thm:xi-localized-finite-index-profinite-completion-classification}
\leanverified{paper\_xi\_localized\_finite\_index\_profinite\_completion\_classification}
沿用定理 \ref{thm:xi-localized-basic-extension-strictly-nonsplit} 的记号。对任意子群 \(H\le G_S\)，以下两条等价：
\begin{enumerate}
  \item \(H\) 在 \(G_S\) 中具有有限指数；
  \item 存在唯一正整数 \(m\) 使得对全部 \(p\in S\) 都有 \(p\nmid m\)，并且
  \[
  H=mG_S.
  \]
\end{enumerate}
在此情形下必有
\[
G_S/H\cong \ZZ/m\ZZ.
\]
进一步，若以
\[
\widehat{G_S}^{\,\mathrm{pf}}
\]
表示 \(G_S\) 的 profinite 完成，则
\[
\widehat{G_S}^{\,\mathrm{pf}}
\cong
\varprojlim_{\substack{m\ge 1\\ p\nmid m\ \forall p\in S}}\ZZ/m\ZZ
\cong
\prod_{\ell\notin S}\ZZ_\ell.
\]
\end{theorem}
\begin{proof}
若 \(H=mG_S\) 且对全部 \(p\in S\) 都有 \(p\nmid m\)，则定理 \ref{thm:xi-localized-finite-quotient-phase-layer-classification} 第 \(1\) 条给出
\[
G_S/H=G_S/mG_S\cong \ZZ/m\ZZ,
\]
故 \(H\) 具有有限指数。

反过来，设 \(H\le G_S\) 具有有限指数，并记
\[
Q:=G_S/H.
\]
由定理 \ref{thm:xi-localized-finite-quotient-phase-layer-classification} 第 \(2\) 条，存在唯一正整数 \(m\) 使得对全部 \(p\in S\) 都有 \(p\nmid m\)，且
\[
Q\cong \ZZ/m\ZZ.
\]
由于 \(Q\) 的指数为 \(m\)，任意同态 \(G_S\to Q\) 都杀掉 \(mG_S\)。因此自然商映射 \(G_S\twoheadrightarrow Q\) 经由
\[
G_S/mG_S
\]
因子化。再由定理 \ref{thm:xi-localized-finite-quotient-phase-layer-classification} 第 \(1\) 条，
\[
G_S/mG_S\cong \ZZ/m\ZZ
\]
也恰有 \(m\) 个元素。于是该因子化给出一个从 \(G_S/mG_S\) 到 \(Q\) 的满射，而两者基数同为 \(m\)，故该满射必为同构。其核只能为零，从而
\[
H=mG_S.
\]
唯一性随即由商群阶数的唯一性得到。

由前半部分，\(\{mG_S:\ p\nmid m\ \forall p\in S\}\) 与 \(G_S\) 的全部有限指数子群完全一致，故这些子群构成定义 profinite 完成的共尾系统。于是
\[
\widehat{G_S}^{\,\mathrm{pf}}
\cong
\varprojlim_{\substack{m\ge 1\\ p\nmid m\ \forall p\in S}}G_S/mG_S
\cong
\varprojlim_{\substack{m\ge 1\\ p\nmid m\ \forall p\in S}}\ZZ/m\ZZ.
\]
最后，由中国剩余定理与 profinite 整数的标准分解，
\[
\varprojlim_{\substack{m\ge 1\\ p\nmid m\ \forall p\in S}}\ZZ/m\ZZ
\cong
\prod_{\ell\notin S}\ZZ_\ell.
\]
证毕。
\end{proof}

上述分类表明，Pontryagin 对偶短正合列所显露的 \(S\)-adic 核与 profinite 完成所显露的 \(S^c\)-adic 有限商方向彼此互补；因而 \(G_S\) 的连续圆因子与离散有限商账本在两侧同时被锁定。

\endinput


\paragraph{局部化 solenoid 的周期指数、熵闭式、子群 \texorpdfstring{$\zeta$}{zeta} 与原始素因子强迫}
在 \(\Sigma_S=\widehat{G_S}\) 的连续自同态、周期点计数以及有限指数子群已经分别得到完全闭式之后，还可把增长指数、拓扑熵、子群 \(\zeta\) 函数与新素数强迫进一步压缩为下列几条纯算术后果。其共同特征是：\(S\)-内素数只改变有限误差、Euler 因子或有限例外，而指数级复杂度完全由分子--分母的最大尺度与 \(S\)-外账本决定。

\begin{theorem}[solenoid 标量自同态的周期点指数由分子--分母最大尺度精确控制]\label{thm:xi-localized-solenoid-periodic-growth-max-scale}
设 \(S\subset\mathbb P\) 为有限素数集，并取
\[
x=\frac{a}{b}\in G_S=\ZZ[S^{-1}]\setminus\{\pm1\},
\qquad
\gcd(a,b)=1,
\qquad
b\in\langle S\rangle.
\]
记
\[
\alpha_x:=\widehat{m_x}:\Sigma_S\to\Sigma_S,
\qquad
F_n(x):=\#\mathrm{Fix}(\alpha_x^n).
\]
则当 \(n\to\infty\) 时有
\[
\boxed{\;
\log F_n(x)
=
n\log \max\{|a|,|b|\}+O(\log n).\;}
\]
从而
\[
\boxed{\;
\lim_{n\to\infty}\frac{1}{n}\log F_n(x)
=
\log \max\{|a|,|b|\}.\;}
\]
\leanverified{paper\_xi\_localized\_solenoid\_periodic\_growth\_max\_scale}
\end{theorem}
\begin{proof}
由定理 \ref{thm:xi-localized-solenoid-general-endomorphism-periodic-zeta}，
\[
F_n(x)=|a^n-b^n|_{S^\perp}.
\]
因此
\[
\log F_n(x)
=
\log|a^n-b^n|-\sum_{p\in S}v_p(a^n-b^n)\log p.
\]

先估计 \(S\)-主部分。固定 \(p\in S\)。
若 \(p\mid ab\)，由于 \(\gcd(a,b)=1\)，\(p\) 只能整除 \(a,b\) 之一，从而
\[
a^n-b^n\equiv \pm b^n \not\equiv 0 \pmod p
\quad\text{或}\quad
a^n-b^n\equiv a^n \not\equiv 0 \pmod p,
\]
故
\[
v_p(a^n-b^n)=0.
\]

以下设 \(p\nmid ab\)。令 \(r_p\) 为 \(ab^{-1}\) 在 \((\ZZ/p\ZZ)^\times\) 中的乘法阶。
若 \(p\mid a^n-b^n\)，则 \(r_p\mid n\)；又因 \(r_p\mid p-1\)，故 \(p\nmid r_p\)。
把 \(n\) 写成
\[
n=r_p\,p^t s,
\qquad
t=v_p(n),
\qquad
p\nmid s.
\]
再令
\[
u_p:=a^{r_p}b^{-r_p}\in \ZZ_p^\times.
\]
由于 \(a^{r_p}\equiv b^{r_p}\pmod p\)，有 \(u_p\equiv 1\pmod p\)。
并且
\[
a^n-b^n
=
b^n\bigl(u_p^{p^t s}-1\bigr),
\]
而 \(b\in \ZZ_p^\times\)，故
\[
v_p(a^n-b^n)=v_p\bigl(u_p^{p^t s}-1\bigr).
\]
又因 \(u_p^{p^t}\equiv 1\pmod p\) 且 \(p\nmid s\)，几何级数
\[
1+u_p^{p^t}+\cdots+\bigl(u_p^{p^t}\bigr)^{s-1}
\equiv s \not\equiv 0 \pmod p
\]
是 \(p\)-进单位，于是
\[
v_p\bigl(u_p^{p^t s}-1\bigr)=v_p\bigl(u_p^{p^t}-1\bigr).
\]

若 \(p\) 为奇素数，写 \(u_p^{p^j}=1+\pi_j\) 且 \(v_p(\pi_j)\ge 1\)。则二项式展开给出
\[
(1+\pi_j)^p-1
=
p\pi_j+\sum_{k=2}^{p-1}\binom{p}{k}\pi_j^k+\pi_j^p,
\]
其中第一项赋值为 \(v_p(\pi_j)+1\)，其余各项赋值严格更大，故
\[
v_p\bigl(u_p^{p^{j+1}}-1\bigr)
=
v_p\bigl(u_p^{p^j}-1\bigr)+1.
\]
迭代可得
\[
v_p\bigl(u_p^{p^t}-1\bigr)=v_p(u_p-1)+t.
\]

若 \(p=2\)，则 \(u_p\) 为奇 \(2\)-进单位，从而 \(u_p^2\equiv 1\pmod 8\)。因此对每个 \(j\ge 1\)，有
\[
v_2\bigl(u_p^{2^{j+1}}-1\bigr)
=
v_2\bigl(u_p^{2^j}-1\bigr)+v_2\bigl(u_p^{2^j}+1\bigr)
=
v_2\bigl(u_p^{2^j}-1\bigr)+1,
\]
因为 \(u_p^{2^j}+1\equiv 2\pmod 8\)。于是存在只依赖于 \(a,b,p\) 的常数 \(C_p\) 使
\[
v_2\bigl(u_p^{2^t}-1\bigr)\le C_p+t.
\]

综上，对每个 \(p\in S\) 都存在常数 \(C_p'\) 使
\[
v_p(a^n-b^n)\le C_p'+v_p(n).
\]
由于 \(S\) 有限，
\[
\sum_{p\in S}v_p(a^n-b^n)\log p
\le
\sum_{p\in S}C_p'\log p+\sum_{p\in S}v_p(n)\log p
=
O(1)+O(\log n).
\]
故
\[
\log F_n(x)=\log|a^n-b^n|+O(\log n).
\]

最后设
\[
M:=\max\{|a|,|b|\},
\qquad
m:=\min\{|a|,|b|\}.
\]
由 \(x\notin\{\pm1\}\) 与 \(\gcd(a,b)=1\) 可知 \(m<M\)，记 \(\theta:=m/M\in(0,1)\)。
于是
\[
|a^n-b^n|
=
M^n\Bigl|1-\varepsilon\,\theta^n\Bigr|
\]
对某个 \(\varepsilon\in\{\pm1\}\) 成立，从而
\[
\log|a^n-b^n|
=
n\log M+O(1).
\]
代回即得
\[
\log F_n(x)=n\log M+O(\log n),
\]
进而极限公式成立。证毕。
\end{proof}

\begin{corollary}[周期点计数级数与 logarithmic 级数的收敛半径同由最大尺度决定]\label{cor:xi-localized-solenoid-periodic-series-radius-max-scale}
沿用定理 \ref{thm:xi-localized-solenoid-periodic-growth-max-scale} 的记号，并记
\[
M:=\max\{|a|,|b|\}.
\]
定义两类周期点级数
\[
\mathcal P_x(z):=\sum_{n\ge 1}F_n(x)\,z^n,
\qquad
\mathcal L_x(z):=\sum_{n\ge 1}\frac{F_n(x)}{n}\,z^n.
\]
则二者的收敛半径都恰为
\[
\boxed{\;
R_x=\frac{1}{M}.\;}
\]
\end{corollary}
\begin{proof}
由定理 \ref{thm:xi-localized-solenoid-periodic-growth-max-scale}，
\[
\log F_n(x)=n\log M+O(\log n).
\]
因此
\[
\limsup_{n\to\infty}F_n(x)^{1/n}=M.
\]
又因为
\[
\left(\frac{F_n(x)}{n}\right)^{1/n}
=
F_n(x)^{1/n}\,n^{-1/n},
\]
而 \(n^{-1/n}\to 1\)，故
\[
\limsup_{n\to\infty}\left(\frac{F_n(x)}{n}\right)^{1/n}=M.
\]
对 \(\mathcal P_x\) 与 \(\mathcal L_x\) 分别应用 Cauchy--Hadamard 公式，即得二者的收敛半径都等于 \(M^{-1}\)。证毕。
\end{proof}

\begin{corollary}[单位情形的拓扑熵仅由有理尺度 \texorpdfstring{$\max\{|a|,|b|\}$}{max\{|a|,|b|\}} 决定]\label{cor:xi-localized-solenoid-unit-entropy-max-scale}
在定理 \ref{thm:xi-localized-solenoid-periodic-growth-max-scale} 的设定下，若 \(x\in G_S^\times\)，则 \(\alpha_x\) 为连续自同构，并且
\[
\boxed{\;
h_{\mathrm{top}}(\alpha_x)=\log\max\{|a|,|b|\}.\;}
\]
\end{corollary}
\begin{proof}
由定理 \ref{thm:xi-localized-solenoid-endomorphism-kernel-cokernel-sfree-numerator} 第 \(3\) 条，\(x\in G_S^\times\) 当且仅当 \(\alpha_x\) 为同构。对 solenoid 自同构，Lind--Ward 的 \(p\)-进熵分解给出
\[
h_{\mathrm{top}}(\alpha_x)
=
\sum_{v\le \infty}\max\{\log|x|_v,0\},
\]
其中 \(v\) 遍历 \(\QQ\) 的 Archimedean 与非 Archimedean 赋值 \cite{LindWard1988SolenoidsPadicEntropy}。

若 \(|a|\ge |b|\)，则
\[
\max\{\log|x|_\infty,0\}=\log\frac{|a|}{|b|}.
\]
又因 \(\gcd(a,b)=1\)，对任意素数 \(p\) 有
\[
\max\{\log|x|_p,0\}=
\begin{cases}
v_p(b)\log p,& p\mid b,\\
0,& p\nmid b.
\end{cases}
\]
故
\[
h_{\mathrm{top}}(\alpha_x)
=
\log\frac{|a|}{|b|}
+\sum_{p\mid b}v_p(b)\log p
=
\log|a|.
\]

若 \(|a|<|b|\)，则 Archimedean 贡献为 \(0\)，并且只有 \(p\mid b\) 时才有正的 \(p\)-进贡献，因此
\[
h_{\mathrm{top}}(\alpha_x)
=
\sum_{p\mid b}v_p(b)\log p
=
\log|b|.
\]
两种情形合并即得
\[
h_{\mathrm{top}}(\alpha_x)=\log\max\{|a|,|b|\}.
\]
证毕。
\end{proof}

\begin{theorem}[局部化数系有限指数子群的 \texorpdfstring{$\zeta$}{zeta} 函数闭式]\label{thm:xi-localized-subgroup-zeta-euler-product}
\leanverified{paper\_xi\_localized\_subgroup\_zeta\_euler\_product}
设
\[
P_S:=\prod_{p\in S}p.
\]
则 \(G_S=\ZZ[S^{-1}]\) 的有限指数子群恰为
\[
H=mG_S,
\qquad
m\ge 1,
\qquad
\gcd(m,P_S)=1,
\]
并且该表示唯一。在此基础上，其子群 \(\zeta\) 函数
\[
\zeta_{G_S}^{\le}(s)
:=
\sum_{H\le_f G_S}[G_S:H]^{-s}
\]
在 \(\Re(s)>1\) 上满足精确 Euler 乘积
\[
\boxed{\;
\zeta_{G_S}^{\le}(s)
=
\sum_{\substack{m\ge 1\\ (m,P_S)=1}}m^{-s}
=
\zeta(s)\prod_{p\in S}(1-p^{-s}).\;}
\]
\end{theorem}
\begin{proof}
有限指数子群的完全分类已由定理 \ref{thm:xi-localized-finite-index-profinite-completion-classification} 给出：\(H\le G_S\) 具有有限指数，当且仅当
\[
H=mG_S
\]
且 \(\gcd(m,P_S)=1\)，并且该 \(m\) 唯一。又由同一定理，
\[
G_S/mG_S\cong \ZZ/m\ZZ,
\]
故
\[
[G_S:mG_S]=m.
\]
因此
\[
\zeta_{G_S}^{\le}(s)
=
\sum_{\substack{m\ge 1\\ (m,P_S)=1}}m^{-s}.
\]
在 \(\Re(s)>1\) 上按素数分解 Euler 乘积，
\[
\sum_{\substack{m\ge 1\\ (m,P_S)=1}}m^{-s}
=
\prod_{p\notin S}(1-p^{-s})^{-1}
=
\zeta(s)\prod_{p\in S}(1-p^{-s}).
\]
证毕。
\end{proof}

\begin{corollary}[有限指数子群计数的主项、显式误差与极点留数]\label{cor:xi-localized-subgroup-counting-main-term}
\leanverified{paper\_xi\_localized\_subgroup\_counting\_main\_term}
沿用定理 \ref{thm:xi-localized-subgroup-zeta-euler-product} 的记号。定义
\[
N_S(X)
:=
\#\{H\le_f G_S:\ [G_S:H]\le X\}.
\]
则对任意 \(X\ge 1\)，有精确公式
\[
N_S(X)
=
\#\{1\le m\le X:\ \gcd(m,P_S)=1\}
=
\sum_{d\mid P_S}\mu(d)\left\lfloor\frac{X}{d}\right\rfloor.
\]
从而
\[
\boxed{\;
N_S(X)
=
\frac{\varphi(P_S)}{P_S}\,X+O(2^{|S|}),
\qquad
X\to\infty.\;}
\]
并且
\[
\boxed{\;
\Res_{s=1}\zeta_{G_S}^{\le}(s)
=
\frac{\varphi(P_S)}{P_S}.\;}
\]
\end{corollary}
\begin{proof}
由定理 \ref{thm:xi-localized-subgroup-zeta-euler-product}，
\[
N_S(X)=\#\{1\le m\le X:\ \gcd(m,P_S)=1\}.
\]
对互素条件施行 Möbius 反演，
\[
\mathbf 1_{\gcd(m,P_S)=1}
=
\sum_{d\mid \gcd(m,P_S)}\mu(d)
=
\sum_{\substack{d\mid P_S\\ d\mid m}}\mu(d).
\]
对 \(m\le X\) 求和并交换顺序，得到
\[
N_S(X)
=
\sum_{d\mid P_S}\mu(d)\sum_{\substack{m\le X\\ d\mid m}}1
=
\sum_{d\mid P_S}\mu(d)\left\lfloor\frac{X}{d}\right\rfloor.
\]
再用
\[
\left\lfloor\frac{X}{d}\right\rfloor=\frac{X}{d}+O(1)
\]
得到
\[
N_S(X)
=
X\sum_{d\mid P_S}\frac{\mu(d)}{d}
+O\!\left(\sum_{d\mid P_S}1\right).
\]
由于 \(P_S\) 为平方自由数，
\[
\sum_{d\mid P_S}1=2^{|S|},
\qquad
\sum_{d\mid P_S}\frac{\mu(d)}{d}
=
\prod_{p\in S}\left(1-\frac1p\right)
=
\frac{\varphi(P_S)}{P_S},
\]
故主项展开成立。

留数公式由定理 \ref{thm:xi-localized-subgroup-zeta-euler-product} 的 Euler 乘积与 \(\zeta(s)\) 在 \(s=1\) 的单极点立即得到。证毕。
\end{proof}

\begin{theorem}[周期点账本的 \texorpdfstring{$S$}{S}-外原始素因子强迫]\label{thm:xi-localized-solenoid-periodic-zsigmondy-forcing}
在定理 \ref{thm:xi-localized-solenoid-periodic-growth-max-scale} 的设定下，存在有限集合
\[
E=E(a,b,S)\subset \NN
\]
使得对每个 \(n\notin E\)，都存在素数
\[
q_n\notin S
\]
满足
\[
q_n\mid F_n(x),
\qquad
q_n\nmid F_k(x)\quad(1\le k<n).
\]
换言之，周期点计数序列 \(\{F_n(x)\}_{n\ge 1}\) 在 \(S\)-外素数方向上具有原始素因子强迫性。
\end{theorem}
\begin{proof}
记
\[
U_n:=\frac{a^n-b^n}{a-b}\qquad(n\ge 1).
\]
由于 \(a/b\notin\{\pm1\}\)，序列 \((U_n)\) 是非退化 Lucas 型序列。Bilu--Hanrot--Voutier 关于 Lucas 与 Lehmer 序列 primitive divisor 的定理表明：存在有限集合 \(E_0=E_0(a,b)\subset\NN\)，使得对每个 \(n\notin E_0\)，\(U_n\) 都存在 primitive prime divisor \(q_n\) \cite{BiluHanrotVoutier2001PrimitiveDivisorsLucas}。

对这样的 \(n\)，由 primitive 性可知
\[
q_n\mid U_n,
\qquad
q_n\nmid (a-b)U_1U_2\cdots U_{n-1}.
\]
因此
\[
q_n\mid a^n-b^n,
\qquad
q_n\nmid a^k-b^k\quad(1\le k<n),
\]
即 \(q_n\) 也是 \(a^n-b^n\) 的 primitive prime divisor。

现证明当 \(n\) 充分大时 \(q_n\notin S\)。若反设 \(q_n\in S\)，则上式自动给出 \(q_n\nmid ab\)。于是 \(ab^{-1}\) 在 \((\ZZ/q_n\ZZ)^\times\) 中有定义，并且 primitive 性推出其乘法阶恰为 \(n\)。故
\[
n\mid (q_n-1).
\]
记
\[
M_S:=
\begin{cases}
0,& S=\varnothing,\\
\max_{q\in S}(q-1),& S\neq\varnothing.
\end{cases}
\]
则当 \(n>M_S\) 时，上述整除关系不可能成立。于是对全部
\[
n\notin E_1:=\{1,\dots,M_S\},
\]
primitive prime divisor 必不属于 \(S\)。

最后令
\[
E:=E_0\cup E_1.
\]
对任意 \(n\notin E\)，已有 \(q_n\notin S\) 且 \(q_n\mid a^n-b^n\)。再由定理 \ref{thm:xi-localized-solenoid-general-endomorphism-periodic-zeta}，
\[
F_n(x)=|a^n-b^n|_{S^\perp},
\]
所以
\[
q_n\mid F_n(x).
\]
另一方面，若对某个 \(k<n\) 有 \(q_n\mid F_k(x)\)，由于 \(q_n\notin S\)，便必有 \(q_n\mid a^k-b^k\)，这与 \(q_n\) 的 primitive 性矛盾。故
\[
q_n\nmid F_k(x)\qquad(1\le k<n).
\]
证毕。
\end{proof}

\noindent
这组结果把 \(G_S\) 的有限指数层、\(\Sigma_S\) 的周期层以及单位自同态的熵层压缩到同一条可计算账本上：指数主项由 \(\max\{|a|,|b|\}\) 给出，有限指数复杂度由删去 \(S\)-Euler 因子的整数计数控制，而足够高的周期层则必然在 \(S\)-外方向引入新的原始素因子。

\endinput


\paragraph{局部化有限指数子群格、补素完成与全素数账本的闭合}
在有限指数子群的完全分类、相应子群 \(\zeta\) 函数的 Euler 乘积以及局部化对偶短正合列的严格不分裂都已建立之后，同一 rank-one 局部化数系的有限指数算术还可进一步压缩为下列三条 lattice/profinite 后果。

\begin{theorem}[局部化数系有限指数子群格的指标保持分类]\label{thm:xi-time-part56e-localized-finite-index-lattice-classification}
\leanverified{paper\_xi\_time\_part56e\_localized\_finite\_index\_lattice\_classification}
设 \(S\subset \mathbb P\) 为有限素数集，并记
\[
P_S:=\prod_{p\in S}p.
\]
则映射
\[
n\longmapsto nG_S
\qquad
\bigl(n\ge 1,\ \gcd(n,P_S)=1\bigr)
\]
在满足 \(\gcd(n,P_S)=1\) 的正整数与 \(G_S\) 的全部有限指数子群之间给出双射，并且保持指标：
\[
[G_S:nG_S]=n.
\]
等价地，\(\operatorname{Sub}_f(G_S)\) 与 \(S\)-自由正整数按整除关系形成的偏序格反向同构。
\end{theorem}
\begin{proof}
定理 \ref{thm:xi-localized-finite-index-profinite-completion-classification} 已证明：\(H\le G_S\) 具有有限指数，当且仅当存在唯一正整数 \(n\) 使
\[
H=nG_S,
\qquad
\gcd(n,P_S)=1,
\]
并且
\[
G_S/H\cong \ZZ/n\ZZ.
\]
因此
\[
[G_S:H]=n=[G_S:nG_S].
\]
双射与指标保持性遂同时成立。最后，由
\[
nG_S\subseteq mG_S
\iff
m\mid n
\]
可知该对应把子群包含关系转化为整数整除关系的反向，故得到所述反向同构。证毕。
\end{proof}

\begin{corollary}[局部化有限指数子群格的 \texorpdfstring{$\gcd/\mathrm{lcm}$}{gcd/lcm} 算术与区间 Möbius 闭式]\label{cor:xi-time-part56e-localized-finite-index-lattice-gcd-lcm-mobius}
沿用定理 \ref{thm:xi-time-part56e-localized-finite-index-lattice-classification} 的记号。对任意满足
\[
\gcd(m,P_S)=\gcd(n,P_S)=1
\]
的正整数，记
\[
H_m:=mG_S,
\qquad
H_n:=nG_S.
\]
则在有限指数子群格 \(\operatorname{Sub}_f(G_S)\) 中有
\[
H_m\wedge H_n
=
H_m\cap H_n
=
\operatorname{lcm}(m,n)\,G_S,
\]
\[
H_m\vee H_n
=
H_m+H_n
=
\gcd(m,n)\,G_S.
\]
进一步，若 \(m\mid n\) 且 \(\gcd(n,P_S)=1\)，则区间 Möbius 函数满足
\[
\mu_{\operatorname{Sub}_f(G_S)}(H_n,H_m)=\mu(n/m),
\]
其中右端为经典数论 Möbius 函数。
\end{corollary}
\begin{proof}
由定理 \ref{thm:xi-time-part56e-localized-finite-index-lattice-classification}，
\(\operatorname{Sub}_f(G_S)\) 与 \(S\)-自由正整数的整除格反向同构。故其 meet 与 join 分别对应于整数侧的 \(\operatorname{lcm}\) 与 \(\gcd\)，从而得到
\[
H_m\cap H_n=\operatorname{lcm}(m,n)\,G_S,
\qquad
H_m+H_n=\gcd(m,n)\,G_S.
\]
若愿意直接验证第二式，只需记 \(d:=\gcd(m,n)\)。由 Bézout 恒等式可取 \(u,v\in\ZZ\) 使
\[
um+vn=d,
\]
故 \(dG_S\subseteq mG_S+nG_S\)；反向包含由 \(d\mid m,n\) 立即成立。

再设 \(m\mid n\)。则
\[
H_n\subseteq H_m,
\]
且区间 \([H_n,H_m]\) 与正整数除子格
\[
\{\,d:\ d\mid n/m\,\}
\]
反向同构。有限除子格的 Möbius 函数正是经典 Möbius 函数 \(\mu\)，故
\[
\mu_{\operatorname{Sub}_f(G_S)}(H_n,H_m)=\mu(n/m).
\]
证毕。
\end{proof}

\begin{theorem}[局部化数系的子群 \texorpdfstring{$\zeta$}{zeta}、profinite 完成与补素核的全素数分裂]\label{thm:xi-time-part56e-localized-subgroup-zeta-profinite-complement-splitting}
设 \(S\subset\mathbb P\) 为有限素数集，并记
\[
\pi_S:\widehat{G_S}\twoheadrightarrow \TT,
\qquad
K_S:=\ker(\pi_S).
\]
则同时有
\[
\zeta_{G_S}^{\le}(s)
=
\zeta(s)\prod_{p\in S}(1-p^{-s}),
\]
\[
\widehat{G_S}^{\,\mathrm{pf}}
\cong
\prod_{\ell\notin S}\ZZ_\ell,
\qquad
K_S
\cong
\prod_{p\in S}\ZZ_p,
\]
从而得到规范拓扑群同构
\[
\widehat{G_S}^{\,\mathrm{pf}}\times K_S
\cong
\prod_{p\in\mathbb P}\ZZ_p
=
\widehat{\ZZ}^{\,\mathrm{pf}}.
\]
等价地，\(G_S\) 的有限商方向与其 Pontryagin 对偶上的全不连通核方向恰按 \(S^c\)-素数与 \(S\)-素数互补分裂，并在合并后恢复完整的全素数 profinite 账本。
\end{theorem}
\begin{proof}
第一式正是定理 \ref{thm:xi-localized-subgroup-zeta-euler-product}。
第二式中的 profinite 完成分类由定理
\ref{thm:xi-localized-finite-index-profinite-completion-classification}
给出：
\[
\widehat{G_S}^{\,\mathrm{pf}}
\cong
\prod_{\ell\notin S}\ZZ_\ell.
\]
另一方面，定理 \ref{thm:xi-localized-basic-extension-strictly-nonsplit}
给出 Pontryagin 对偶短正合列
\[
0\longrightarrow \prod_{p\in S}\ZZ_p
\longrightarrow \widehat{G_S}
\xrightarrow{\ \pi_S\ }
\TT
\longrightarrow 0,
\]
故
\[
K_S\cong \prod_{p\in S}\ZZ_p.
\]
将两式相乘即可得到
\[
\widehat{G_S}^{\,\mathrm{pf}}\times K_S
\cong
\left(\prod_{\ell\notin S}\ZZ_\ell\right)\times
\left(\prod_{p\in S}\ZZ_p\right)
\cong
\prod_{p\in\mathbb P}\ZZ_p.
\]
证毕。
\end{proof}

\endinput

\paragraph{有限素数支撑与紧核族之间的 Galois 型偏序恢复}
在有限指数子群格、profinite 完成以及补素核的全素数分裂都已闭合之后，有限素数支撑 \(S\subset\mathbb P\) 与紧全不连通核
\[
K_S\cong \prod_{p\in S}\ZZ_p
\]
之间还存在一条更强的偏序对应：不仅每个有限素数集可由核本身恢复，而且包含关系也可由核间的可商性精确读取。

\begin{theorem}[有限素数支撑的 Galois 型对应]\label{thm:xi-time-part56ea-prime-support-kernel-galois-correspondence}
\leanverified{paper\_xi\_time\_part56ea\_prime\_support\_kernel\_galois\_correspondence}
对任意有限素数集 \(S,T\subset\mathbb P\)，记
\[
G_S:=\ZZ[S^{-1}],
\qquad
K_S:=\ker(\widehat{G_S}\to \TT).
\]
则下列命题严格等价：
\begin{enumerate}
  \item \(S\subseteq T\)；
  \item 存在连续开满射
  \[
  \pi_{T,S}:K_T\to K_S;
  \]
  \item 存在短正合列
  \[
  0\longrightarrow \prod_{p\in T\setminus S}\ZZ_p
  \longrightarrow K_T
  \xrightarrow{\ \pi_{T,S}\ }
  K_S
  \longrightarrow 0.
  \]
\end{enumerate}
特别地，有限素数集按包含形成的偏序与紧核族 \(\{K_S\}\) 按连续开商映射形成的偏序完全一致。
\end{theorem}
\begin{proof}
由定理~\ref{thm:xi-time-part56e-localized-subgroup-zeta-profinite-complement-splitting}，
\[
K_S\cong \prod_{p\in S}\ZZ_p,
\qquad
K_T\cong \prod_{p\in T}\ZZ_p.
\]
若 \(S\subseteq T\)，则坐标投影
\[
\prod_{p\in T}\ZZ_p\longrightarrow \prod_{p\in S}\ZZ_p
\]
给出连续开满射，其核正是
\[
\prod_{p\in T\setminus S}\ZZ_p.
\]
故 \((1)\Rightarrow(2)\Rightarrow(3)\)。

反过来，设存在连续开满射
\[
\pi_{T,S}:K_T\to K_S.
\]
对任意素数 \(p\)，取模 \(p\) 商得到诱导满射
\[
K_T/pK_T\longrightarrow K_S/pK_S.
\]
而由
\[
K_U\cong \prod_{\ell\in U}\ZZ_\ell
\]
可直接计算
\[
K_U/pK_U\cong
\begin{cases}
\FF_p, & p\in U,\\
0, & p\notin U.
\end{cases}
\]
因此，若 \(p\in S\)，则 \(K_S/pK_S\neq 0\)，从而满射迫使 \(K_T/pK_T\neq 0\)，故 \(p\in T\)。
这就证明了 \(S\subseteq T\)，即 \((2)\Rightarrow(1)\)。
至于 \((3)\Rightarrow(2)\) 则为正合列定义的直接内容。证毕。
\end{proof}

\endinput


\paragraph{局部化数系的有限商、torsion 谱与 Dirichlet 账本}
在有限商分类、solenoid 的整数乘法核度公式与有限指数子群 Euler 乘积已经确立之后，局部化数系的 quotient 层与 torsion 层还可进一步压缩为同一组 Dirichlet 账本。

\begin{theorem}[局部化数系有限商的指数账本与循环实现判据]\label{thm:xi-localized-quotient-index-ledger}
沿用前述记号。对任意正整数 \(n\)，记
\[
n_{S^\perp}:=\frac{n}{\prod_{p\in S}p^{v_p(n)}}.
\]
则有
\[
G_S/nG_S\cong \ZZ/n_{S^\perp}\ZZ,
\qquad
[G_S:nG_S]=n_{S^\perp}.
\]
特别地，有限循环群 \(\ZZ/m\ZZ\) 出现在 \(G_S\) 的有限商中，当且仅当 \(m\) 的全部素因子都不属于 \(S\)。
\end{theorem}
\begin{proof}
定理 \ref{thm:xi-localized-finite-quotient-phase-layer-classification} 已给出
\[
G_S/nG_S\cong \ZZ/n_{S^\perp}\ZZ.
\]
两边取基数即得指数公式。最后一句正是同一定理对有限商分类的重述。证毕。
\end{proof}

\begin{theorem}[Pontryagin 对偶的 \texorpdfstring{$n$}{n}-挠谱与 exact-order 计数]\label{thm:xi-localized-torsion-exact-order-ledger}
沿用上述记号，并记 \(\Sigma_S:=\widehat{G_S}\)。对任意正整数 \(n\)，有
\[
\Sigma_S[n]\cong \ZZ/n_{S^\perp}\ZZ\cong \mu_{n_{S^\perp}},
\qquad
|\Sigma_S[n]|=n_{S^\perp}.
\]
其中 \(\mu_m:=\{z\in\TT:\ z^m=1\}\)。更进一步，恰为 \(n\) 阶的元素个数满足
\[
\#\{x\in\Sigma_S:\ \operatorname{ord}(x)=n\}
=
\begin{cases}
\varphi(n),& \text{\(n\) 的全部素因子都不属于 }S,\\[1mm]
0,& \text{否则}.
\end{cases}
\]
\end{theorem}
\begin{proof}
由定理 \ref{thm:xi-localized-solenoid-finite-subgroups-are-kernels}，
\[
\Sigma_S[n]=\ker[n]\cong \ZZ/n_{S^\perp}\ZZ.
\]
有限循环群与 \(\mu_{n_{S^\perp}}\) 同构，故第一式成立，取基数即得第二式。
又因 \(\Sigma_S[n]\) 为阶 \(n_{S^\perp}\) 的循环群，其中恰为 \(n\) 阶的元素个数在 \(n\mid n_{S^\perp}\) 时为 \(\varphi(n)\)，否则为 \(0\)。而 \(n\mid n_{S^\perp}\) 当且仅当 \(n\) 的全部素因子都不属于 \(S\)。证毕。
\end{proof}

\begin{theorem}[quotient-zeta 与 torsion-zeta 的 Euler 乘积闭式]\label{thm:xi-localized-quotient-torsion-zeta-euler-product}
沿用上述记号，定义 Dirichlet 级数
\[
Z_S^{\mathrm{quot}}(s):=\sum_{n\ge 1}\frac{[G_S:nG_S]}{n^s},
\qquad
Z_S^{\mathrm{tor}}(s):=\sum_{n\ge 1}\frac{|\Sigma_S[n]|}{n^s}.
\]
则二者在 \(\Re(s)>2\) 上绝对收敛，并满足
\[
Z_S^{\mathrm{quot}}(s)=Z_S^{\mathrm{tor}}(s)
=
\prod_{p\in S}\frac{1}{1-p^{-s}}
\prod_{p\notin S}\frac{1}{1-p^{1-s}}
=
\zeta(s-1)\prod_{p\in S}\frac{1-p^{1-s}}{1-p^{-s}}.
\]
若再定义 exact-order Dirichlet 级数
\[
E_S(s):=\sum_{n\ge 1}\frac{\#\{x\in\Sigma_S:\ \operatorname{ord}(x)=n\}}{n^s},
\]
则在 \(\Re(s)>2\) 上有
\[
E_S(s)
=
\frac{\zeta(s-1)}{\zeta(s)}
\prod_{p\in S}\frac{1-p^{1-s}}{1-p^{-s}}.
\]
\end{theorem}
\begin{proof}
由定理 \ref{thm:xi-localized-quotient-index-ledger} 与定理
\ref{thm:xi-localized-torsion-exact-order-ledger}，
\[
[G_S:nG_S]=n_{S^\perp}=|\Sigma_S[n]|.
\]
故两条 Dirichlet 级数逐项相等。
又因
\[
0\le n_{S^\perp}\le n,
\]
故当 \(\sigma:=\Re(s)>2\) 时，
\[
\sum_{n\ge 1}\left|\frac{[G_S:nG_S]}{n^s}\right|
=
\sum_{n\ge 1}\frac{n_{S^\perp}}{n^\sigma}
\le
\sum_{n\ge 1}n^{1-\sigma}<\infty,
\]
从而 \(Z_S^{\mathrm{quot}}\) 与 \(Z_S^{\mathrm{tor}}\) 都绝对收敛。

函数 \(n\mapsto n_{S^\perp}\) 为乘法函数，且
\[
p^e\longmapsto
\begin{cases}
1,& p\in S,\\[1mm]
p^e,& p\notin S.
\end{cases}
\]
因此 Euler 因子分别为
\[
\sum_{e\ge 0}\frac{1}{p^{es}}=\frac{1}{1-p^{-s}}
\qquad (p\in S),
\]
\[
\sum_{e\ge 0}\frac{p^e}{p^{es}}=\frac{1}{1-p^{1-s}}
\qquad (p\notin S),
\]
从而得到所示 Euler 乘积；再把 \(\prod_{p\notin S}(1-p^{1-s})^{-1}\) 写成 \(\zeta(s-1)\) 去掉 \(S\)-内 Euler 因子后的形式，即得第二个闭式。

对 \(E_S(s)\)，由定理 \ref{thm:xi-localized-torsion-exact-order-ledger}，其系数函数为
\[
e_S(n):=\#\{x\in\Sigma_S:\ \operatorname{ord}(x)=n\}
=
\begin{cases}
\varphi(n),& \text{\(n\) 的全部素因子都不属于 }S,\\[1mm]
0,& \text{否则}.
\end{cases}
\]
又因 \(0\le e_S(n)\le \varphi(n)\le n\)，故 \(E_S(s)\) 在 \(\Re(s)>2\) 上亦绝对收敛。
故 \(e_S\) 为乘法函数，且对应 Euler 因子为
\[
1
\qquad (p\in S),
\]
以及
\[
1+\sum_{e\ge 1}\frac{\varphi(p^e)}{p^{es}}
=
1+\sum_{e\ge 1}\frac{p^e-p^{e-1}}{p^{es}}
=
\frac{1-p^{-s}}{1-p^{1-s}}
\qquad (p\notin S).
\]
将全部素数相乘，并把 \(p\notin S\) 的乘积重写为
\[
\prod_{p}\frac{1-p^{-s}}{1-p^{1-s}}
\prod_{p\in S}\frac{1-p^{1-s}}{1-p^{-s}}
=
\frac{\zeta(s-1)}{\zeta(s)}
\prod_{p\in S}\frac{1-p^{1-s}}{1-p^{-s}},
\]
便得所示公式。证毕。
\end{proof}

\begin{corollary}[有限商指数函数在素数点上的完全恢复]\label{cor:xi-localized-quotient-index-prime-recovery}
定义
\[
I_S(n):=[G_S:nG_S].
\]
则对每个 \(n\ge 1\) 都有
\[
I_S(n)=n_{S^\perp}.
\]
特别地，对任意素数 \(p\) 有
\[
p\in S \iff I_S(p)=1,
\qquad
p\notin S \iff I_S(p)=p.
\]
因此，仅由素数点取值模式 \(p\mapsto I_S(p)\) 即可完全且唯一恢复 \(S\)。
\end{corollary}
\begin{proof}
第一式是定理 \ref{thm:xi-localized-quotient-index-ledger} 的指数公式。对 \(n=p\) 代入即得第二式，故结论成立。证毕。
\end{proof}


\begin{theorem}[连通 Lie 商的圆周刚性]\label{thm:xi-localized-connected-lie-quotient-circle-rigidity}
沿用定理 \ref{thm:xi-localized-basic-extension-strictly-nonsplit} 的记号。设 \(K\) 是 \(\widehat{G_S}\) 的连通紧 Lie 商。则必有
\[
K\cong \TT^r
\qquad\text{且}\qquad
r\le 1.
\]
特别地，任何非平凡连通紧 Lie 商都同构于 \(\TT\)。
\end{theorem}
\begin{proof}
连通紧 Lie 交换群必同构于某个环面 \(\TT^r\)。设
\[
\pi:\widehat{G_S}\twoheadrightarrow K
\]
为满射，则对偶化得到单射
\[
\widehat K\hookrightarrow G_S.
\]
而
\[
\widehat K\cong \ZZ^r.
\]
由于 \(G_S=\ZZ[S^{-1}]\subset \QQ\)，有
\[
G_S\otimes_{\ZZ}\QQ\cong \QQ,
\]
故 \(G_S\) 的有理秩为 \(1\)。于是其任意自由阿贝尔子群的秩至多为 \(1\)，从而 \(r\le 1\)。若 \(K\) 非平凡，则 \(r\neq 0\)，故必有 \(r=1\)，即
\[
K\cong \TT.
\]
证毕。
\end{proof}

\begin{theorem}[非平凡局部化数系与有限生成离散交换群的相位采样不可混同]\label{thm:xi-localized-phase-sampling-not-finitely-generated}
\leanverified{paper\_xi\_localized\_phase\_sampling\_not\_finitely\_generated}
设 \(S\subset\mathbb P\) 为非空有限集。则不存在有限生成离散交换群 \(H\) 使得
\[
\card{\operatorname{Hom}(H,\ZZ/N\ZZ)}
=
\Sigma^{\mathrm{ph}}_{G_S}(N)
\qquad(\forall\,N\ge 1).
\]
\end{theorem}
\begin{proof}
反设存在
\(
H\cong \ZZ^{r}\oplus F
\)
满足上述恒等式。对任意素数 \(p\) 与 \(a\ge 1\)，由注 \ref{rem:xi-cdim-hom-quotient-consistency}，
\[
\card{\operatorname{Hom}(H,\ZZ/p^a\ZZ)}=\card{H/p^aH}.
\]
再由命题 \ref{prop:xi-cdim-local-inversion-delta}，定义
\[
c_p(a):=
v_p\!\bigl(\card{H/p^aH}\bigr)-ar
=
\sum_{j=1}^{t_p}\min(\lambda_{p,j},a),
\]
其中 \(F_p\cong\bigoplus_{j=1}^{t_p}\ZZ/p^{\lambda_{p,j}}\ZZ\) 为 \(F\) 的 \(p\)-primary 分量。因而 \(c_p(a)\) 对 \(a\) 有界，从而
\[
\lim_{a\to\infty}\frac1a
v_p\!\bigl(\card{\operatorname{Hom}(H,\ZZ/p^a\ZZ)}\bigr)
=
r
\qquad\text{对每个素数 }p\text{ 都成立}.
\]
另一方面，由定理 \ref{thm:xi-localized-phase-sampling-closed-form}，
\[
\Sigma^{\mathrm{ph}}_{G_S}(p^a)=
\begin{cases}
1,& p\in S,\\[1mm]
p^a,& p\notin S,
\end{cases}
\]
故
\[
\frac1a v_p\!\bigl(\Sigma^{\mathrm{ph}}_{G_S}(p^a)\bigr)
=
\begin{cases}
0,& p\in S,\\[1mm]
1,& p\notin S.
\end{cases}
\]
由于 \(S\) 非空且有限，可同时取到 \(p\in S\) 与 \(q\notin S\)，从而上述极限在素数间取两个不同值 \(0\) 与 \(1\)。这与有限生成离散交换群在全部素数上共享同一秩斜率 \(r\) 矛盾。故不存在这样的 \(H\)。证毕。
\end{proof}

\paragraph{有限值连续 prime-audit 的圆柱因子化与指数级歧义}
局部化数系的 Pontryagin 对偶已把 prime-register 核显式识别为有限个 \(p\)-进坐标的乘积；在此基础上，任何有限值连续审计都必定降到有限精度的圆柱商。

\begin{theorem}[有限值连续 prime-audit 必因子化到有限圆柱商]\label{thm:xi-localized-finite-prime-audit-cylinder-factorization}
\leanverified{paper\_xi\_localized\_finite\_prime\_audit\_cylinder\_factorization}
设 \(S\subset\mathbb P\) 为有限素数集，并记
\[
\mathfrak K_S:=\prod_{p\in S}\ZZ_p.
\]
由定理 \ref{thm:killo-localization-circle-dimension-prime-ledger-splitting}，\(\mathfrak K_S\) 正是 \(\widehat{G_S}\) 的 prime-register 核。令 \(A\) 为有限离散集合，\(F:\mathfrak K_S\to A\) 为连续映射。则存在子集 \(T\subset S\) 与整数族 \((n_p)_{p\in T}\)（其中 \(n_p\ge 1\)），使得 \(F\) 因子化为
\[
\mathfrak K_S
\longrightarrow
\prod_{p\in T}\ZZ/p^{n_p}\ZZ
\xrightarrow{\ \overline F\ }
A.
\]
换言之，任何有限输出的连续 prime-audit 都只依赖有限精度的 \(p\)-进截断；在有限输出与连续性同时成立时，不存在 genuinely 全精度的 prime-audit。
\end{theorem}
\begin{proof}
\(\mathfrak K_S\) 是紧全不连通 Hausdorff 群，而 \(A\) 为有限离散空间。故对每个 \(a\in A\)，纤维
\[
F^{-1}(a)
\]
都是紧开集。

取任意 \(a\in A\) 与任意 \(x=(x_p)_{p\in S}\in F^{-1}(a)\)。由乘积拓扑的基可知，存在整数族
\[
n_{x,p}\in \ZZ_{\ge 0}\qquad(p\in S)
\]
使得圆柱邻域
\[
U_x:=\prod_{p\in S}\bigl(x_p+p^{n_{x,p}}\ZZ_p\bigr)
\]
满足
\[
x\in U_x\subseteq F^{-1}(a).
\]
由于 \(F^{-1}(a)\) 紧，可从 \(\{U_x\}_{x\in F^{-1}(a)}\) 中抽取有限子覆盖
\[
F^{-1}(a)\subseteq U_{x^{(a,1)}}\cup\cdots\cup U_{x^{(a,r_a)}}.
\]
对全部 \(a\in A\) 与全部所选圆柱，把每个素数坐标上的指数取最大值：
\[
n_p:=\max_{a\in A}\max_{1\le j\le r_a} n_{x^{(a,j)},p}
\qquad(p\in S).
\]
再记
\[
T:=\{p\in S:\ n_p\ge 1\}.
\]
现设 \(y,z\in\mathfrak K_S\) 满足
\[
y\equiv z \pmod{p^{n_p}\ZZ_p}
\qquad(\forall\,p\in S).
\]
若 \(F(y)=a\)，则 \(y\in U_{x^{(a,j)}}\) 对某个 \(j\) 成立。由于
\[
n_p\ge n_{x^{(a,j)},p}
\qquad(\forall\,p\in S),
\]
并且 \(z\) 与 \(y\) 在每个坐标上满足更强的同余条件，必有
\[
z\in U_{x^{(a,j)}}\subseteq F^{-1}(a).
\]
故 \(F(z)=a=F(y)\)。这说明 \(F\) 在同一有限商
\[
\prod_{p\in T}\ZZ/p^{n_p}\ZZ
\]
的每条纤维上常值，从而存在唯一映射 \(\overline F\) 使上述因子化成立。证毕。
\end{proof}

\begin{corollary}[有限 prime-audit 的不可完备性具有任意指数级歧义]\label{cor:xi-localized-finite-prime-audit-arbitrary-exponential-ambiguity}
\leanverified{paper\_xi\_localized\_finite\_prime\_audit\_arbitrary\_exponential\_ambiguity}
固定一个经定理 \ref{thm:xi-localized-finite-prime-audit-cylinder-factorization} 因子化到有限圆柱商
\[
Q_T:=\prod_{p\in T}\ZZ/p^{n_p}\ZZ
\]
的 finite prime-audit，其中 \(T\subset\mathbb P\) 有限。则对任意 \(N\ge 1\)，都存在至少 \(2^N\) 个两两非同构的局部化群
\[
G_{S_E}:=\ZZ[S_E^{-1}],
\qquad
S_E:=T\cup\{q_i:\ i\in E\},
\qquad
E\subseteq\{1,\dots,N\},
\]
其中 \(q_1,\dots,q_N\in\mathbb P\setminus T\) 互异，并且这些局部化群都诱导同一 audited finite quotient \(Q_T\)。因而在给定 audit 下，它们不可区分。
\end{corollary}
\begin{proof}
任选互异素数
\[
q_1,\dots,q_N\in\mathbb P\setminus T,
\]
并对每个 \(E\subseteq\{1,\dots,N\}\) 定义
\[
S_E:=T\cup\{q_i:\ i\in E\},
\qquad
G_{S_E}:=\ZZ[S_E^{-1}].
\]
由定理 \ref{thm:killo-localization-circle-dimension-prime-ledger-splitting}，
\[
\widehat{G_{S_E}}
\]
的 prime-register 核恰为
\[
\mathfrak K_{S_E}=\prod_{p\in S_E}\ZZ_p
=
\left(\prod_{p\in T}\ZZ_p\right)\times
\left(\prod_{q_i\in E}\ZZ_{q_i}\right).
\]
而审计因子商 \(Q_T\) 只读取 \(T\) 中素数的有限截断，因此对全部 \(E\) 都有同一投影
\[
\mathfrak K_{S_E}\longrightarrow Q_T,
\]
额外加入的素数坐标 \(q_i\) 全部被遗忘。故这 \(2^N\) 个局部化群诱导完全相同的 audited finite quotient \(Q_T\)，从而在该 audit 下不可区分。

另一方面，若 \(E\neq E'\)，则 \(S_E\neq S_{E'}\)。由定理 \ref{thm:xi-localized-integers-padic-kernel-rigidity}，
\[
G_{S_E}\cong G_{S_{E'}}
\iff
S_E=S_{E'},
\]
故这些 \(G_{S_E}\) 两两非同构。由于子集总数恰为 \(2^N\)，结论成立。证毕。
\end{proof}

\endinput


\begin{theorem}[有限 prime-audit 的指数级不完备性]\label{thm:xi-localized-finite-prime-audit-exponential-incompleteness}
设 \(T\subset\mathbb P\) 为有限素数集。对任意 \(N\ge 1\)，存在至少 \(2^N\) 个两两非同构的有限素数局部化群
\[
G_E:=\ZZ[S_E^{-1}],
\qquad
S_E:=T\cup\{q_i:\ i\in E\},
\qquad
E\subseteq\{1,\dots,N\},
\]
其中 \(q_1,\dots,q_N\in\mathbb P\setminus T\) 互异，并且满足
\[
\cdim_{\mathrm{fr}}(G_E)=1,
\qquad
K_E^{(T)}:=\prod_{p\in S_E\cap T}\ZZ_p=\prod_{p\in T}\ZZ_p.
\]
因此，任何仅审计有限秩圆维与有限素数窗口 \(T\) 上可见 \(p\)-进核的制度，至少会把 \(2^N\) 个互不同构类压缩为同一个审计像。
\end{theorem}
\begin{proof}
任选互异素数
\[
q_1,\dots,q_N\in\mathbb P\setminus T.
\]
对每个子集 \(E\subseteq\{1,\dots,N\}\)，定义
\[
S_E:=T\cup\{q_i:\ i\in E\},
\qquad
G_E:=\ZZ[S_E^{-1}].
\]
由定理 \ref{thm:xi-localized-phase-sampling-closed-form}，对一切有限素数集 \(S\) 都有
\[
\cdim_{\mathrm{fr}}(\ZZ[S^{-1}])=1,
\]
故全部 \(G_E\) 具有相同的有限秩圆维。

另一方面，由构造
\[
S_E\cap T=T,
\]
于是
\[
K_E^{(T)}=\prod_{p\in S_E\cap T}\ZZ_p=\prod_{p\in T}\ZZ_p
\]
与 \(E\) 无关。故在有限素数窗口 \(T\) 上可见的 \(p\)-进核对这 \(2^N\) 个对象完全一致。

若 \(E\neq E'\)，则 \(S_E\neq S_{E'}\)。由定理 \ref{thm:xi-localized-integers-padic-kernel-rigidity}，
\[
G_E\cong G_{E'}
\iff
S_E=S_{E'},
\]
故这些 \(G_E\) 两两非同构。由于子集总数恰为 \(2^N\)，结论成立。证毕。
\end{proof}

\begin{theorem}[局部化对偶的同构刚性]\label{thm:xi-time-part56-localized-pontryagin-isomorphism-rigidity}
对任意有限素数集 \(S,T\subset \mathbb P\)，记
\[
G_S:=\ZZ[S^{-1}],
\qquad
\Sigma_S:=\widehat{G_S},
\]
则
\[
\boxed{\;
\Sigma_S\cong \Sigma_T
\quad\Longleftrightarrow\quad
S=T.\;}
\]
并且每个 \(\Sigma_S\) 都满足短正合列
\[
\boxed{\;
0\longrightarrow \prod_{p\in S}\ZZ_p
\longrightarrow \Sigma_S
\longrightarrow \TT
\longrightarrow 0,\;}
\]
同时
\[
\boxed{\;
\cdim_{\mathrm{fr}}(G_S)=1.\;}
\]
因此，所有有限素数局部化在连通圆维与有限秩扩展维数层上都塌缩到同一个值，而全部算术差异只能由全不连通的 \(p\)-进核恢复。
\end{theorem}
\begin{proof}
由定理 \ref{thm:xi-localized-phase-sampling-closed-form}，每个 \(G_S\) 的相位采样函数
\[
\Sigma^{\mathrm{ph}}_{G_S}(N)=\card{\operatorname{Hom}(G_S,\ZZ/N\ZZ)}
\]
唯一恢复素数集 \(S\)，并且同一定理已给出
\[
0\longrightarrow \prod_{p\in S}\ZZ_p
\longrightarrow \Sigma_S
\longrightarrow \TT
\longrightarrow 0,
\qquad
\cdim_{\mathrm{fr}}(G_S)=1.
\]
若 \(\Sigma_S\cong \Sigma_T\)，则由 Pontryagin 对偶的反变等价得到
\[
G_S\cong G_T.
\]
于是对每个 \(N\ge 1\) 都有
\[
\card{\operatorname{Hom}(G_S,\ZZ/N\ZZ)}
=
\card{\operatorname{Hom}(G_T,\ZZ/N\ZZ)},
\]
即 \(\Sigma^{\mathrm{ph}}_{G_S}=\Sigma^{\mathrm{ph}}_{G_T}\)。由相位采样闭式的唯一恢复性，必有 \(S=T\)。反向 implication 显然。证毕。
\end{proof}

\begin{corollary}[单圆维下的指数级局部化多样性]\label{cor:xi-time-part56-single-circle-dimension-exponential-diversity}
对每个整数 \(n\ge 1\)，都存在
\[
2^n
\]
个两两不同构的紧交换群，它们都具有相同的连通商 \(\TT\)，并且其 Pontryagin 对偶离散群都满足
\[
\cdim_{\mathrm{fr}}=1.
\]
\end{corollary}
\begin{proof}
取互异素数
\[
p_1,\dots,p_n,
\]
并对每个子集 \(E\subseteq\{1,\dots,n\}\) 定义
\[
S_E:=\{p_i:\ i\in E\},
\qquad
\Sigma_E:=\widehat{\ZZ[S_E^{-1}]}.
\]
由于子集总数为 \(2^n\)，只需证明这些 \(\Sigma_E\) 两两不同构。若 \(E\neq E'\)，则 \(S_E\neq S_{E'}\)，由定理 \ref{thm:xi-time-part56-localized-pontryagin-isomorphism-rigidity} 立得
\[
\Sigma_E\not\cong \Sigma_{E'}.
\]
另一方面，同一定理又给出每个 \(\Sigma_E\) 都满足短正合列
\[
0\longrightarrow \prod_{p\in S_E}\ZZ_p
\longrightarrow \Sigma_E
\longrightarrow \TT
\longrightarrow 0,
\]
故其连通商恒为 \(\TT\)，并且
\[
\cdim_{\mathrm{fr}}\!\bigl(\ZZ[S_E^{-1}]\bigr)=1.
\]
结论成立。证毕。
\end{proof}

\begin{theorem}[分层 prime-register 外置的最坏情形下界与自然扩张闭合]\label{thm:xi-layered-prime-register-sharp-natural-extension}
设 \((X_j)_{j=0}^{k}\) 为可数集合列，\(\pi_j:X_j\twoheadrightarrow X_{j+1}\) 为有限纤维满射，并假设对每个 \(j=0,\dots,k-1\) 都存在 \(y_j\in X_{j+1}\) 使
\[
\#\pi_j^{-1}(y_j)\ge 2.
\]
考虑任何分层外置
\[
\Psi(x_0)=(x_k,H(x_0)),
\qquad
H(x_0)=\prod_{j=0}^{k-1}h_j(x_0),
\]
其中每个 \(h_j(x_0)\) 的全部素因子均被限制在预先指定的 prime slice \(\mathcal P_j\)。若 \(\Psi\) 为单射，则最坏情形下每一层都必须有效占用对应的 prime slice；等价地，对每个 \(j\) 都存在微态 \(x_0^{(j)}\) 使
\[
h_j\bigl(x_0^{(j)}\bigr)\neq 1.
\]
另一方面，定理 \ref{thm:killo-layered-prime-slices-finite-depth} 给出恰在每层写入一枚素数的单射构造。故在“每层至少一处分支”的最坏情形下，分层 prime-register 外置所需的最小 prime-slice 数目恰等于层数 \(k\)。

特别地，若 \(T:X\twoheadrightarrow X\) 为非单射有限纤维满射，则任意深度 \(k\) 的前史消歧都至少需要 \(k\) 个有效 prime slices；而定理 \ref{thm:killo-natural-extension-branch-register} 的分支指标寄存器实现把整个无穷前史歧义共轭为右移插入系统，从而把上述有限深度最优外置在逆极限上精确封闭。
\end{theorem}
\begin{proof}
推论 \ref{cor:killo-layered-necessity-minimality} 断言：在每层都存在分支点的前提下，任何单射分层外置都迫使每个 prime slice \(\mathcal P_j\) 在某个微态上非平凡使用。故最坏情形下至少需要 \(k\) 个彼此独立的有效 prime slices。
\par
反向上，定理 \ref{thm:killo-layered-prime-slices-finite-depth} 已构造出单射映射
\[
\Phi(x_0)=(x_k,G(x_0)),
\qquad
G(x_0)=\prod_{j=0}^{k-1}q_j(x_0),
\qquad
q_j(x_0)\in\mathcal P_j,
\]
并且每层恰写入一枚素数。故上述下界可达，因而最小 prime-slice 数目正是 \(k\)。
\par
若 \(T:X\twoheadrightarrow X\) 为非单射有限纤维满射，则存在 \(y\in X\) 使 \(\#T^{-1}(y)\ge 2\)。对任意固定深度 \(k\)，取 \(X_j:=X\) 且 \(\pi_j:=T\)（\(0\le j\le k-1\)），便满足前述最坏情形假设，故至少需要 \(k\) 个有效 prime slices。
最后，定理 \ref{thm:killo-natural-extension-branch-register} 把自然扩张
\(
X^{\leftarrow}
\)
与右移插入系统
\[
(x,(a_1,a_2,\dots))\longmapsto (T(x),(\beta(x),a_1,a_2,\dots))
\]
共轭，从而说明有限深度上的最优分层外置在逆极限上恰由规范分支指标寄存器闭合。证毕。
\end{proof}

\paragraph{\texorpdfstring{$S$}{S}-solenoid 的泛终性质与 prime-register 轴的函子恢复}
在局部化数系对偶 \(\Sigma_S=\widehat{G_S}\) 的刚性、短正合列
\[
0\longrightarrow \prod_{p\in S}\ZZ_p\longrightarrow \Sigma_S\longrightarrow \TT\longrightarrow 0
\]
以及 \(\Sigma_S\cong \Sigma_T\iff S=T\) 的判据已经确定之后，还可把“prime-register 轴并非实现偶然，而是结构强制不变量”压缩为一条范畴论表述：在固定 \(S\)-局部化纪律的实现范畴内，\(\Sigma_S\) 是唯一的终对象，而其账本核的 \(p\)-主分量恰把素数集 \(S\) 函子地恢复出来；对偶地，离散侧的 \(G_S=\ZZ[S^{-1}]\) 则是同一局部化纪律下的初对象。

\begin{theorem}[\texorpdfstring{$S$}{S}-solenoid 的终对象泛性质与 prime-register 轴的函子恢复]\label{thm:xi-time-part56-solenoid-terminal-prime-register-recovery}
\leanverified{paper\_xi\_time\_part56\_solenoid\_terminal\_prime\_register\_recovery}
固定有限素数集 \(S\subset\mathbb P\)，并记
\[
G_S:=\ZZ[S^{-1}],
\qquad
\Sigma_S:=\widehat{G_S}.
\]
记 \(\pi_S:\Sigma_S\twoheadrightarrow\TT\) 为标准投影，并定义范畴 \(\mathcal C_S\) 如下：
\begin{enumerate}
  \item 对象为一对 \((K,\pi)\)，其中 \(K\) 为紧交换群，\(\pi:K\twoheadrightarrow\TT\) 为连续满射；
  \item 若记 \(A:=\widehat K\) 且 \(j:\ZZ\hookrightarrow A\) 为 \(\pi\) 的对偶嵌入，则要求对每个 \(p\in S\)，乘法映射
  \[
  [p]_A:A\longrightarrow A
  \]
  为双射；
  \item 态射 \(f:(K,\pi)\to(K',\pi')\) 为满足
  \[
  \pi' \circ f=\pi
  \]
  的连续群同态。
\end{enumerate}
则成立：
\begin{enumerate}
  \item \((\Sigma_S,\pi_S)\) 在 \(\mathcal C_S\) 中为终对象；等价地，对任意 \((K,\pi)\in\mathcal C_S\)，都存在唯一连续群同态
  \[
  f:K\longrightarrow \Sigma_S
  \]
  使
  \[
  \pi_S\circ f=\pi.
  \]
  \item 其账本核
  \[
  K_S^{\mathrm{reg}}:=\ker(\pi_S)
  \]
  满足
  \[
  K_S^{\mathrm{reg}}\cong \prod_{p\in S}\ZZ_p;
  \]
  \item 因而素数集 \(S\) 可由终对象的账本核函子地恢复为
  \[
  \boxed{\;
  S=\{\,p\in\mathbb P:\ (K_S^{\mathrm{reg}})_p\neq 0\,\},\;}
  \]
  其中 \((K_S^{\mathrm{reg}})_p\) 表示其 Sylow pro-\(p\) 分量。
\end{enumerate}
\end{theorem}

\begin{proof}
第 \(1\) 条正是局部化的泛性质在紧侧的对偶表述，而其离散对偶表述正是 \(G_S=\ZZ[S^{-1}]\) 在含 \(\ZZ\) 的 \(S\)-局部离散群范畴中具有初对象性质。
确切地说，定理 \ref{thm:cdim-s-solenoid-terminal-object} 已证明：对任意对象 \((K,\pi)\in\mathcal C_S\)，存在唯一连续群同态
\[
f:K\longrightarrow \Sigma_S
\]
使
\[
\pi_S\circ f=\pi.
\]
故 \((\Sigma_S,\pi_S)\) 为终对象。

第 \(2\) 条由定理 \ref{thm:xi-localized-phase-sampling-closed-form} 与定理 \ref{thm:xi-time-part56-localized-pontryagin-isomorphism-rigidity} 已给出的标准短正合列直接得到：
\[
0\longrightarrow \prod_{p\in S}\ZZ_p
\longrightarrow \Sigma_S
\xrightarrow{\ \pi_S\ }
\TT
\longrightarrow 0.
\]
因此
\[
\ker(\pi_S)\cong \prod_{p\in S}\ZZ_p.
\]

对第 \(3\) 条，\(\prod_{p\in S}\ZZ_p\) 的 Sylow pro-\(p\) 分量在且仅在 \(p\in S\) 时非平凡；
而不同素数坐标彼此正交，故核的 \(p\)-主分量结构唯一决定 \(S\)。
这正说明 prime-register 轴并非额外实现选择，而是由终对象的核函子强制记录的结构不变量。证毕。
\end{proof}


\paragraph{冻结 escort 谱、全微态位率阈值与三重算术独立的乘法密度}
以下条目把固定冻结相中的 escort Rényi 熵谱与速率压力差公式、bin-fold 本征推前的绝对 Rényi 常数项、dyadic 容量临界窗与二相自由能变分、全微态精确反演所需比特率的有限层闭式，以及分枝三次场、零温四次场与 RH 临界五次场之间的算术独立性，压缩为一组可直接调用的终端结论。

\begin{theorem}[escort Rényi 速率的压力差公式、双阈值塌缩与 \texorpdfstring{$q=0$}{q=0} 跃迁]\label{thm:xi-fixed-freezing-escort-renyi-spectrum-collapse}
\leanverified{paper\_xi\_fixed\_freezing\_escort\_renyi\_spectrum\_collapse}
沿用结论 \ref{con:xi-terminal-fixed-freezing-escort-oracle-sharp-threshold} 的记号。记
\[
a_{\mathrm c}:=\frac{\log\varphi-g_*}{\alpha_*-\alpha_*^{(2)}}<\infty,
\]
并记
\[
S_t(m):=\sum_{x\in X_m}d_m(x)^t
\qquad (t>0).
\]
对任意固定 \(a>0\) 定义 escort 分布
\[
\pi_{m,a}(x):=\frac{d_m(x)^a}{S_a(m)},
\qquad
A_m^{\max}:=\{x\in X_m:\ d_m(x)=M_m\},
\qquad
K_m:=|A_m^{\max}|.
\]
若极限
\[
P_t:=\lim_{m\to\infty}\frac1m\log S_t(m)
\]
存在，则称其为 \(t\)-阶压力。
记 \(H_q\) 为以自然对数定义的 Rényi 熵族，其中
\[
H_q(\mu):=\frac{1}{1-q}\log\sum_x \mu(x)^q\quad(q\in(0,\infty)\setminus\{1\}),
\]
\[
H_1(\mu):=-\sum_x \mu(x)\log\mu(x),
\qquad
H_\infty(\mu):=-\log\max_x \mu(x),
\qquad
H_0(\mu):=\log|\supp \mu|.
\]
则成立：
\begin{enumerate}
  \item 对任意 \(q\in(0,\infty)\setminus\{1\}\)，若 \(P_a\) 与 \(P_{aq}\) 存在，则
  \[
  \boxed{\;
  \lim_{m\to\infty}\frac1m H_q(\pi_{m,a})
  =
  \frac{P_{aq}-qP_a}{1-q}.\;}
  \]
  \item 若进一步落在双阈值区
  \[
  a>a_{\mathrm c},
  \qquad
  aq>a_{\mathrm c},
  \]
  则对每个 \(q\in(0,\infty)\setminus\{1\}\) 都有
  \[
  \lim_{m\to\infty}\frac1m H_q(\pi_{m,a})=g_*.
  \]
  当 \(q>1\) 时，第二个阈值由 \(aq>a>a_{\mathrm c}\) 自动成立。
  \item 对 \(q\in\{1,\infty\}\)，若 \(a>a_{\mathrm c}\)，则同样有
  \[
  \lim_{m\to\infty}\frac1m H_q(\pi_{m,a})=g_*.
  \]
  \item 若 \(a>a_{\mathrm c}\) 且记
  \[
  U_m^{\max}:=\mathrm{Unif}(A_m^{\max}),
  \]
  则总变差距离满足精确恒等式
  \[
  \boxed{\;
  \|\pi_{m,a}-U_m^{\max}\|_{\mathrm{TV}}
  =
  1-\pi_{m,a}(A_m^{\max})
  \le
  \exp\!\bigl(-m\Delta(a)+o(m)\bigr),\;}
  \]
  其中
  \[
  \Delta(a):=a(\alpha_*-\alpha_*^{(2)})-(\log\varphi-g_*)>0.
  \]
  \item Hartley 熵率不塌缩，而满足
  \[
  \boxed{\;
  \lim_{m\to\infty}\frac1m H_0(\pi_{m,a})
  =
  \lim_{m\to\infty}\frac1m\log|X_m|
  =
  \log\varphi.\;}
  \]
\end{enumerate}
因此，正阶 Rényi 速率先由压力差公式统一表示；一旦 \((a,aq)\) 同时越过冻结阈值，全部正阶 Rényi 熵率便共同塌缩到 \(g_*\)；而当 \(\log\varphi>g_*\) 时，Hartley 熵率仍停留在 \(\log\varphi\)，从而在 \(q=0\) 处发生真实跃迁。
\end{theorem}
\begin{proof}
先取 \(q\in(0,\infty)\setminus\{1\}\)。由定义
\[
\sum_{x\in X_m}\pi_{m,a}(x)^q
=
\frac{S_{aq}(m)}{S_a(m)^q},
\]
故
\[
H_q(\pi_{m,a})
=
\frac{1}{1-q}\Bigl(\log S_{aq}(m)-q\log S_a(m)\Bigr).
\]
若 \(P_a\) 与 \(P_{aq}\) 存在，两边同除以 \(m\) 并令 \(m\to\infty\)，便得到第一条。

若进一步 \(a>a_{\mathrm c}\) 且 \(aq>a_{\mathrm c}\)，则定理 \ref{thm:fold-pressure-freezing-threshold} 同时适用于指数 \(a\) 与 \(aq\)，从而
\[
\frac1m\log S_a(m)\to a\alpha_*+g_*,
\qquad
\frac1m\log S_{aq}(m)\to aq\alpha_*+g_*.
\]
代入即得
\[
\frac1m H_q(\pi_{m,a})
\to
\frac{(aq\alpha_*+g_*)-q(a\alpha_*+g_*)}{1-q}
=
g_*.
\]
这便证明了第二条。
\par
对 \(q=\infty\)，推论 \ref{cor:fold-escort-minentropy-rate} 已给出
\[
\lim_{m\to\infty}\frac1m H_\infty(\pi_{m,a})=h_\infty(a)=g_*.
\]
这给出第三条中 \(q=\infty\) 的情形。
\par
再证 Shannon 情形 \(q=1\)。记
\[
p_m:=\pi_{m,a}(A_m^{\max}).
\]
由于 \(d_m(x)=M_m\) 在 \(A_m^{\max}\) 上常值，\(\pi_{m,a}\) 在该集合上的条件分布本身就是均匀分布 \(U_m^{\max}\)。因此
\[
\|\pi_{m,a}-U_m^{\max}\|_{\mathrm{TV}}
=
\frac12\sum_{x\in A_m^{\max}}\left|\frac{p_m}{K_m}-\frac1{K_m}\right|
+\frac12\sum_{x\notin A_m^{\max}}\pi_{m,a}(x)
=
1-p_m.
\]
再由定理 \ref{thm:fold-escort-groundstate-concentration}，
\[
1-p_m\le \exp\!\bigl(-m\Delta(a)+o(m)\bigr),
\]
即得第四条中的总变差公式。
\par
若 \(p_m<1\)，记 \(\eta_m\) 为 \(\pi_{m,a}\) 在补集 \(X_m\setminus A_m^{\max}\) 上的条件分布；若 \(p_m=1\)，任取 \(\eta_m\) 即可。于是
\[
\pi_{m,a}=p_m\,U_m^{\max}+(1-p_m)\eta_m
\]
且两项支撑互不相交，从而 Shannon 熵满足精确分解
\[
H(\pi_{m,a})
=
h_{\mathrm{bin}}(p_m)+p_m\log K_m+(1-p_m)H(\eta_m),
\]
其中
\[
h_{\mathrm{bin}}(t):=-t\log t-(1-t)\log(1-t).
\]
又因为
\[
0\le H(\eta_m)\le \log|X_m|=\log F_{m+2}=m\log\varphi+O(1),
\]
并且 \(1-p_m\) 指数衰减，故
\[
h_{\mathrm{bin}}(p_m)=o(m),
\qquad
(1-p_m)H(\eta_m)=o(m),
\qquad
p_m\log K_m=\log K_m+o(m).
\]
于是
\[
H(\pi_{m,a})=\log K_m+o(m),
\]
从而
\[
\lim_{m\to\infty}\frac1m H(\pi_{m,a})=g_*.
\]
这便完成了第三条中 \(q=1\) 的情形。
\par
最后，由定义对一切 \(x\in X_m\) 都有 \(d_m(x)\ge 1\)，故 \(\pi_{m,a}(x)>0\) 在整个 \(X_m\) 上成立。因此
\[
H_0(\pi_{m,a})=\log|\supp(\pi_{m,a})|=\log|X_m|=\log F_{m+2},
\]
从而
\[
\lim_{m\to\infty}\frac1m H_0(\pi_{m,a})=\log\varphi.
\]
这正是第五条。
证毕。
\end{proof}

\begin{corollary}[冻结相中 Rényi 阶阈值的显式闭式]\label{cor:xi-fixed-freezing-renyi-critical-order-threshold}
\leanverified{paper\_xi\_fixed\_freezing\_renyi\_critical\_order\_threshold}
沿用定理 \ref{thm:xi-fixed-freezing-escort-renyi-spectrum-collapse} 的记号。对固定
\[
a>a_{\mathrm c},
\]
定义
\[
r_{\mathrm c}(a):=\frac{a_{\mathrm c}}{a}\in(0,1).
\]
则对任意
\[
r\in[r_{\mathrm c}(a),\infty]
\]
都有
\[
\lim_{m\to\infty}\frac1m H_r(\pi_{m,a})=g_*.
\]
另一方面，
\[
\lim_{m\to\infty}\frac1m H_0(\pi_{m,a})=\log\varphi.
\]
因此，冻结相内部的 Rényi 谱在阶参数上具有显式阈值 \(r_{\mathrm c}(a)\)：阈值以上的全部信息型熵率统一塌缩到基态简并指数 \(g_*\)，而支撑熵率仍停留在环境熵率 \(\log\varphi\)。
\end{corollary}
\begin{proof}
对 \(r>r_{\mathrm c}(a)\) 且 \(r\notin\{1,\infty\}\)，由于
\[
ar>a_{\mathrm c},
\]
定理 \ref{thm:xi-fixed-freezing-escort-renyi-spectrum-collapse} 的第二条直接给出
\[
\lim_{m\to\infty}\frac1m H_r(\pi_{m,a})=g_*.
\]
对 \(r\in\{1,\infty\}\)，同一定理第三条给出同样结论；对 \(r=0\)，第五条给出
\[
\lim_{m\to\infty}\frac1m H_0(\pi_{m,a})=\log\varphi.
\]
因此只需处理临界点
\[
r=r_{\mathrm c}(a)=\frac{a_{\mathrm c}}{a}\in(0,1).
\]

记
\[
\varepsilon_m:=1-\pi_{m,a}(A_m^{\max}).
\]
由定理 \ref{thm:xi-fixed-freezing-escort-renyi-spectrum-collapse} 的第四条，
\[
\varepsilon_m\le \exp\!\bigl(-m\Delta(a)+o(m)\bigr),
\qquad
\Delta(a):=a(\alpha_*-\alpha_*^{(2)})-(\log\varphi-g_*)>0.
\]
并且 \(\pi_{m,a}\) 在 \(A_m^{\max}\) 上均匀，故
\[
\sum_{x\in A_m^{\max}}\pi_{m,a}(x)^r
=
K_m\left(\frac{1-\varepsilon_m}{K_m}\right)^r
=
K_m^{1-r}(1-\varepsilon_m)^r.
\]
另一方面，由 \(0<r<1\) 时的凹性不等式
\[
\sum_{i=1}^{N}b_i^r\le N^{1-r}\Bigl(\sum_{i=1}^{N}b_i\Bigr)^r
\qquad (b_i\ge 0),
\]
可得
\[
\sum_{x\notin A_m^{\max}}\pi_{m,a}(x)^r
\le
|X_m|^{1-r}\varepsilon_m^r.
\]
由于 \(|X_m|=\exp(m\log\varphi+o(m))\)，上式右端满足
\[
|X_m|^{1-r}\varepsilon_m^r
\le
\exp\!\bigl(m[(1-r)\log\varphi-r\Delta(a)]+o(m)\bigr).
\]
当 \(r=r_{\mathrm c}(a)\) 时，由
\[
ra(\alpha_*-\alpha_*^{(2)})=a_{\mathrm c}(\alpha_*-\alpha_*^{(2)})=\log\varphi-g_*
\]
可得
\[
r\Delta(a)=r\bigl[a(\alpha_*-\alpha_*^{(2)})-(\log\varphi-g_*)\bigr]
=(1-r)(\log\varphi-g_*),
\]
故
\[
(1-r)\log\varphi-r\Delta(a)=(1-r)g_*.
\]
同时
\[
K_m^{1-r}(1-\varepsilon_m)^r
=
\exp\!\bigl((1-r)g_*\,m+o(m)\bigr).
\]
于是
\[
\exp\!\bigl((1-r)g_*\,m+o(m)\bigr)
\le
\sum_x\pi_{m,a}(x)^r
\le
2\exp\!\bigl((1-r)g_*\,m+o(m)\bigr),
\]
从而
\[
\sum_x\pi_{m,a}(x)^r
=
\exp\!\bigl((1-r)g_*\,m+o(m)\bigr).
\]
两边取对数并除以 \(1-r\)，便得到
\[
H_r(\pi_{m,a})=g_*\,m+o(m),
\]
即
\[
\lim_{m\to\infty}\frac1m H_r(\pi_{m,a})=g_*.
\]
证毕。
\end{proof}

\begin{corollary}[冻结 escort 态到基态均匀分布的精确总变差公式]\label{cor:xi-fixed-freezing-escort-groundstate-tv-identity}
\leanverified{paper\_xi\_fixed\_freezing\_escort\_groundstate\_tv\_identity}
沿用定理 \ref{thm:xi-fixed-freezing-escort-renyi-spectrum-collapse} 的记号。对任意整数
\[
a>a_{\mathrm c},
\]
记
\[
U_m^{\max}:=\mathrm{Unif}(A_m^{\max}).
\]
则有精确恒等式
\[
\|\pi_{m,a}-U_m^{\max}\|_{\mathrm{TV}}
=
1-\pi_{m,a}(A_m^{\max}),
\]
并且
\[
\|\pi_{m,a}-U_m^{\max}\|_{\mathrm{TV}}
\le
\exp\!\bigl(-m\Delta(a)+o(m)\bigr).
\]
\end{corollary}
\begin{proof}
这正是定理 \ref{thm:xi-fixed-freezing-escort-renyi-spectrum-collapse} 第四条在整数
\(
a>a_{\mathrm c}
\)
上的直接重述。证毕。
\end{proof}

\begin{theorem}[冻结 escort 态对一切有界可观测量的指数统计塌缩]\label{thm:xi-fixed-freezing-escort-bounded-observable-collapse}
沿用推论 \ref{cor:xi-fixed-freezing-escort-groundstate-tv-identity} 的记号。对任意固定
\[
a>a_{\mathrm c}
\]
与任意有界可观测量 \(f_m:X_m\to\CC\)，都有
\[
\boxed{\;
\left|
\EE_{\pi_{m,a}}f_m-\EE_{U_m^{\max}}f_m
\right|
\le
2\|f_m\|_\infty
\exp\!\bigl(-m\Delta(a)+o(m)\bigr).\;}
\]
若进一步 \(f_m\) 在 \(A_m^{\max}\) 上为常数 \(c_m\)，则有精确恒等式
\[
\boxed{\;
\EE_{\pi_{m,a}}f_m-\EE_{U_m^{\max}}f_m
=
\sum_{x\notin A_m^{\max}}\pi_{m,a}(x)\bigl(f_m(x)-c_m\bigr).\;}
\]
因此在冻结区内，一切有界可观测量都以指数速度塌缩到基态均匀分布的统计值；而对基态层上常值的可观测量，全部误差恰由基态外质量单独承担。
\leanverified{paper\_xi\_fixed\_freezing\_escort\_bounded\_observable\_collapse}
\end{theorem}
\begin{proof}
由总变差距离与有界测试函数的标准对偶不等式，
\[
\left|
\EE_{\pi_{m,a}}f_m-\EE_{U_m^{\max}}f_m
\right|
\le
2\|f_m\|_\infty
\|\pi_{m,a}-U_m^{\max}\|_{\mathrm{TV}}.
\]
再由推论 \ref{cor:xi-fixed-freezing-escort-groundstate-tv-identity}，
\[
\|\pi_{m,a}-U_m^{\max}\|_{\mathrm{TV}}
\le
\exp\!\bigl(-m\Delta(a)+o(m)\bigr),
\]
便得到第一条估计。

若 \(f_m|_{A_m^{\max}}\equiv c_m\)，则
\[
\EE_{U_m^{\max}}f_m=c_m.
\]
另一方面，由于 \(\pi_{m,a}\) 在 \(A_m^{\max}\) 上本身均匀，故
\[
\sum_{x\in A_m^{\max}}\pi_{m,a}(x)f_m(x)
=
\pi_{m,a}(A_m^{\max})\,c_m.
\]
于是
\[
\begin{aligned}
\EE_{\pi_{m,a}}f_m-\EE_{U_m^{\max}}f_m
&=
\sum_{x\in A_m^{\max}}\pi_{m,a}(x)f_m(x)
+
\sum_{x\notin A_m^{\max}}\pi_{m,a}(x)f_m(x)-c_m\\
&=
\bigl(\pi_{m,a}(A_m^{\max})-1\bigr)c_m
+
\sum_{x\notin A_m^{\max}}\pi_{m,a}(x)f_m(x)\\
&=
\sum_{x\notin A_m^{\max}}\pi_{m,a}(x)\bigl(f_m(x)-c_m\bigr).
\end{aligned}
\]
证毕。
\end{proof}

\begin{corollary}[冻结相内部不同 escort 态的指数全变差趋同]\label{cor:xi-fixed-freezing-escort-pairwise-tv-coalescence}
沿用定理 \ref{thm:xi-fixed-freezing-escort-renyi-spectrum-collapse} 的记号。对任意固定
\[
a,b>a_{\mathrm c},
\]
都有
\[
\|\pi_{m,a}-\pi_{m,b}\|_{\mathrm{TV}}
\le
\exp\!\bigl(-m\Delta(a)+o(m)\bigr)
+
\exp\!\bigl(-m\Delta(b)+o(m)\bigr).
\]
特别地，
\[
\|\pi_{m,a}-\pi_{m,b}\|_{\mathrm{TV}}
\le
\exp\!\bigl(-m\min\{\Delta(a),\Delta(b)\}+o(m)\bigr).
\]
\end{corollary}
\begin{proof}
记
\[
U_m^{\max}:=\mathrm{Unif}(A_m^{\max}).
\]
由三角不等式，
\[
\|\pi_{m,a}-\pi_{m,b}\|_{\mathrm{TV}}
\le
\|\pi_{m,a}-U_m^{\max}\|_{\mathrm{TV}}
+
\|U_m^{\max}-\pi_{m,b}\|_{\mathrm{TV}}.
\]
再由定理 \ref{thm:xi-fixed-freezing-escort-renyi-spectrum-collapse} 的第四条，
\[
\|\pi_{m,a}-U_m^{\max}\|_{\mathrm{TV}}
\le
\exp\!\bigl(-m\Delta(a)+o(m)\bigr),
\]
\[
\|U_m^{\max}-\pi_{m,b}\|_{\mathrm{TV}}
\le
\exp\!\bigl(-m\Delta(b)+o(m)\bigr).
\]
相加即得第一条。第二条由
\[
u_m+v_m
\le
2\max\{u_m,v_m\}
\]
应用于
\[
u_m:=\exp\!\bigl(-m\Delta(a)+o(m)\bigr),
\qquad
v_m:=\exp\!\bigl(-m\Delta(b)+o(m)\bigr)
\]
后直接得到。证毕。
\end{proof}


\begin{theorem}[全微态精确反演的有限层比特阈值与线性速率]\label{thm:xi-full-microstate-exact-inversion-bitrate-threshold}
\leanverified{paper\_xi\_full\_microstate\_exact\_inversion\_bitrate\_threshold}
对整数预算 \(B\ge 0\)，定义 dyadic 预言机容量
\[
C_m(B):=\sum_{x\in X_m}\min\bigl(d_m(x),2^B\bigr),
\]
并记最大纤维规模
\[
M_m:=\max_{x\in X_m}d_m(x).
\]
定义全微态精确反演所需的最小比特数
\[
B_m^{\mathrm{full}}
:=
\min\{B\in\ZZ_{\ge 0}: C_m(B)=2^m\}.
\]
则对每个 \(m\) 都有严格恒等式
\[
\boxed{\;
B_m^{\mathrm{full}}=\left\lceil \log_2 M_m\right\rceil.\;}
\]
特别地，若
\[
\alpha_*=\lim_{m\to\infty}\frac1m\log M_m
\]
存在，则
\[
\boxed{\;
\lim_{m\to\infty}\frac{B_m^{\mathrm{full}}}{m}
=
\frac{\alpha_*}{\log 2}.\;}
\]
更进一步，对任意 \(\varepsilon>0\) 与任意整数序列 \(B_m\ge 0\)，成立：
\begin{enumerate}
  \item 若
  \[
  B_m\ge \left(\frac{\alpha_*}{\log2}+\varepsilon\right)m
  \]
  对充分大 \(m\) 成立，则最终
  \[
  C_m(B_m)=2^m.
  \]
  \item 若
  \[
  B_m\le \left(\frac{\alpha_*}{\log2}-\varepsilon\right)m
  \]
  对充分大 \(m\) 成立，则
  \[
  \boxed{\;
  2^m-C_m(B_m)\ \ge\ K_m\bigl(M_m-2^{B_m}\bigr),\;}
  \]
  其中 \(K_m:=|A_m^{\max}|\)。若再假设
  \[
  g_*=\lim_{m\to\infty}\frac1m\log K_m
  \]
  存在，则
  \[
  \boxed{\;
  \liminf_{m\to\infty}\frac1m\log\bigl(2^m-C_m(B_m)\bigr)
  \ge
  \alpha_*+g_*.\;}
  \]
\end{enumerate}
因此 \(\alpha_*/\log 2\) 不是近似阈值，而是全微态精确反演的精确信息速率阈值。
\end{theorem}
\begin{proof}
由定理 \ref{thm:fold-oracle-capacity-closed-form}，
\[
C_m(B)=\sum_{x\in X_m}\min\bigl(d_m(x),2^B\bigr).
\]
注意到
\[
\sum_{x\in X_m}d_m(x)=2^m.
\]
因此 \(C_m(B)=2^m\) 当且仅当每一项都未被截断，即
\[
\min(d_m(x),2^B)=d_m(x)\qquad(\forall x\in X_m),
\]
等价于
\[
2^B\ge d_m(x)\qquad(\forall x\in X_m),
\]
也即
\[
2^B\ge M_m.
\]
满足此条件的最小非负整数 \(B\) 正是 \(\lceil \log_2 M_m\rceil\)，故得到第一式。
\par
若 \(\alpha_*=\lim m^{-1}\log M_m\) 存在，则
\[
\frac1m\log_2 M_m=\frac{1}{m\log2}\log M_m\to \frac{\alpha_*}{\log2}.
\]
而上取整仅引入 \(O(1)\) 误差，故
\[
\frac{B_m^{\mathrm{full}}}{m}\to \frac{\alpha_*}{\log2}.
\]
\par
若
\[
B_m\ge \left(\frac{\alpha_*}{\log2}+\varepsilon\right)m,
\]
则
\[
2^{B_m}\ge \exp\!\bigl((\alpha_*+\varepsilon\log2)m\bigr).
\]
由 \(M_m=\exp(\alpha_*m+o(m))\)，对充分大 \(m\) 必有 \(2^{B_m}\ge M_m\)，于是 \(C_m(B_m)=2^m\)。
\par
反之，若
\[
B_m\le \left(\frac{\alpha_*}{\log2}-\varepsilon\right)m,
\]
则
\[
2^m-C_m(B_m)
=
\sum_{x\in X_m}\bigl(d_m(x)-2^{B_m}\bigr)_+,
\]
其中 \((t)_+:=\max\{t,0\}\)。对每个 \(x\in A_m^{\max}\) 都有 \(d_m(x)=M_m\)，故
\[
2^m-C_m(B_m)\ge \sum_{x\in A_m^{\max}}\bigl(M_m-2^{B_m}\bigr)
=
K_m\bigl(M_m-2^{B_m}\bigr).
\]
这便得到所示有限层下界。
\par
若 \(g_*\) 存在，则 \(K_m=\exp(g_*m+o(m))\)。同时由预算假设
\[
2^{B_m}\le \exp\!\bigl((\alpha_*-\varepsilon\log2)m\bigr),
\]
故
\[
M_m-2^{B_m}
=
\exp(\alpha_*m+o(m))
\left(1-\exp\!\bigl(-\varepsilon(\log2)m+o(m)\bigr)\right)
=
\exp(\alpha_*m+o(m)).
\]
于是
\[
K_m\bigl(M_m-2^{B_m}\bigr)=\exp\!\bigl((\alpha_*+g_*)m+o(m)\bigr),
\]
从而
\[
\liminf_{m\to\infty}\frac1m\log\bigl(2^m-C_m(B_m)\bigr)\ge \alpha_*+g_*.
\]
证毕。
\end{proof}

\paragraph{bin-fold 绝对熵常数、容量临界窗、阈值预算与上界刚性}
在前述冻结 escort 熵率、全微态位率阈值之外，bin-fold 本征推前自身的绝对 Rényi 常数项、dyadic 预算下的临界成功率窗口、阈值预算的常数跃迁、最大纤维组合上界的指数精确性，以及其零温修正的二相变分结构，也可由既有两态渐近与 escort 末位律完全闭合。

\begin{theorem}[bin-fold 本征推前的绝对 Rényi 熵常数项闭式]\label{thm:xi-foldbin-absolute-renyi-entropy-constant-closed}
\leanverified{paper\_xi\_foldbin\_absolute\_renyi\_entropy\_constant\_closed}
令
\[
\mu_m(w):=\frac{d_m^{\mathrm{bin}}(w)}{2^m},
\qquad
w\in X_m.
\]
则对任意固定 \(q>0\)、\(q\neq 1\)，有
\[
H_q(\mu_m)
=
m\log\varphi
+
\frac{1}{1-q}\log\!\Bigl(\frac{\varphi+\varphi^{-q}}{\sqrt5}\Bigr)
+
O\!\Bigl(\Bigl(\frac{\varphi}{2}\Bigr)^m\Bigr).
\]
相应的 Shannon 熵满足
\[
H(\mu_m)
=
m\log\varphi
+
\frac{\log\varphi}{1+\varphi^2}
+
O\!\Bigl(\Bigl(\frac{\varphi}{2}\Bigr)^m\Bigr).
\]
\end{theorem}
\begin{proof}
在定理 \ref{thm:xi-fold-escort-renyi-shannon-constants} 中取 escort 参数为 \(1\)，并把其中的 Rényi 阶改记为 \(q\)，即可得到
\[
H_q(\mu_m)
=
m\log\varphi
+
\frac{1}{1-q}\left[
\log\!\Bigl(\frac{\varphi+\varphi^{-q}}{\sqrt5}\Bigr)
-q\log\!\Bigl(\frac{\varphi+\varphi^{-1}}{\sqrt5}\Bigr)
\right]
+
O\!\Bigl(\Bigl(\frac{\varphi}{2}\Bigr)^m\Bigr).
\]
而 \(\varphi+\varphi^{-1}=\sqrt5\)，故第二项化简为所示闭式。Shannon 情形正是推论 \ref{cor:xi-foldbin-shannon-gauge-constant-mirror} 的内容；再用恒等式 \(1+\varphi^2=\varphi\sqrt5\) 即得此处写法。证毕。
\end{proof}

\begin{theorem}[bin-fold dyadic 预言机容量的临界窗口极限]\label{thm:xi-foldbin-dyadic-capacity-critical-window-limit}
\leanverified{paper\_xi\_foldbin\_dyadic\_capacity\_critical\_window\_limit}
对整数预算 \(B\ge 0\)，定义 dyadic 预言机容量与最优成功率
\[
C_m^{\mathrm{bin}}(B):=\sum_{w\in X_m}\min\!\bigl(d_m^{\mathrm{bin}}(w),2^B\bigr),
\qquad
\mathrm{Succ}_m^{\mathrm{bin}}(B):=\frac{C_m^{\mathrm{bin}}(B)}{2^m}=P_m^\star(2^B).
\]
记
\[
\lambda:=\frac{2}{\varphi}.
\]
若整数序列 \((B_m)_{m\ge 1}\) 满足
\[
\frac{2^{B_m}}{\lambda^m}\longrightarrow \zeta\in(0,\infty),
\]
则极限
\[
\lim_{m\to\infty}\mathrm{Succ}_m^{\mathrm{bin}}(B_m)
\]
存在，且等于
\[
\mathcal S_{\mathrm{bin}}(\zeta)
:=
\begin{cases}
\dfrac{\varphi^2}{\sqrt5}\,\zeta, & 0<\zeta\le \varphi^{-1},\\[3mm]
\dfrac{\varphi}{\sqrt5}\,\zeta+\dfrac{1}{\varphi\sqrt5}, & \varphi^{-1}\le \zeta\le 1,\\[3mm]
1, & \zeta\ge 1.
\end{cases}
\]
\end{theorem}
\begin{proof}
记
\[
A_b^{(m)}:=\{w\in X_m:\ w_m=b\},
\qquad
b\in\{0,1\}.
\]
由定理 \ref{thm:fold-bin-two-state-asymptotic}，存在常数 \(C>0\) 使得对充分大 \(m\) 与任意 \(w\in X_m\) 都有
\[
d_m^{\mathrm{bin}}(w)
=
\lambda^m\varphi^{-w_m}\bigl(1+\varepsilon_{m,w}\bigr),
\qquad
|\varepsilon_{m,w}|\le C\Bigl(\frac{\varphi}{2}\Bigr)^m.
\]
又由稳定词计数，
\[
|A_0^{(m)}|=F_{m+1},
\qquad
|A_1^{(m)}|=F_m.
\]

若 \(0<\zeta<\varphi^{-1}\)，则对充分大 \(m\) 有
\[
2^{B_m}<\lambda^m\varphi^{-1}\left(1-C\Bigl(\frac{\varphi}{2}\Bigr)^m\right),
\]
故全部项都被阈值截断，从而
\[
C_m^{\mathrm{bin}}(B_m)=2^{B_m}|X_m|=2^{B_m}F_{m+2}.
\]
于是
\[
\mathrm{Succ}_m^{\mathrm{bin}}(B_m)
=
\frac{2^{B_m}}{\lambda^m}\cdot \frac{F_{m+2}}{\varphi^m}
\longrightarrow
\zeta\cdot \frac{\varphi^2}{\sqrt5}.
\]

若 \(\varphi^{-1}<\zeta<1\)，则对充分大 \(m\) 有
\[
\lambda^m\varphi^{-1}\left(1+C\Bigl(\frac{\varphi}{2}\Bigr)^m\right)
<
2^{B_m}
<
\lambda^m\left(1-C\Bigl(\frac{\varphi}{2}\Bigr)^m\right).
\]
因此末位 \(0\) 扇区仍被截断，而末位 \(1\) 扇区已全部通过，故
\[
C_m^{\mathrm{bin}}(B_m)
=
F_{m+1}2^{B_m}+\sum_{w\in A_1^{(m)}}d_m^{\mathrm{bin}}(w).
\]
再由上一致展开，
\[
\sum_{w\in A_1^{(m)}}d_m^{\mathrm{bin}}(w)
=
F_m\lambda^m\varphi^{-1}\left(1+O\!\Bigl(\Bigl(\frac{\varphi}{2}\Bigr)^m\Bigr)\right).
\]
两端除以 \(2^m\) 得
\[
\mathrm{Succ}_m^{\mathrm{bin}}(B_m)
=
\frac{2^{B_m}}{\lambda^m}\cdot \frac{F_{m+1}}{\varphi^m}
+
\frac{F_m}{\varphi^{m+1}}
+
o(1)
\longrightarrow
\zeta\frac{\varphi}{\sqrt5}+\frac{1}{\varphi\sqrt5}.
\]

若 \(\zeta>1\)，则对充分大 \(m\) 有
\[
2^{B_m}>\lambda^m\left(1+C\Bigl(\frac{\varphi}{2}\Bigr)^m\right),
\]
故所有纤维都完全通过，从而
\[
C_m^{\mathrm{bin}}(B_m)=\sum_{w\in X_m}d_m^{\mathrm{bin}}(w)=2^m,
\]
即
\[
\mathrm{Succ}_m^{\mathrm{bin}}(B_m)\to 1.
\]
边界点 \(\zeta=\varphi^{-1},1\) 由上述分段公式的连续性直接补上。证毕。
\end{proof}

\begin{corollary}[critical-scale 预言机反演成功曲线的双折点与最小资源反演闭式]\label{cor:xi-foldbin-critical-scale-success-curve-inversion}
\leanverified{paper\_xi\_foldbin\_critical\_scale\_success\_curve\_inversion}
沿用定理 \ref{thm:xi-foldbin-dyadic-capacity-critical-window-limit} 的记号。对 \(s\ge 0\) 定义连续临界尺度
\[
T_m(s):=s\Bigl(\frac{2}{\varphi}\Bigr)^m,
\]
以及连续化成功曲线
\[
\mathcal S_m(s):=
2^{-m}\,\mathcal C_m^{\mathrm{bin,cont}}\!\bigl(T_m(s)\bigr).
\]
则 \(\mathcal S_m\) 在每个紧集 \(K\Subset(0,\infty)\) 上一致收敛到分段线性极限
\[
\boxed{\;
\mathcal S_\infty(s)=
\begin{cases}
\dfrac{\varphi^2}{\sqrt5}\,s, & 0\le s\le \varphi^{-1},\\[6pt]
\dfrac{\varphi}{\sqrt5}\,s+\dfrac{1}{\varphi\sqrt5}, & \varphi^{-1}\le s\le 1,\\[8pt]
1, & s\ge 1.
\end{cases}}
\]
从而极限成功曲线恰在
\[
s=\varphi^{-1},
\qquad
s=1
\]
处出现且仅出现两处刚性折点。

若进一步定义其最小资源广义逆
\[
\mathcal S_\infty^{-1}(\alpha):=\inf\{s\ge 0:\ \mathcal S_\infty(s)\ge \alpha\}
\qquad (0\le \alpha\le 1),
\]
则有显式公式
\[
\boxed{\;
\mathcal S_\infty^{-1}(\alpha)=
\begin{cases}
\dfrac{\sqrt5}{\varphi^2}\,\alpha, & 0\le \alpha\le \dfrac{\varphi}{\sqrt5},\\[8pt]
\dfrac{\sqrt5\,\alpha-\varphi^{-1}}{\varphi}, & \dfrac{\varphi}{\sqrt5}\le \alpha\le 1.
\end{cases}}
\]
\end{corollary}
\begin{proof}
由定理 \ref{thm:xi-foldbin-underresolution-two-atomic-curvature-measure}，
\[
\mathcal S_m(s)=\widehat C_m(s)
\]
在每个紧集 \(K\Subset(0,\infty)\) 上一致收敛到
\[
\widehat C_\infty(s)=\frac{\varphi}{\sqrt5}\min\{1,s\}+\frac{1}{\sqrt5}\min\{\varphi^{-1},s\}.
\]
再注意 \(\mathcal S_m(0)=0=\mathcal S_\infty(0)\)。将右端按 \(s\in[0,\varphi^{-1}]\)、\(s\in[\varphi^{-1},1]\) 与 \(s\ge 1\) 分段展开，即得所示三段式。因此折点恰位于两个最小值函数发生切换的
\[
s=\varphi^{-1},\qquad s=1.
\]

又因 \(\mathcal S_\infty\) 在 \([0,1]\) 上分段线性且严格递增，而在 \([1,\infty)\) 上恒等于 \(1\)，故其最小资源广义逆由两段线性方程直接求解而得。对第一段，
\[
\alpha=\frac{\varphi^2}{\sqrt5}\,s
\]
给出
\[
s=\frac{\sqrt5}{\varphi^2}\,\alpha.
\]
对第二段，
\[
\alpha=\frac{\varphi}{\sqrt5}\,s+\frac{1}{\varphi\sqrt5}
\]
给出
\[
s=\frac{\sqrt5\,\alpha-\varphi^{-1}}{\varphi}.
\]
证毕。
\end{proof}

\begin{corollary}[critical-scale 资源边际收益在首折点处恰按 \texorpdfstring{$\varphi$}{phi} 倍坍缩]\label{cor:xi-foldbin-critical-scale-marginal-gain-phi-drop}
\leanverified{paper\_xi\_foldbin\_critical\_scale\_marginal\_gain\_phi\_drop}
沿用推论 \ref{cor:xi-foldbin-critical-scale-success-curve-inversion} 的记号。则极限成功曲线的边际收益满足
\[
\mathcal S_\infty'(s)=
\begin{cases}
\dfrac{\varphi^2}{\sqrt5}, & 0<s<\varphi^{-1},\\[6pt]
\dfrac{\varphi}{\sqrt5}, & \varphi^{-1}<s<1,\\[6pt]
0, & s>1.
\end{cases}
\]
特别地，在首折点 \(s=\varphi^{-1}\) 处，左右导数之比满足
\[
\boxed{\;
\frac{\mathcal S_\infty'(\varphi^{-1}-0)}{\mathcal S_\infty'(\varphi^{-1}+0)}=\varphi.\;}
\]
亦即，一旦末位 \(1\) 层率先饱和，继续增加临界尺度资源所获得的边际成功收益恰好下降一个黄金比例因子。
\end{corollary}
\begin{proof}
对推论 \ref{cor:xi-foldbin-critical-scale-success-curve-inversion} 中的分段线性表达式逐段求导即可。证毕。
\end{proof}

\begin{theorem}[bin-fold 零温自由能修正的二相变分公式]\label{thm:xi-foldbin-two-phase-free-energy-variational}
定义
\[
\Psi(q):=\log\!\Bigl(\frac{\varphi+\varphi^{-q}}{\sqrt5}\Bigr),
\qquad
q\ge 0.
\]
则有精确变分公式
\[
\Psi(q)
=
\log\!\Bigl(\frac{\varphi}{\sqrt5}\Bigr)
+
\sup_{\theta\in[0,1]}
\Bigl\{
h(\theta)-\theta(q+1)\log\varphi
\Bigr\},
\]
其中
\[
h(\theta):=-\theta\log\theta-(1-\theta)\log(1-\theta)
\]
为二元 Shannon 熵。该上确界由唯一点
\[
\theta(q)=\frac{1}{1+\varphi^{q+1}}
\]
取得，并满足
\[
\Psi'(q)=-\theta(q)\log\varphi.
\]
若进一步沿用定理 \ref{thm:xi-fold-escort-renyi-shannon-constants} 的 escort 记号 \(\mu_{m,q}\)，并定义
\[
\kappa_{m,q}^{\mathrm{esc}}
:=
\mathbb E_{\mu_{m,q}}\!\bigl[\log d_m^{\mathrm{bin}}(W_{m,q})\bigr],
\qquad
W_{m,q}\sim \mu_{m,q},
\]
则对每个固定 \(q\ge 0\) 都有
\[
\kappa_{m,q}^{\mathrm{esc}}
=
m\log\!\Bigl(\frac{2}{\varphi}\Bigr)
+
\Psi'(q)
+
O\!\Bigl(\Bigl(\frac{\varphi}{2}\Bigr)^m\Bigr).
\]
\end{theorem}
\begin{proof}
设
\[
F_q(\theta):=h(\theta)-\theta(q+1)\log\varphi.
\]
则
\[
F_q'(\theta)=\log\frac{1-\theta}{\theta}-(q+1)\log\varphi,
\qquad
F_q''(\theta)=-\frac{1}{\theta(1-\theta)}<0.
\]
因此 \(F_q\) 在 \((0,1)\) 上严格凹，且唯一驻点满足
\[
\frac{1-\theta}{\theta}=\varphi^{q+1},
\]
亦即
\[
\theta=\theta(q)=\frac{1}{1+\varphi^{q+1}}.
\]
这就给出了唯一极大点。再用标准二元自由能恒等式
\[
\sup_{\theta\in[0,1]}\{h(\theta)-\theta L\}=\log(1+e^{-L}),
\qquad
L\in\RR,
\]
并在此处代入 \(L=(q+1)\log\varphi\)，得到
\[
\sup_{\theta\in[0,1]}F_q(\theta)
=
\log\bigl(1+\varphi^{-(q+1)}\bigr).
\]
因此
\[
\log\!\Bigl(\frac{\varphi}{\sqrt5}\Bigr)+\sup_{\theta\in[0,1]}F_q(\theta)
=
\log\!\Bigl(\frac{\varphi+\varphi^{-q}}{\sqrt5}\Bigr)
=
\Psi(q).
\]

直接对 \(\Psi\) 求导可得
\[
\Psi'(q)
=
-\frac{\varphi^{-q}\log\varphi}{\varphi+\varphi^{-q}}
=
-\frac{\log\varphi}{1+\varphi^{q+1}}
=
-\theta(q)\log\varphi.
\]

最后，由定理 \ref{thm:fold-bin-two-state-asymptotic} 的一致对数展开，
\[
\log d_m^{\mathrm{bin}}(w)
=
m\log\!\Bigl(\frac{2}{\varphi}\Bigr)
-w_m\log\varphi
+
O\!\Bigl(\Bigl(\frac{\varphi}{2}\Bigr)^m\Bigr)
\]
在 \(w\in X_m\) 上一致成立。另一方面，命题 \ref{prop:fold-bin-escort-lastbit} 给出
\[
\mu_{m,q}(w_m=1)=\theta(q)+O\!\Bigl(\Bigl(\frac{\varphi}{2}\Bigr)^m\Bigr).
\]
对上式在 \(\mu_{m,q}\) 下取期望，得到
\[
\kappa_{m,q}^{\mathrm{esc}}
=
m\log\!\Bigl(\frac{2}{\varphi}\Bigr)
-\theta(q)\log\varphi
+
O\!\Bigl(\Bigl(\frac{\varphi}{2}\Bigr)^m\Bigr).
\]
再代入 \(\Psi'(q)=-\theta(q)\log\varphi\) 即得所述展开。证毕。
\end{proof}

\begin{theorem}[bin-fold 隐藏位预算、可见熵与均匀基线 KL 缺口共享同一常数核]\label{thm:xi-foldbin-hidden-budget-visible-entropy-kl-one-constant}
沿用定理 \ref{thm:xi-foldbin-absolute-renyi-entropy-constant-closed} 与定理
\ref{thm:xi-foldbin-two-phase-free-energy-variational} 的记号。令
\[
\kappa_m:=\EE_{\mu_m}\!\bigl[\log d_m^{\mathrm{bin}}(W_m)\bigr],
\qquad
W_m\sim\mu_m,
\]
并记稳定层均匀分布为
\[
u_m(w):=\frac{1}{|X_m|}
\qquad (w\in X_m).
\]
则有三条共享同一隐藏位常数核的精确展开：
\[
\boxed{\;
\kappa_m
=
m\log\!\Bigl(\frac{2}{\varphi}\Bigr)
-\frac{\log\varphi}{\varphi\sqrt5}
+O\!\Bigl(\Bigl(\frac{\varphi}{2}\Bigr)^m\Bigr),\;}
\]
\[
\boxed{\;
H(\mu_m)
=
m\log\varphi
+\frac{\log\varphi}{\varphi\sqrt5}
+O\!\Bigl(\Bigl(\frac{\varphi}{2}\Bigr)^m\Bigr),\;}
\]
\[
\boxed{\;
\log|X_m|-H(\mu_m)
=
2\log\varphi-\frac12\log5-\frac{\log\varphi}{\varphi\sqrt5}
+O\!\Bigl(\Bigl(\frac{\varphi}{2}\Bigr)^m\Bigr).\;}
\]
因此
\[
\boxed{\;
D_{\mathrm{KL}}(\mu_m\Vert u_m)
=
2\log\varphi-\frac12\log5-\frac{\log\varphi}{\varphi\sqrt5}
+O\!\Bigl(\Bigl(\frac{\varphi}{2}\Bigr)^m\Bigr),\;}
\]
从而极限
\[
\boxed{\;
\lim_{m\to\infty}D_{\mathrm{KL}}(\mu_m\Vert u_m)
=
2\log\varphi-\frac12\log5-\frac{\log\varphi}{\varphi\sqrt5}>0.\;}
\]
存在且为严格正数。换言之，隐藏位预算、可见 Shannon 熵与相对于稳定层均匀基线的 KL 缺口并非三条彼此独立的数值链，而共享同一个刚性常数核
\[
\frac{\log\varphi}{\varphi\sqrt5}.
\]
\end{theorem}
\begin{proof}
在定理 \ref{thm:xi-foldbin-two-phase-free-energy-variational} 中取 \(q=1\)，并注意
\[
\kappa_m=\kappa_{m,1}^{\mathrm{esc}},
\qquad
\theta(1)=\frac{1}{1+\varphi^2}=\frac{1}{\varphi\sqrt5},
\]
便得到
\[
\kappa_m
=
m\log\!\Bigl(\frac{2}{\varphi}\Bigr)
-\frac{\log\varphi}{\varphi\sqrt5}
+O\!\Bigl(\Bigl(\frac{\varphi}{2}\Bigr)^m\Bigr).
\]
这证明了第一式。

第二式由 Shannon 熵恒等式
\[
H(\mu_m)=m\log2-\kappa_m
\]
直接推出。

再由稳定词计数
\[
|X_m|=F_{m+2}
=
\frac{\varphi^{m+2}}{\sqrt5}+O(\varphi^{-m}),
\]
可得
\[
\log|X_m|
=
(m+2)\log\varphi-\frac12\log5+O(\varphi^{-2m}).
\]
将其减去第二式，便得到第三式；而
\[
D_{\mathrm{KL}}(\mu_m\Vert u_m)=\log|X_m|-H(\mu_m)
\]
是均匀基线下的标准熵差公式，因此第四式与其极限立即成立。严格正性则可由
\[
\frac{1}{\varphi\sqrt5}\neq \varphi^{-2}
\]
以及推论 \ref{cor:xi-fold-uniform-baseline-classical-fdiv-bernoulli-collapse} 在
\[
f(u)=u\log u
\]
处的 Bernoulli 闭式看出。证毕。
\end{proof}

\begin{theorem}[bin-fold 容错阈值预算的 \texorpdfstring{$\log_2\varphi$}{log2 phi} 常数跃迁与指数斜率锁定]\label{thm:xi-foldbin-threshold-budget-logphi-jump}
沿用推论 \ref{cor:fold-bin-quantile-threshold-constant} 的记号，定义规范阈值预算
\[
P_m^{\mathrm{thr}}(\varepsilon):=2B^{\mathrm{bin}}_m(\varepsilon),
\qquad
\varepsilon\in(0,1),
\]
并记临界误差阈值
\[
\varepsilon_{\mathrm c}:=\frac{\varphi}{\sqrt5}.
\]
则有：
\begin{enumerate}
  \item 对任意固定
  \[
  0<\varepsilon_-<\varepsilon_{\mathrm c}<\varepsilon_+<1,
  \]
  都有
  \[
  \frac{P_m^{\mathrm{thr}}(\varepsilon_+)}{P_m^{\mathrm{thr}}(\varepsilon_-)}
  \longrightarrow
  \varphi^{-1},
  \]
  从而二进制对数预算差满足
  \[
  \log_2 P_m^{\mathrm{thr}}(\varepsilon_-)
  -
  \log_2 P_m^{\mathrm{thr}}(\varepsilon_+)
  \longrightarrow
  \log_2\varphi.
  \]
  \item 对任意固定
  \[
  \varepsilon\in(0,\varepsilon_{\mathrm c})\cup(\varepsilon_{\mathrm c},1),
  \]
  都有
  \[
  \lim_{m\to\infty}\frac1m\log P_m^{\mathrm{thr}}(\varepsilon)
  =
  \log\frac{2}{\varphi}
  =
  \lim_{m\to\infty}\frac1m\log d_{\max}^{\mathrm{bin}}(m).
  \]
\end{enumerate}
因此，阈值 \(\varepsilon_{\mathrm c}\) 只引入 \(\log_2\varphi\) 量级的常数预算跃迁，而不改变由最大纤维指数锁定的指数主斜率。
\end{theorem}
\begin{proof}
由推论 \ref{cor:fold-bin-quantile-threshold-constant}，
\[
\Bigl(\frac{\varphi}{2}\Bigr)^m B^{\mathrm{bin}}_m(\varepsilon)\longrightarrow
\begin{cases}
1,& 0<\varepsilon<\varepsilon_{\mathrm c},\\
\varphi^{-1},& \varepsilon_{\mathrm c}<\varepsilon<1.
\end{cases}
\]
乘以 \(2\) 即得
\[
P_m^{\mathrm{thr}}(\varepsilon_-)
=
2\Bigl(\frac{2}{\varphi}\Bigr)^m(1+o(1)),
\qquad
P_m^{\mathrm{thr}}(\varepsilon_+)
=
2\varphi^{-1}\Bigl(\frac{2}{\varphi}\Bigr)^m(1+o(1)),
\]
从而
\[
\frac{P_m^{\mathrm{thr}}(\varepsilon_+)}{P_m^{\mathrm{thr}}(\varepsilon_-)}
\longrightarrow
\varphi^{-1}.
\]
再取二进制对数即可得到
\[
\log_2 P_m^{\mathrm{thr}}(\varepsilon_-)
-
\log_2 P_m^{\mathrm{thr}}(\varepsilon_+)
\longrightarrow
\log_2\varphi.
\]
这就证明了第一条。

同样由上式，对任意固定
\[
\varepsilon\in(0,\varepsilon_{\mathrm c})\cup(\varepsilon_{\mathrm c},1)
\]
都有
\[
\lim_{m\to\infty}\frac1m\log P_m^{\mathrm{thr}}(\varepsilon)
=
\log\frac{2}{\varphi}
\]
另一方面，定理 \ref{thm:fold-bin-max-fiber-exponent} 给出
\[
\lim_{m\to\infty}\frac1m\log d_{\max}^{\mathrm{bin}}(m)
=
\log\frac{2}{\varphi}.
\]
两式合并即得所示指数斜率锁定。证毕。
\end{proof}

\begin{theorem}[bin-fold 最大纤维组合上界的指数精确性与最终严格失饱和]\label{thm:xi-foldbin-maxfiber-upper-bound-exponential-sharp-nonsaturation}
沿用推论 \ref{cor:fold-bin-saturation} 的记号，定义缺口
\[
\delta_m:=F_{K(m)-m+2}-d_{\max}^{\mathrm{bin}}(m).
\]
则有：
\begin{enumerate}
  \item 存在 \(m_0\)，使得对全部 \(m\ge m_0\) 都有
  \[
  \delta_m>0.
  \]
  \item 并且
  \[
  \lim_{m\to\infty}\frac1m\log F_{K(m)-m+2}
  =
  \lim_{m\to\infty}\frac1m\log d_{\max}^{\mathrm{bin}}(m)
  =
  \log\frac{2}{\varphi},
  \]
  从而
  \[
  \lim_{m\to\infty}\frac1m\log\frac{F_{K(m)-m+2}}{d_{\max}^{\mathrm{bin}}(m)}=0.
  \]
\end{enumerate}
换言之，组合上界在指数尺度上是精确的，但在算术精确尺度上最终严格失饱和。
\end{theorem}
\begin{proof}
由注 \ref{rem:fold-bin-finite}，上界饱和
\[
d_{\max}^{\mathrm{bin}}(m)=F_{K(m)-m+2}
\]
只能发生有限次。故存在 \(m_0\)，使得对全部 \(m\ge m_0\) 都有
\[
\delta_m=F_{K(m)-m+2}-d_{\max}^{\mathrm{bin}}(m)>0.
\]
这就证明了第一条。

对第二条，由定义 \eqref{eq:K-of-m}
\[
F_{K(m)+1}\le 2^m-1<F_{K(m)+2}.
\]
结合 Binet 渐近
\[
\log F_n=n\log\varphi+O(1)
\]
可得
\[
K(m)=m\log_{\varphi}2+O(1).
\]
于是
\[
\log F_{K(m)-m+2}
=
\bigl(K(m)-m+2\bigr)\log\varphi+O(1)
=
m(\log2-\log\varphi)+O(1),
\]
从而
\[
\lim_{m\to\infty}\frac1m\log F_{K(m)-m+2}
=
\log\frac{2}{\varphi}.
\]
另一方面，由定理 \ref{thm:fold-bin-max-fiber-exponent}，
\[
\lim_{m\to\infty}\frac1m\log d_{\max}^{\mathrm{bin}}(m)
=
\log\frac{2}{\varphi}.
\]
两式相减即得
\[
\lim_{m\to\infty}\frac1m\log\frac{F_{K(m)-m+2}}{d_{\max}^{\mathrm{bin}}(m)}=0.
\]
证毕。
\end{proof}

\endinput

\paragraph{escort Rényi 熵常数的 \texorpdfstring{$\theta(q)$}{theta(q)}-闭式化}
在 Rényi 熵率塌缩与 escort 常数谱已经确立之后，温度参数 \(q\) 的全部 \(O(1)\) 熵信息还可继续压缩为末位 Bernoulli 质量
\(
\theta(q)=1/(1+\varphi^{q+1})
\)。

\begin{theorem}[escort Rényi 熵的速率塌缩与常数项全解]\label{thm:xi-fold-escort-renyi-constant-theta-closed}
\leanverified{paper\_xi\_fold\_escort\_renyi\_constant\_theta\_closed}
沿用定理 \ref{thm:xi-fold-escort-renyi-shannon-constants} 的记号。记
\[
\theta(q):=\frac{1}{1+\varphi^{q+1}}.
\]
则对任意固定 \(q\ge 0\) 与任意 \(r>0\)、\(r\neq 1\)，有
\[
H_r(\mu_{m,q})
=
m\log\varphi
-\frac12\log 5
+
\frac{1}{1-r}
\log\!\Bigl(
\varphi^{1-r}(1-\theta(q))^r+\theta(q)^r
\Bigr)
+
O\!\Bigl(\Bigl(\frac{\varphi}{2}\Bigr)^m\Bigr).
\]
等价地，
\[
H_r(\mu_{m,q})
=
m\log\varphi
-\frac12\log 5
+
\frac{1}{1-r}
\Bigl[
\log\!\bigl(1+\varphi^{1+qr}\bigr)
-r\log\!\bigl(1+\varphi^{q+1}\bigr)
\Bigr]
+
O\!\Bigl(\Bigl(\frac{\varphi}{2}\Bigr)^m\Bigr).
\]
相应的 Shannon 熵满足
\[
H(\mu_{m,q})
=
m\log\varphi
-\frac12\log 5
+
h_2(\theta(q))
+
(1-\theta(q))\log\varphi
+
O\!\Bigl(\Bigl(\frac{\varphi}{2}\Bigr)^m\Bigr),
\]
其中
\[
h_2(p):=-p\log p-(1-p)\log(1-p).
\]
因此 escort 温度 \(q\) 的全部 \(O(1)\) 熵信息都被压缩为单变量 \(\theta(q)\)。
\end{theorem}
\begin{proof}
由定理 \ref{thm:xi-fold-escort-renyi-shannon-constants}，
\[
H_r(\mu_{m,q})
=
m\log\varphi
+
\frac{1}{1-r}\left[
\log\!\Bigl(\frac{\varphi+\varphi^{-qr}}{\sqrt5}\Bigr)
-r\log\!\Bigl(\frac{\varphi+\varphi^{-q}}{\sqrt5}\Bigr)
\right]
+
O\!\Bigl(\Bigl(\frac{\varphi}{2}\Bigr)^m\Bigr).
\]
再用
\[
\varphi+\varphi^{-a}=\varphi^{-a}(1+\varphi^{a+1}),
\]
便得到第二个 Rényi 公式。

另一方面，由
\[
1-\theta(q)=\frac{\varphi^{q+1}}{1+\varphi^{q+1}}
\]
可直接化简出
\[
\varphi^{1-r}(1-\theta(q))^r+\theta(q)^r
=
\frac{1+\varphi^{1+qr}}{(1+\varphi^{q+1})^r},
\]
故第一个 Rényi 公式与第二个完全等价。

对 Shannon 熵，定理 \ref{thm:xi-fold-escort-renyi-shannon-constants} 同样给出
\[
H(\mu_{m,q})
=
m\log\varphi
+
\log\!\Bigl(\frac{\varphi+\varphi^{-q}}{\sqrt5}\Bigr)
+
\frac{q\log\varphi}{1+\varphi^{q+1}}
+
O\!\Bigl(\Bigl(\frac{\varphi}{2}\Bigr)^m\Bigr).
\]
再注意到
\[
h_2(\theta(q))
=
\log\!\bigl(1+\varphi^{q+1}\bigr)
-(q+1)(1-\theta(q))\log\varphi,
\]
于是
\[
-\frac12\log 5
+
h_2(\theta(q))
+
(1-\theta(q))\log\varphi
=
\log\!\Bigl(\frac{\varphi+\varphi^{-q}}{\sqrt5}\Bigr)
+
\frac{q\log\varphi}{1+\varphi^{q+1}}.
\]
这就给出了 Shannon 常数项的所述写法。证毕。
\end{proof}

\endinput

\paragraph{冻结锥平台、极大纤维骨架与纤维指数自平均}
在固定冻结相的 Rényi 熵率塌缩、基态质量指数集中以及精确反演位率阈值已经建立之后，还可把冻结锥上的平台律、极大纤维骨架投影、辨路径预算相变与纤维指数自平均进一步压缩为下列四条结果。

\begin{theorem}[冻结锥上的 escort Rényi 平台]\label{thm:xi-fixed-freezing-cone-escort-renyi-plateau}
\leanverified{paper\_xi\_fixed\_freezing\_cone\_escort\_renyi\_plateau}
沿用定理 \ref{thm:xi-fixed-freezing-escort-renyi-spectrum-collapse} 的记号。对 \(a>0\) 定义 escort 分布
\[
\pi_{m,a}(x):=\frac{d_m(x)^a}{S_a(m)},
\qquad
S_a(m):=\sum_{x\in X_m}d_m(x)^a,
\]
并对 \(q>0\)、\(q\neq 1\) 定义 escort Rényi 熵
\[
H_q^{\mathrm{esc}}(m;a):=
\frac{1}{1-q}\log\sum_{x\in X_m}\pi_{m,a}(x)^q.
\]
若
\[
a>a_{\mathrm c},
\qquad
aq>a_{\mathrm c},
\]
则极限
\[
h_q^{\mathrm{esc}}(a):=
\lim_{m\to\infty}\frac1m H_q^{\mathrm{esc}}(m;a)
\]
存在，并且满足
\[
h_q^{\mathrm{esc}}(a)=g_*.
\]
\end{theorem}
\begin{proof}
由定义
\[
\sum_{x\in X_m}\pi_{m,a}(x)^q
=
\frac{S_{aq}(m)}{S_a(m)^q},
\]
故
\[
\frac1m H_q^{\mathrm{esc}}(m;a)
=
\frac{1}{1-q}
\left(
\frac1m\log S_{aq}(m)-q\frac1m\log S_a(m)
\right).
\]
当 \(a>a_{\mathrm c}\) 且 \(aq>a_{\mathrm c}\) 时，定理 \ref{thm:fold-pressure-freezing-threshold} 给出
\[
\lim_{m\to\infty}\frac1m\log S_t(m)=t\alpha_*+g_*
\qquad
(t>a_{\mathrm c}).
\]
将 \(t=a\) 与 \(t=aq\) 代回，便得到
\[
h_q^{\mathrm{esc}}(a)
=
\frac{(aq\alpha_*+g_*)-q(a\alpha_*+g_*)}{1-q}
=
g_*.
\]
证毕。
\end{proof}

\begin{theorem}[冻结相的极大纤维骨架投影]\label{thm:xi-fixed-freezing-maxfiber-skeleton-projection}
\leanverified{paper\_xi\_fixed\_freezing\_maxfiber\_skeleton\_projection}
记
\[
A_m^{\max}:=\{x\in X_m:\ d_m(x)=M_m\},
\qquad
K_m:=|A_m^{\max}|,
\]
并令 \(\nu_m\) 为 \(A_m^{\max}\) 上的均匀分布，即
\[
\nu_m(x):=
\begin{cases}
K_m^{-1}, & x\in A_m^{\max},\\
0, & x\notin A_m^{\max}.
\end{cases}
\]
若 \(a>a_{\mathrm c}\)，则对任意满足
\[
\|f_m\|_{\infty}\le 1
\qquad
\bigl(f_m:X_m\to\CC\bigr)
\]
的测试函数，都有
\[
\left|
\sum_{x\in X_m}\pi_{m,a}(x)f_m(x)
-
\sum_{x\in X_m}\nu_m(x)f_m(x)
\right|
\le
2\exp\!\bigl(-m\Delta(a)+o(m)\bigr),
\]
其中
\[
\Delta(a):=a(\alpha_*-\alpha_*^{(2)})-(\log\varphi-g_*)>0.
\]
\end{theorem}
\begin{proof}
由定理 \ref{thm:xi-fixed-freezing-escort-renyi-spectrum-collapse} 的第四条，
\[
\|\pi_{m,a}-\nu_m\|_{\mathrm{TV}}
=
1-\pi_{m,a}(A_m^{\max})
\le
\exp\!\bigl(-m\Delta(a)+o(m)\bigr).
\]
再由总变差距离的对偶表述，对任意 \(\|f_m\|_\infty\le 1\) 都有
\[
\left|
\sum_{x\in X_m}\pi_{m,a}(x)f_m(x)
-
\sum_{x\in X_m}\nu_m(x)f_m(x)
\right|
\le
2\|\pi_{m,a}-\nu_m\|_{\mathrm{TV}}.
\]
将前式代回即得结论。证毕。
\end{proof}

\begin{theorem}[冻结相中的辨路径位率相变]\label{thm:xi-fixed-freezing-path-budget-phase-transition}
设 \(a>a_{\mathrm c}\)。对整数预算 \(B\ge 0\)，定义 escort-可解码质量
\[
\Gamma_{m,a}(B):=
\pi_{m,a}\bigl(\{x\in X_m:\ d_m(x)\le 2^B\}\bigr).
\]
令
\[
B_m:=\lfloor \rho m\rfloor.
\]
并记
\[
\Delta(a):=a(\alpha_*-\alpha_*^{(2)})-(\log\varphi-g_*)>0.
\]
则存在严格相变：
\begin{enumerate}
  \item 若
  \[
  \rho<\frac{\alpha_*}{\log 2},
  \]
  则
  \[
  \Gamma_{m,a}(B_m)\le \exp\!\bigl(-m\Delta(a)+o(m)\bigr)\xrightarrow[m\to\infty]{}0.
  \]
  \item 若
  \[
  \rho>\frac{\alpha_*}{\log 2},
  \]
  则对充分大的 \(m\) 有
  \[
  \Gamma_{m,a}(B_m)=1.
  \]
\end{enumerate}
\end{theorem}
\begin{proof}
若 \(\rho<\alpha_*/\log 2\)，则由
\[
B_m=\rho m+O(1),
\qquad
M_m=\exp(\alpha_*m+o(m))
\]
可得对充分大的 \(m\)，
\[
2^{B_m}=\exp\!\bigl((\rho\log 2)m+o(m)\bigr)<M_m.
\]
因此每个 \(x\in A_m^{\max}\) 都满足 \(d_m(x)=M_m>2^{B_m}\)，从而
\[
\Gamma_{m,a}(B_m)\le \pi_{m,a}(X_m\setminus A_m^{\max}).
\]
再由定理 \ref{thm:xi-fixed-freezing-escort-renyi-spectrum-collapse} 的第四条，
\[
\pi_{m,a}(X_m\setminus A_m^{\max})
=
1-\pi_{m,a}(A_m^{\max})
\le
\exp\!\bigl(-m\Delta(a)+o(m)\bigr),
\]
即得第一条。

若 \(\rho>\alpha_*/\log 2\)，则对充分大的 \(m\) 有
\[
2^{B_m}>M_m.
\]
由于对一切 \(x\in X_m\) 都有 \(d_m(x)\le M_m\)，故
\[
\{x\in X_m:\ d_m(x)\le 2^{B_m}\}=X_m,
\]
于是
\[
\Gamma_{m,a}(B_m)=1.
\]
证毕。
\end{proof}

\begin{theorem}[冻结相中的纤维指数自平均]\label{thm:xi-fixed-freezing-fiber-exponent-selfaveraging}
设 \(a>a_{\mathrm c}\)，并定义
\[
Y_m(x):=\frac{1}{m}\log d_m(x),
\qquad
x\in X_m.
\]
对任意满足
\[
a+t>a_{\mathrm c}
\]
的实数 \(t\)，极限
\[
\Lambda_a(t):=
\lim_{m\to\infty}
\frac1m
\log \EE_{\pi_{m,a}}\!\left[e^{mtY_m}\right]
\]
存在，且
\[
\Lambda_a(t)=\alpha_* t.
\]
进一步，对任意有限 \(p>0\)，都有
\[
\lim_{m\to\infty}
\EE_{\pi_{m,a}}
\left[
\left|Y_m-\alpha_*\right|^p
\right]
=0.
\]
\end{theorem}
\begin{proof}
先计算指数母函数。由定义
\[
\EE_{\pi_{m,a}}\!\left[e^{mtY_m}\right]
=
\sum_{x\in X_m}\frac{d_m(x)^a}{S_a(m)}\,d_m(x)^t
=
\frac{S_{a+t}(m)}{S_a(m)}.
\]
因此
\[
\frac1m
\log \EE_{\pi_{m,a}}\!\left[e^{mtY_m}\right]
=
\frac1m\log S_{a+t}(m)-\frac1m\log S_a(m).
\]
由于 \(a>a_{\mathrm c}\) 且 \(a+t>a_{\mathrm c}\)，定理 \ref{thm:fold-pressure-freezing-threshold} 给出
\[
\lim_{m\to\infty}\frac1m\log S_s(m)=s\alpha_*+g_*
\qquad
(s>a_{\mathrm c}),
\]
故
\[
\Lambda_a(t)
=
(a+t)\alpha_*+g_*-(a\alpha_*+g_*)
=
\alpha_* t.
\]

再证 \(L^p\) 收敛。由 \(1\le d_m(x)\le 2^m\)，对全部 \(x\in X_m\) 都有
\[
0\le Y_m(x)\le \log 2.
\]
定义有界测试函数
\[
f_m(x):=\frac{|Y_m(x)-\alpha_*|^p}{(\log 2+|\alpha_*|)^p},
\]
则 \(\|f_m\|_\infty\le 1\)。另一方面，在 \(A_m^{\max}\) 上 \(Y_m\) 常值，故
\[
\sum_{x\in X_m}\nu_m(x)f_m(x)
=
\left|
\frac1m\log M_m-\alpha_*
\right|^p
(\log 2+|\alpha_*|)^{-p}.
\]
由定理 \ref{thm:xi-fixed-freezing-maxfiber-skeleton-projection}，
\[
\left|
\EE_{\pi_{m,a}}[f_m]-\sum_{x\in X_m}\nu_m(x)f_m(x)
\right|
\le
2\exp\!\bigl(-m\Delta(a)+o(m)\bigr).
\]
再乘回 \((\log 2+|\alpha_*|)^p\)，便得到
\[
\EE_{\pi_{m,a}}\!\left[|Y_m-\alpha_*|^p\right]
\le
\left|
\frac1m\log M_m-\alpha_*
\right|^p
+2(\log 2+|\alpha_*|)^p
\exp\!\bigl(-m\Delta(a)+o(m)\bigr).
\]
由 \(\frac1m\log M_m\to\alpha_*\)，右端趋于 \(0\)，故结论成立。证毕。
\end{proof}

\endinput

\paragraph{冻结尾部的凸延拓刚性、Legendre 硬壁与 Abel 有限部余项}
在定理 \ref{thm:xi-time-part57b-pressure-spectrum-discrete-convexity-monotone-excess} 已把整数压力谱锚定为离散凸序列、定理 \ref{thm:xi-freezing-affine-rigidity-two-point-recovery} 已把冻结相整数尾部压缩为两点即可恢复的仿射律之后，还可继续把连续凸延拓、Legendre 对偶边界与 Möbius--Euler 截断余项进一步收束为下列四条后果。

\begin{theorem}[冻结阈值之后任意凸压力延拓的仿射刚性]\label{thm:xi-time-part57f-convex-extension-affine-tail}
\leanverified{paper\_xi\_time\_part57f\_convex\_extension\_affine\_tail}
在严格基态间隙
\[
\alpha_*^{(2)}<\alpha_*
\]
下，记
\[
a_{\mathrm c}:=\frac{\log\varphi-g_*}{\alpha_*-\alpha_*^{(2)}},
\qquad
a_0:=\lfloor a_{\mathrm c}\rfloor+1.
\]
设
\[
\Lambda:[a_0,\infty)\to\RR
\]
为任意凸函数，并满足
\[
\Lambda(n)=P_n
\qquad
\bigl(n\in\ZZ,\ n\ge a_0\bigr).
\]
则必有
\[
\boxed{\;
\Lambda(t)=\alpha_* t+g_*
\qquad
(t\ge a_0).\;}
\]
特别地，任何保持凸性并在冻结区整数锚点与离散压力序列对齐的连续压力外推，在阈值之后都不再允许残留曲率；因此，定理 \ref{thm:xi-projective-pressure-canonical-convex-majorant} 所给出的连续压力接口一旦在这些整数锚点上实现精确对齐，其尾部也只能是同一条仿射直线。
\end{theorem}
\begin{proof}
由定理 \ref{thm:fold-pressure-freezing-threshold}，对每个整数 \(n\ge a_0\) 都有
\[
P_n=n\alpha_*+g_*,
\]
故
\[
G(t):=\Lambda(t)-(\alpha_* t+g_*)
\]
是定义在 \([a_0,\infty)\) 上的凸函数，并满足
\[
G(n)=0
\qquad
\bigl(n\in\ZZ,\ n\ge a_0\bigr).
\]

先看首个区间 \([a_0,a_0+1]\)。端点 \(t=a_0,a_0+1\) 时结论显然，故只需取
\[
t\in(a_0,a_0+1).
\]
由凸函数割线斜率单调性，对
\[
t<a_0+1<a_0+2
\]
应用
\[
\frac{G(a_0+1)-G(t)}{a_0+1-t}
\le
\frac{G(a_0+2)-G(a_0+1)}{1},
\]
便得
\[
\frac{-G(t)}{a_0+1-t}\le 0,
\]
从而 \(G(t)\ge 0\)。另一方面，由
\[
a_0\le t\le a_0+1
\]
与凸性，
\[
G(t)\le
\frac{a_0+1-t}{1}G(a_0)+\frac{t-a_0}{1}G(a_0+1)=0,
\]
故 \(G(t)=0\)。于是
\[
G\equiv 0
\qquad\text{于 }[a_0,a_0+1].
\]

再取任意整数 \(n\ge a_0+1\)。端点 \(t=n,n+1\) 时同样平凡，故只需考虑
\[
t\in(n,n+1).
\]
同样由割线斜率单调性，对
\[
n-1<n\le t
\]
有
\[
\frac{G(n)-G(n-1)}{1}
\le
\frac{G(t)-G(n)}{t-n},
\]
即
\[
0\le \frac{G(t)}{t-n},
\]
故 \(G(t)\ge 0\)。而由 \(G(n)=G(n+1)=0\) 与凸性，
\[
G(t)\le
\frac{n+1-t}{1}G(n)+\frac{t-n}{1}G(n+1)=0.
\]
因此 \(G(t)=0\)。这表明 \(G\equiv 0\) 于每个区间 \([n,n+1]\) 上成立。并合全部相邻整数区间，即得
\[
G(t)\equiv 0
\qquad
(t\ge a_0),
\]
亦即
\[
\Lambda(t)=\alpha_* t+g_*.
\]
证毕。
\end{proof}

\begin{theorem}[冻结尾部的 Legendre 硬壁与边界代价]\label{thm:xi-time-part57f-legendre-hard-wall-boundary-cost}
沿用定理 \ref{thm:xi-time-part57f-convex-extension-affine-tail} 的记号。在 \(X_m\) 上取均匀分布 \(u_m\)，并定义随机变量
\[
Y_m(x):=\frac{1}{m}\log d_m(x)
\qquad
(x\in X_m).
\]
则对每个整数 \(a\ge 1\)，有
\[
\frac1m\log \EE_{u_m}\!\left[e^{amY_m}\right]
=
\frac1m\log\frac{S_a(m)}{|X_m|}
\longrightarrow
P_a-\log\varphi.
\]
再设
\[
I(y):=
\sup_{t\ge a_0}
\Bigl\{
ty-\bigl(\Lambda(t)-\log\varphi\bigr)
\Bigr\}
\]
为任一此类凸延拓 \(\Lambda\) 所对应的 Legendre--Fenchel 对偶。则成立：
\begin{enumerate}
  \item 对任意 \(y>\alpha_*\)，有
  \[
  \boxed{\;I(y)=+\infty.\;}
  \]
  \item 在边界点 \(y=\alpha_*\) 处，恰有
  \[
  \boxed{\;I(\alpha_*)=\log\varphi-g_*.\;}
  \]
\end{enumerate}
换言之，\(\alpha_*\) 是冻结尾部在 Legendre 对偶中的硬上壁，而抵达该边界的指数代价恰由环境熵率与基态简并指数之差给出。
\end{theorem}
\begin{proof}
由定义
\[
\EE_{u_m}\!\left[e^{amY_m}\right]
=
\frac{1}{|X_m|}
\sum_{x\in X_m} d_m(x)^a
=
\frac{S_a(m)}{|X_m|}.
\]
又因 \(|X_m|=F_{m+2}\) 且
\[
\lim_{m\to\infty}\frac1m\log F_{m+2}=\log\varphi,
\]
故
\[
\frac1m\log \EE_{u_m}\!\left[e^{amY_m}\right]
=
\frac1m\log S_a(m)-\frac1m\log|X_m|
\longrightarrow
P_a-\log\varphi.
\]

另一方面，由定理 \ref{thm:xi-time-part57f-convex-extension-affine-tail}，
\[
\Lambda(t)=\alpha_* t+g_*
\qquad
(t\ge a_0).
\]
于是对任意 \(t\ge a_0\)，都有
\[
ty-\bigl(\Lambda(t)-\log\varphi\bigr)
=
t(y-\alpha_*)+\log\varphi-g_*.
\]
若 \(y>\alpha_*\)，右端随 \(t\to+\infty\) 线性发散到 \(+\infty\)，从而
\[
I(y)=+\infty.
\]
若 \(y=\alpha_*\)，则上式对所有 \(t\ge a_0\) 都恒等于
\[
\log\varphi-g_*,
\]
故其上确界亦为该常数，即
\[
I(\alpha_*)=\log\varphi-g_*.
\]
证毕。
\end{proof}

\begin{theorem}[Abel 有限部常数的内禀 \texorpdfstring{$\eta$}{eta}-截断误差界]\label{thm:xi-time-part57f-abel-finite-part-eta-tail}
沿用命题 \ref{prop:real-input-40-logM-infty} 的记号。记
\[
H(x):=(1-x)(1-6x+9x^2-x^3-x^4)
=
1-7x+15x^2-10x^3+x^5.
\]
并定义
\[
\eta:=\min_{0\le x\le b^2}H(x)>0.
\]
对每个整数 \(N\ge 2\)，记
\[
\log\mathfrak M_\infty^{(N)}
:=
\log C_\infty
-\sum_{m=2}^{N}\frac{\mu(m)}{m}\log H(b^m).
\]
则截断误差满足显式上界
\[
\boxed{\;
\bigl|\log\mathfrak M_\infty-\log\mathfrak M_\infty^{(N)}\bigr|
\le
\frac{7}{\eta(1-b)}\cdot \frac{b^{N+1}}{N+1}.\;}
\]
因此，Abel 有限部常数的 Möbius--Euler 展开不只几何收敛，而且其余项可由区间 \([0,b^2]\) 上的内禀最小值 \(\eta\) 直接审计。
\end{theorem}
\begin{proof}
由于 \(b\) 是多项式
\[
1-6x+9x^2-x^3-x^4
\]
的最大正根，故对 \(0\le x\le b^2<b\) 有
\[
1-6x+9x^2-x^3-x^4>0
\qquad\text{且}\qquad
1-x>0.
\]
因此
\[
H(x)>0
\qquad
(0\le x\le b^2),
\]
从而 \(\eta>0\)。

又由命题 \ref{prop:real-input-40-logM-infty} 的同一根定位，可知
\[
\frac14<b<\frac13,
\]
于是
\[
b^m\le b^2<\frac19
\qquad
(m\ge 2).
\]
对 \(x\in[0,b^2]\) 直接计算得
\[
1-H(x)=7x-15x^2+10x^3-x^5\le 7x.
\]
另一方面，由 \(x\le 1/9\) 可知
\[
1-H(x)
=
x\bigl(7-15x+10x^2-x^4\bigr)
\ge
x\Bigl(7-\frac{15}{9}\Bigr)\ge 0,
\]
故
\[
0<H(x)\le 1
\qquad
(0\le x\le b^2).
\]
因此对每个 \(m\ge 2\)，都可应用基本不等式
\[
-\log t\le \frac{1-t}{t}
\qquad
(0<t\le 1),
\]
得到
\[
\bigl|\log H(b^m)\bigr|
=
-\log H(b^m)
\le
\frac{1-H(b^m)}{H(b^m)}
\le
\frac{7b^m}{\eta}.
\]
于是
\[
\bigl|\log\mathfrak M_\infty-\log\mathfrak M_\infty^{(N)}\bigr|
\le
\sum_{m>N}\frac{|\mu(m)|}{m}\,\bigl|\log H(b^m)\bigr|
\le
\frac{7}{\eta}\sum_{m>N}\frac{b^m}{m}.
\]
而对 \(m>N\) 有 \(1/m\le 1/(N+1)\)，故
\[
\sum_{m>N}\frac{b^m}{m}
\le
\frac{1}{N+1}\sum_{m>N}b^m
=
\frac{1}{N+1}\cdot \frac{b^{N+1}}{1-b}.
\]
代回即得所示盒装上界。证毕。
\end{proof}

\begin{proposition}[Möbius--Euler 余项的首项由下一平方自由层主导]\label{prop:xi-time-part57f-mobius-euler-next-squarefree-layer}
沿用定理 \ref{thm:xi-time-part57f-abel-finite-part-eta-tail} 的记号。定义截断余项
\[
R_N:=
\log\mathfrak M_\infty-\log\mathfrak M_\infty^{(N)}
=
-\sum_{m>N}\frac{\mu(m)}{m}\log H(b^m).
\]
则当 \(N\to\infty\) 时，有统一渐近展开
\[
\boxed{\;
R_N
=
7\sum_{m>N}\frac{\mu(m)}{m}b^m
+
O\!\left(\frac{b^{2N+2}}{N+1}\right).\;}
\]
特别地，若 \(N+1\) 为平方自由数，则
\[
\boxed{\;
R_N
=
\frac{7\,\mu(N+1)}{N+1}b^{N+1}
+
O\!\left(\frac{b^{N+2}}{N+2}\right).\;}
\]
因此，余项的首个非零主项由下一层平方自由指数唯一决定，其符号直接由 \(\mu(N+1)\) 控制。
\end{proposition}
\begin{proof}
由
\[
H(x)=1-7x+O(x^2)
\qquad
(x\to 0),
\]
可得
\[
\log H(x)=-7x+O(x^2)
\qquad
(x\to 0).
\]
把 \(x=b^m\) 代入并带回余项展开，得到
\[
\begin{aligned}
R_N
&=
-\sum_{m>N}\frac{\mu(m)}{m}\Bigl(-7b^m+O(b^{2m})\Bigr)\\
&=
7\sum_{m>N}\frac{\mu(m)}{m}b^m
+
O\!\left(\sum_{m>N}\frac{b^{2m}}{m}\right).
\end{aligned}
\]
又因
\[
\sum_{m>N}\frac{b^{2m}}{m}
\le
\frac{1}{N+1}\sum_{m>N}b^{2m}
=
O\!\left(\frac{b^{2N+2}}{N+1}\right),
\]
便得第一条渐近式。

若 \(N+1\) 为平方自由数，则 \(\mu(N+1)\neq 0\)，从而
\[
\sum_{m>N}\frac{\mu(m)}{m}b^m
=
\frac{\mu(N+1)}{N+1}b^{N+1}
+
O\!\left(\sum_{m\ge N+2}\frac{b^m}{m}\right).
\]
而
\[
\sum_{m\ge N+2}\frac{b^m}{m}
\le
\frac{1}{N+2}\sum_{m\ge N+2}b^m
=
O\!\left(\frac{b^{N+2}}{N+2}\right).
\]
将其代回第一条式子，即得
\[
R_N
=
\frac{7\,\mu(N+1)}{N+1}b^{N+1}
+
O\!\left(\frac{b^{N+2}}{N+2}\right).
\]
证毕。
\end{proof}

\noindent 结合定理 \ref{thm:xi-time-part57b-pressure-spectrum-discrete-convexity-monotone-excess} 的离散凸性、定理 \ref{thm:xi-time-part57f-convex-extension-affine-tail} 的凸延拓刚性、定理 \ref{thm:xi-time-part57f-legendre-hard-wall-boundary-cost} 的 Legendre 边界代价，以及定理 \ref{thm:xi-time-part57f-abel-finite-part-eta-tail} 与命题 \ref{prop:xi-time-part57f-mobius-euler-next-squarefree-layer} 对 Möbius--Euler 余项的显式控制可见，冻结相在压力对偶与 Abel 有限部两侧都呈现出同一种硬壁结构：一侧把可达纤维指数锁定在 \(y\le \alpha_*\) 并在边界处留下精确代价 \(\log\varphi-g_*\)，另一侧则把有限部余项压缩为由下一平方自由层主导的首项与可闭式审计的尾界。这为进一步把谱硬壁与有限部硬余项统一到同一 Legendre--Witt 框架提供了直接接口。

\endinput

\paragraph{冻结预算相图、尾计数平台与 escort 预言机竞争律}
在冻结锥上的极大纤维骨架、连续容量曲线的 Mellin--Stieltjes 对偶与全微态位率阈值已经建立之后，还可把顶端平台刚性、容量缺口闭式与 escort 预算竞争进一步压缩为下列四条结果。它们把最大纤维指数 \(\alpha_*\)、最大层复杂度 \(g_*\)、次大纤维指数 \(\alpha_*^{(2)}\) 与预算速率 \(\beta\) 统一组织为一张可审计的冻结预算相图。

\begin{theorem}[尾计数函数的顶端平台刚性]\label{thm:xi-time-part57ea-tailcount-top-platform-rigidity}
\leanverified{paper\_xi\_time\_part57ea\_tailcount\_top\_platform\_rigidity}
沿用定理 \ref{thm:xi-time-part57b-escort-maxfiber-layer-concentration} 的记号。设
\[
\alpha_*^{(2)}<\beta<\alpha_*,
\qquad
B_m:=\left\lfloor \frac{\beta}{\log 2}\,m\right\rfloor,
\qquad
T_m:=2^{B_m}.
\]
并记连续阈值尾计数函数
\[
N_m(T):=\#\{x\in X_m:\ d_m(x)\ge T\}.
\]
再记
\[
A_m^{\max}:=\{x\in X_m:\ d_m(x)=M_m\}.
\]
则对充分大的 \(m\)，有严格恒等式
\[
\boxed{\;
N_m(T_m)=K_m.\;}
\]
换言之，只要指数阈值 \(T_m\) 最终落在 \(M_m^{(2)}\) 与 \(M_m\) 之间，则超过该阈值的纤维恰好就是全部最大纤维。
\end{theorem}
\begin{proof}
由 \(\beta<\alpha_*\) 及
\[
\frac1m\log M_m\to\alpha_*,
\]
可得对充分大的 \(m\)，
\[
\frac1m\log T_m=\frac{B_m\log 2}{m}
=
\beta+O(m^{-1})
<
\alpha_*+o(1)
=
\frac1m\log M_m,
\]
从而
\[
T_m<M_m.
\]
另一方面，由 \(\beta>\alpha_*^{(2)}\) 及
\[
\limsup_{m\to\infty}\frac1m\log M_m^{(2)}=\alpha_*^{(2)},
\]
可得对充分大的 \(m\)，
\[
\frac1m\log M_m^{(2)}<\beta+O(m^{-1})=\frac1m\log T_m,
\]
从而
\[
M_m^{(2)}<T_m.
\]
于是对 \(x\in A_m^{\max}\) 有 \(d_m(x)=M_m>T_m\)，而对 \(x\notin A_m^{\max}\) 有 \(d_m(x)\le M_m^{(2)}<T_m\)。因此
\[
N_m(T_m)=\#A_m^{\max}=K_m.
\]
证毕。
\end{proof}

\begin{theorem}[预言机容量缺口的顶端精确公式]\label{thm:xi-time-part57ea-oracle-capacity-defect-top-formula}
在定理 \ref{thm:xi-time-part57ea-tailcount-top-platform-rigidity} 的假设下，对充分大的 \(m\) 有
\[
\boxed{\;
C_m(B_m)=2^m-K_m(M_m-T_m),\;}
\]
其中
\[
C_m(B):=\sum_{x\in X_m}\min(d_m(x),2^B),
\qquad
\mathrm{Succ}_m(B):=\frac{C_m(B)}{2^m}.
\]
从而
\[
\boxed{\;
1-\mathrm{Succ}_m(B_m)
=
\frac{K_m(M_m-T_m)}{2^m}.\;}
\]
并且
\[
\boxed{\;
\lim_{m\to\infty}\frac1m
\log\bigl(1-\mathrm{Succ}_m(B_m)\bigr)
=
\alpha_*+g_*-\log 2.\;}
\]
\end{theorem}
\begin{proof}
由定理 \ref{thm:xi-time-part57ea-tailcount-top-platform-rigidity}，对充分大的 \(m\) 有
\[
M_m^{(2)}<T_m<M_m.
\]
于是对 \(x\notin A_m^{\max}\) 有 \(\min(d_m(x),T_m)=d_m(x)\)，对 \(x\in A_m^{\max}\) 有 \(\min(d_m(x),T_m)=T_m\)。故
\[
C_m(B_m)
=
\sum_{x\notin A_m^{\max}}d_m(x)+\sum_{x\in A_m^{\max}}T_m.
\]
再用总质量恒等式 \(\sum_x d_m(x)=2^m\)，得到
\[
C_m(B_m)=2^m-K_mM_m+K_mT_m
=
2^m-K_m(M_m-T_m).
\]
两边除以 \(2^m\) 即得成功率缺口公式。

最后，由 \(\beta<\alpha_*\) 可知
\[
\frac{T_m}{M_m}
=
\exp\!\bigl(-(\alpha_*-\beta)m+o(m)\bigr)\to 0,
\]
故
\[
M_m-T_m=M_m(1+o(1)).
\]
再结合
\[
M_m=\exp(\alpha_*m+o(m)),
\qquad
K_m=\exp(g_*m+o(m)),
\]
得到
\[
1-\mathrm{Succ}_m(B_m)
=
\exp\!\bigl((\alpha_*+g_*-\log 2)m+o(m)\bigr),
\]
从而结论成立。证毕。
\end{proof}

\begin{theorem}[冻结相 escort 预言机成功率的精确分解与竞争上界]\label{thm:xi-time-part57ea-frozen-escort-oracle-success-decomposition}
设 \(a>a_{\mathrm c}\)，并沿用定理 \ref{thm:xi-fixed-freezing-escort-renyi-spectrum-collapse} 的 escort 分布
\[
\pi_{m,a}(x):=\frac{d_m(x)^a}{S_a(m)}.
\]
定义 escort 预算 \(B\) 下的预言机成功率
\[
\mathrm{Succ}_{m,a}^{\mathrm{esc}}(B)
:=
\sum_{x\in X_m}\pi_{m,a}(x)\min\!\left(1,\frac{2^B}{d_m(x)}\right).
\]
仍取
\[
\alpha_*^{(2)}<\beta<\alpha_*,
\qquad
B_m=\left\lfloor \frac{\beta}{\log 2}\,m\right\rfloor,
\qquad
T_m=2^{B_m}.
\]
则对充分大的 \(m\)，有严格分解
\[
\boxed{\;
\mathrm{Succ}_{m,a}^{\mathrm{esc}}(B_m)
=
\frac{T_m}{M_m}\,\pi_{m,a}(A_m^{\max})
+
\bigl(1-\pi_{m,a}(A_m^{\max})\bigr).\;}
\]
从而
\[
\boxed{\;
\mathrm{Succ}_{m,a}^{\mathrm{esc}}(B_m)
=
\frac{T_m}{M_m}
+
O\!\bigl(e^{-m\Delta(a)+o(m)}\bigr),\;}
\]
其中
\[
\Delta(a):=a(\alpha_*-\alpha_*^{(2)})-(\log\varphi-g_*)>0.
\]
特别地，
\[
\boxed{\;
\limsup_{m\to\infty}\frac1m
\log \mathrm{Succ}_{m,a}^{\mathrm{esc}}(B_m)
\le
-\min\{\alpha_*-\beta,\Delta(a)\}.\;}
\]
并且总有下界
\[
\boxed{\;
\mathrm{Succ}_{m,a}^{\mathrm{esc}}(B_m)
\ge
\exp\!\bigl(-(\alpha_*-\beta)m+o(m)\bigr).\;}
\]
\end{theorem}
\begin{proof}
由定理 \ref{thm:xi-time-part57ea-tailcount-top-platform-rigidity}，对充分大的 \(m\) 只有最大纤维满足 \(d_m(x)>T_m\)，而对 \(x\notin A_m^{\max}\) 都有 \(d_m(x)\le T_m\)。因此
\[
\min\!\left(1,\frac{T_m}{d_m(x)}\right)
=
\begin{cases}
T_m/M_m, & x\in A_m^{\max},\\
1, & x\notin A_m^{\max}.
\end{cases}
\]
代回定义即得
\[
\mathrm{Succ}_{m,a}^{\mathrm{esc}}(B_m)
=
\frac{T_m}{M_m}\,\pi_{m,a}(A_m^{\max})
+
\bigl(1-\pi_{m,a}(A_m^{\max})\bigr).
\]

再由定理 \ref{thm:xi-fixed-freezing-escort-renyi-spectrum-collapse} 第四条，
\[
1-\pi_{m,a}(A_m^{\max})
\le
\exp\!\bigl(-m\Delta(a)+o(m)\bigr).
\]
于是
\[
\mathrm{Succ}_{m,a}^{\mathrm{esc}}(B_m)
=
\frac{T_m}{M_m}
+
\left(1-\frac{T_m}{M_m}\right)
\bigl(1-\pi_{m,a}(A_m^{\max})\bigr),
\]
从而
\[
\mathrm{Succ}_{m,a}^{\mathrm{esc}}(B_m)
=
\frac{T_m}{M_m}
+
O\!\bigl(e^{-m\Delta(a)+o(m)}\bigr).
\]
又因为
\[
\frac{T_m}{M_m}
=
\exp\!\bigl(-(\alpha_*-\beta)m+o(m)\bigr),
\]
故上式立即给出
\[
\limsup_{m\to\infty}\frac1m
\log \mathrm{Succ}_{m,a}^{\mathrm{esc}}(B_m)
\le
-\min\{\alpha_*-\beta,\Delta(a)\}.
\]

另一方面，由 \(\pi_{m,a}(A_m^{\max})\to 1\) 可知对充分大的 \(m\) 有 \(\pi_{m,a}(A_m^{\max})\ge \tfrac12\)。因此
\[
\mathrm{Succ}_{m,a}^{\mathrm{esc}}(B_m)
\ge
\frac{T_m}{M_m}\,\pi_{m,a}(A_m^{\max})
\ge
\frac12\,\frac{T_m}{M_m},
\]
即得到所示下界。证毕。
\end{proof}

\begin{theorem}[冻结预算相图中的 oracle--escort 临界线]\label{thm:xi-time-part57ea-oracle-escort-critical-line}
在定理 \ref{thm:xi-time-part57ea-frozen-escort-oracle-success-decomposition} 的设定下，定义
\[
\beta_{\mathrm c}(a)
:=
\alpha_*-\Delta(a)
=
\alpha_*^{(2)}+\frac{\log\varphi-g_*}{a}.
\]
则成立：
\begin{enumerate}
  \item 若
  \[
  \beta_{\mathrm c}(a)<\beta<\alpha_*,
  \]
  则 oracle 截断项严格主导，并且
  \[
  \boxed{\;
  \mathrm{Succ}_{m,a}^{\mathrm{esc}}(B_m)
  =
  \exp\!\bigl(-(\alpha_*-\beta)m+o(m)\bigr).\;}
  \]
  \item 若
  \[
  \alpha_*^{(2)}<\beta\le \beta_{\mathrm c}(a),
  \]
  则有统一上界
  \[
  \boxed{\;
  \mathrm{Succ}_{m,a}^{\mathrm{esc}}(B_m)
  \le
  \exp\!\bigl(-\Delta(a)m+o(m)\bigr).\;}
  \]
  在临界点 \(\beta=\beta_{\mathrm c}(a)\) 处，oracle 项与 escort 泄漏上界恰处于同一指数尺度。
  \item 若冻结相的 min-entropy 率存在，则由定理 \ref{thm:xi-fixed-freezing-escort-renyi-spectrum-collapse} 的第三条，
  \[
  h_\infty(a)=g_*,
  \]
  因而临界线可重写为
  \[
  \boxed{\;
  \beta_{\mathrm c}(a)
  =
  \alpha_*^{(2)}+\frac{\log\varphi-h_\infty(a)}{a}.\;}
  \]
\end{enumerate}
\end{theorem}
\begin{proof}
由定义
\[
\alpha_*-\beta_{\mathrm c}(a)=\Delta(a).
\]
若 \(\beta>\beta_{\mathrm c}(a)\)，则
\[
\alpha_*-\beta<\Delta(a).
\]
于是定理 \ref{thm:xi-time-part57ea-frozen-escort-oracle-success-decomposition} 中
\[
\frac{T_m}{M_m}
=
\exp\!\bigl(-(\alpha_*-\beta)m+o(m)\bigr)
\]
比误差项
\(
O(e^{-m\Delta(a)+o(m)})
\)
衰减更慢，从而成为主导项，故得到第一条。

若 \(\beta\le \beta_{\mathrm c}(a)\)，则
\[
\alpha_*-\beta\ge \Delta(a),
\]
故
\[
\frac{T_m}{M_m}
=
O\!\bigl(e^{-m\Delta(a)+o(m)}\bigr).
\]
再与定理 \ref{thm:xi-time-part57ea-frozen-escort-oracle-success-decomposition} 中的误差控制合并，即得第二条上界；当 \(\beta=\beta_{\mathrm c}(a)\) 时，两项指数相同。

第三条只需把 \(g_*\) 用
\[
h_\infty(a)=g_*
\]
替换即可。证毕。
\end{proof}

\paragraph{尾计数指数、冻结阈值与 oracle 前沿的变分闭包}
在定理 \ref{thm:xi-time-part57b-capacity-moment-exact-mellin-stieltjes} 已把幂和矩写成尾计数与容量曲线的 Mellin--Stieltjes 积分、定理 \ref{thm:xi-time-part57ea-tailcount-top-platform-rigidity}--\ref{thm:xi-time-part57ea-oracle-escort-critical-line} 已把冻结预算相图压缩为最大纤维平台律之后，还可进一步把尾计数指数函数、压力谱、均匀采样下的对数多重度大偏差、尺寸偏置倾斜以及连续容量指数统一组织为一套纯变分闭包。该闭包把前文的充分条件提升为必要且充分的边界公式，并把预算前沿、倾斜前沿与冻结前沿一并置于同一尾指数函数 \(\Psi\) 的支撑之上。

\begin{theorem}[尾计数指数的 Laplace--Mellin 变分原理]\label{thm:xi-time-part57eb-tail-laplace-mellin-variational}
设存在上半连续函数
\[
\Psi:[0,\alpha_*]\to[-\infty,\log\varphi]
\]
满足：对任意紧区间 \(I\subset[0,\alpha_*]\)，当 \(m\to\infty\) 时都有一致指数渐近
\[
N_m(e^{\alpha m})
=
\exp\!\bigl(m\Psi(\alpha)+o(m)\bigr)
\qquad
(\alpha\in I),
\]
其中
\[
N_m(T):=\#\{x\in X_m:\ d_m(x)\ge T\}.
\]
则对每个实数 \(a>0\)，压力满足
\[
\boxed{\;
P_a=\sup_{0\le \alpha\le \alpha_*}\bigl(a\alpha+\Psi(\alpha)\bigr).\;}
\]
\end{theorem}
\begin{proof}
由定理 \ref{thm:xi-time-part57b-capacity-moment-exact-mellin-stieltjes}，
\[
S_a(m)=a\int_0^\infty T^{a-1}N_m(T)\,\dd T.
\]
又因 \(d_m(x)\le M_m\) 对所有 \(x\in X_m\) 成立，故
\[
S_a(m)=a\int_0^{M_m}T^{a-1}N_m(T)\,\dd T.
\]
令 \(T=e^{\alpha m}\)，则 \(\dd T=me^{\alpha m}\dd\alpha\)，从而
\[
S_a(m)
=
am\int_{-\infty}^{\frac1m\log M_m}
\exp(am\alpha)\,N_m(e^{\alpha m})\,\dd\alpha.
\]

当 \(\alpha<0\) 时，\(e^{\alpha m}<1\)，而 \(d_m(x)\in\ZZ_{\ge 1}\)，故
\[
N_m(e^{\alpha m})=|X_m|
=
\exp\!\bigl((\log\varphi)m+o(m)\bigr).
\]
于是该区间上的指数至多为 \(a\alpha+\log\varphi\le \log\varphi\)，不会优于 \(\alpha=0\) 处的候选值。另一方面，
\[
\frac1m\log M_m\to \alpha_*.
\]
因此只需在 \([0,\alpha_*]\) 上分析积分主项。

由假设，对每个紧区间 \(I\subset[0,\alpha_*]\) 都有
\[
N_m(e^{\alpha m})=\exp\!\bigl(m\Psi(\alpha)+o(m)\bigr)
\qquad(\alpha\in I)
\]
且 \(o(m)\) 对 \(\alpha\in I\) 一致。将 \([0,\alpha_*]\) 以有限个小区间覆盖，并对每段分别应用 Laplace 原理，再利用 \(\Psi\) 的上半连续性与有限覆盖取极大值，得到
\[
S_a(m)
=
\exp\!\Bigl(
m\sup_{0\le \alpha\le \alpha_*}(a\alpha+\Psi(\alpha))+o(m)
\Bigr).
\]
两边取 \(\frac1m\log\) 并令 \(m\to\infty\)，即得
\[
P_a=\sup_{0\le \alpha\le \alpha_*}(a\alpha+\Psi(\alpha)).
\]
证毕。
\end{proof}

\input{subsubsec__xi-time-protocol-conclusions-part57ec-uniform-tilt-oracle-gap-ldp}

\begin{theorem}[连续容量曲线的变分公式与成功率指数]\label{thm:xi-time-part57eb-continuous-capacity-success-variational}
定义连续容量曲线
\[
\mathcal C_m^{\mathrm{cont}}(T):=\sum_{x\in X_m}\min(d_m(x),T)
\qquad(T\ge 0),
\]
以及预算指数
\[
Q(\beta):=\limsup_{m\to\infty}\frac1m\log \mathcal C_m^{\mathrm{cont}}(e^{\beta m})
\qquad(\beta\ge 0).
\]
在定理 \ref{thm:xi-time-part57eb-tail-laplace-mellin-variational} 的同一假设下，对每个 \(\beta\ge 0\) 都有
\[
\boxed{\;
Q(\beta)=\sup_{0\le \alpha\le \min(\beta,\alpha_*)}\bigl(\alpha+\Psi(\alpha)\bigr).\;}
\]
进一步，若取离散预算
\[
B_m:=\left\lfloor \frac{\beta m}{\log 2}\right\rfloor,
\]
并记
\[
\mathrm{Succ}_m(B):=\frac{C_m(B)}{2^m},
\qquad
C_m(B):=\sum_{x\in X_m}\min(d_m(x),2^B),
\]
则有
\[
\boxed{\;
\limsup_{m\to\infty}\frac1m\log \mathrm{Succ}_m(B_m)
=
Q(\beta)-\log 2.\;}
\]
\end{theorem}
\begin{proof}
由定理 \ref{thm:xi-capacity-tail-histogram-moment-complete-statistic}，
\[
\mathcal C_m^{\mathrm{cont}}(T)=\int_0^T N_m(t)\,\dd t.
\]
取 \(T=e^{\beta m}\)，再作变量代换 \(t=e^{\alpha m}\)，得到
\[
\mathcal C_m^{\mathrm{cont}}(e^{\beta m})
=
m\int_{-\infty}^{\beta}\exp(\alpha m)\,N_m(e^{\alpha m})\,\dd\alpha.
\]

当 \(\alpha<0\) 时，仍有 \(N_m(e^{\alpha m})=|X_m|=\exp((\log\varphi)m+o(m))\)，故对应指数不超过 \(\log\varphi+\alpha\le \log\varphi\)，不会优于 \(\alpha=0\) 的候选值。又由于 \(N_m(t)=0\) 当 \(t>M_m\)，并且 \(\frac1m\log M_m\to\alpha_*\)，故只需在 \([0,\min(\beta,\alpha_*)]\) 上取主项。由定理 \ref{thm:xi-time-part57eb-tail-laplace-mellin-variational} 的同一一致指数假设，对该积分再施行 Laplace 原理即得
\[
Q(\beta)=\sup_{0\le \alpha\le \min(\beta,\alpha_*)}\bigl(\alpha+\Psi(\alpha)\bigr).
\]

另一方面，由命题 \ref{prop:fold-capacity-tail-equivalence}，
\[
C_m(B)=\mathcal C_m^{\mathrm{cont}}(2^B)
\qquad(B\in\ZZ_{\ge 0}).
\]
对 \(B=B_m\) 有
\[
2^{B_m}=e^{\beta m+O(1)},
\]
故
\[
\limsup_{m\to\infty}\frac1m\log C_m(B_m)=Q(\beta).
\]
再由 \(\mathrm{Succ}_m(B_m)=C_m(B_m)/2^m\)，得到
\[
\limsup_{m\to\infty}\frac1m\log \mathrm{Succ}_m(B_m)=Q(\beta)-\log 2.
\]
证毕。
\end{proof}

\begin{theorem}[冻结阈值的精确变分刻画]\label{thm:xi-time-part57eb-freezing-threshold-exact-variational}
设在区间 \((\alpha_*^{(2)},\alpha_*)\) 上存在尾计数平台
\[
\Psi(\alpha)=g_*
\qquad
(\alpha_*^{(2)}<\alpha<\alpha_*).
\]
定义
\[
a_{\mathrm{fr}}
:=
\sup_{0\le \alpha<\alpha_*}
\frac{\Psi(\alpha)-g_*}{\alpha_*-\alpha}.
\]
则对任意实数 \(a>0\)，有等价判别
\[
\boxed{\;
P_a=a\alpha_*+g_*
\iff
a\ge a_{\mathrm{fr}}.\;}
\]
并且
\[
\boxed{\;
a_{\mathrm{fr}}
\le
\frac{\log\varphi-g_*}{\alpha_*-\alpha_*^{(2)}}.\;}
\]
因此前文的阈值
\[
a_{\mathrm c}:=\frac{\log\varphi-g_*}{\alpha_*-\alpha_*^{(2)}}
\]
不仅是充分条件，而且是精确冻结阈值的统一上界。
\end{theorem}
\begin{proof}
由定理 \ref{thm:xi-time-part57eb-tail-laplace-mellin-variational}，
\[
P_a=\sup_{0\le \alpha\le \alpha_*}\bigl(a\alpha+\Psi(\alpha)\bigr).
\]
由平台假设，对每个 \(\alpha\in(\alpha_*^{(2)},\alpha_*)\) 都有
\[
a\alpha+\Psi(\alpha)=a\alpha+g_*.
\]
令 \(\alpha\uparrow\alpha_*\)，得到
\[
\sup_{\alpha_*^{(2)}<\alpha<\alpha_*}\bigl(a\alpha+\Psi(\alpha)\bigr)
=
a\alpha_*+g_*.
\]
因此
\[
P_a=a\alpha_*+g_*
\]
成立，当且仅当对所有 \(0\le \alpha<\alpha_*\) 都有
\[
a\alpha+\Psi(\alpha)\le a\alpha_*+g_*,
\]
等价于
\[
a\ge \frac{\Psi(\alpha)-g_*}{\alpha_*-\alpha}.
\]
对全部 \(0\le \alpha<\alpha_*\) 取上确界，即得
\[
P_a=a\alpha_*+g_*
\iff
a\ge a_{\mathrm{fr}}.
\]

再由平台假设，当 \(\alpha\in(\alpha_*^{(2)},\alpha_*)\) 时，
\[
\Psi(\alpha)=g_*,
\]
故该区间对上确界没有正贡献。于是
\[
a_{\mathrm{fr}}
=
\sup_{0\le \alpha\le \alpha_*^{(2)}}
\frac{\Psi(\alpha)-g_*}{\alpha_*-\alpha}.
\]
另一方面，对所有 \(\alpha\in[0,\alpha_*]\) 都有 \(\Psi(\alpha)\le \log\varphi\)，从而
\[
a_{\mathrm{fr}}
\le
\sup_{0\le \alpha\le \alpha_*^{(2)}}
\frac{\log\varphi-g_*}{\alpha_*-\alpha}
=
\frac{\log\varphi-g_*}{\alpha_*-\alpha_*^{(2)}}.
\]
证毕。
\end{proof}

\begin{theorem}[oracle--escort 双前沿的精确相图]\label{thm:xi-time-part57eb-oracle-escort-double-frontier}
保持定理 \ref{thm:xi-time-part57eb-freezing-threshold-exact-variational} 的尾计数平台假设。定义连续容量指数
\[
Q(\beta)=\sup_{0\le \alpha\le \min(\beta,\alpha_*)}\bigl(\alpha+\Psi(\alpha)\bigr),
\]
以及预算饱和阈值
\[
\beta_{\mathrm{sat}}
:=
\sup_{0\le \alpha\le \alpha_*^{(2)}}\bigl(\alpha+\Psi(\alpha)-g_*\bigr).
\]
则对每个 \(\beta\in(\alpha_*^{(2)},\alpha_*)\)，成立如下分类：
\begin{enumerate}
  \item 若 \(\beta\ge \beta_{\mathrm{sat}}\)，则
  \[
  \boxed{\;
  Q(\beta)=\beta+g_*.\;}
  \]
  因而对
  \[
  B_m=\left\lfloor \frac{\beta m}{\log 2}\right\rfloor
  \]
  有
  \[
  \limsup_{m\to\infty}\frac1m\log \mathrm{Succ}_m(B_m)
  =
  \beta+g_*-\log 2.
  \]
  \item 若 \(\beta<\beta_{\mathrm{sat}}\)，则
  \[
  \boxed{\;
  Q(\beta)=
  \sup_{0\le \alpha\le \alpha_*^{(2)}}\bigl(\alpha+\Psi(\alpha)\bigr).\;}
  \]
  换言之，容量指数仍由次大尺度区间决定，尚未进入最大纤维主导相。
  \item 总有显式上界
  \[
  \boxed{\;
  \beta_{\mathrm{sat}}
  \le
  \alpha_*^{(2)}+\log\varphi-g_*.\;}
  \]
  因此只要
  \[
  \beta>\alpha_*^{(2)}+\log\varphi-g_*,
  \]
  就必进入线性饱和相 \(Q(\beta)=\beta+g_*\)。
\end{enumerate}
\end{theorem}
\begin{proof}
对 \(\beta\in(\alpha_*^{(2)},\alpha_*)\)，由平台性质，
\[
\Psi(\alpha)=g_*
\qquad
(\alpha_*^{(2)}<\alpha\le \beta).
\]
故在该区间上
\[
\alpha+\Psi(\alpha)=\alpha+g_*
\]
严格随 \(\alpha\) 递增，其最大值为 \(\beta+g_*\)。于是
\[
Q(\beta)
=
\max\Bigl\{
\beta+g_*,
\sup_{0\le \alpha\le \alpha_*^{(2)}}(\alpha+\Psi(\alpha))
\Bigr\}.
\]

若 \(\beta\ge \beta_{\mathrm{sat}}\)，则按定义有
\[
\beta+g_*
\ge
\sup_{0\le \alpha\le \alpha_*^{(2)}}(\alpha+\Psi(\alpha)),
\]
从而
\[
Q(\beta)=\beta+g_*.
\]
再由定理 \ref{thm:xi-time-part57eb-continuous-capacity-success-variational}，得到
\[
\limsup_{m\to\infty}\frac1m\log \mathrm{Succ}_m(B_m)
=
Q(\beta)-\log 2
=
\beta+g_*-\log 2.
\]

若 \(\beta<\beta_{\mathrm{sat}}\)，则
\[
\beta+g_*
<
\sup_{0\le \alpha\le \alpha_*^{(2)}}(\alpha+\Psi(\alpha)),
\]
故
\[
Q(\beta)=
\sup_{0\le \alpha\le \alpha_*^{(2)}}(\alpha+\Psi(\alpha)).
\]
这就给出第二种相。

最后，由 \(\Psi(\alpha)\le \log\varphi\)，有
\[
\beta_{\mathrm{sat}}
=
\sup_{0\le \alpha\le \alpha_*^{(2)}}(\alpha+\Psi(\alpha)-g_*)
\le
\sup_{0\le \alpha\le \alpha_*^{(2)}}(\alpha+\log\varphi-g_*)
=
\alpha_*^{(2)}+\log\varphi-g_*.
\]
故当
\(
\beta>\alpha_*^{(2)}+\log\varphi-g_*
\)
时，必有 \(\beta>\beta_{\mathrm{sat}}\)，从而落入线性饱和相。证毕。
\end{proof}

\input{subsubsec__xi-time-protocol-conclusions-part57ed-oracle-gap-min-envelope-conjecture}

\endinput


\endinput


\begin{theorem}[分枝三次场、零温四次场与 RH 临界五次场的算术独立]\label{thm:xi-branch-rh-zero-temp-arithmetic-independence}
设
\[
\Delta(u):=4u^3+25u^2+88u+8,
\]
\[
R_{\mathrm{gr}}(y):=y^4-6y^3+9y^2-y-1,
\]
\[
Q_{\mathrm{RH}}(u):=u^5+4u^4+3u^3-96u^2-42u-14.
\]
分别记其分裂域为 \(K_\Delta\)、\(K_{\mathrm{gr}}\)、\(K_{\mathrm{RH}}\)。则
\[
\boxed{\;
\Gal(K_\Delta/\QQ)\cong S_3,\qquad
\Gal(K_{\mathrm{gr}}/\QQ)\cong S_4,\qquad
\Gal(K_{\mathrm{RH}}/\QQ)\cong S_5.\;}
\]
并且三者在 \(\QQ\) 上算术独立，满足
\[
\boxed{\;
\Gal(K_\Delta K_{\mathrm{gr}}K_{\mathrm{RH}}/\QQ)
\cong
S_3\times S_4\times S_5.\;}
\]
等价地，分枝几何、零温地基与 RH 临界阈值并非同一伽罗瓦源头的不同投影，而对应三个彼此正交的对称场。
\end{theorem}
\begin{proof}
先看三次多项式 \(\Delta(u)\)。由有理根判别法，\(\Delta\) 在 \(\QQ\) 上不可约。其判别式为
\[
\Disc(\Delta)=-2^5\cdot 5^3\cdot 11^3,
\]
非平方，故其分裂域伽罗瓦群只能为 \(S_3\) 而非 \(A_3\)。于是
\[
\Gal(K_\Delta/\QQ)\cong S_3,
\qquad
K_\Delta^{A_3}=\QQ(\sqrt{-110}).
\]
\par
再看四次多项式 \(R_{\mathrm{gr}}(y)\)。模 \(3\) 约化后 \(R_{\mathrm{gr}}\) 仍不可约，因此分裂域伽罗瓦群在 \(S_4\) 中传递并含一个 \(4\)-循环。其判别式为
\[
\Disc(R_{\mathrm{gr}})=2777,
\]
非平方，故群不包含于 \(A_4\)。另一方面，模 \(11\) 的分解型为 \([1,3]\)，故 Frobenius 元给出一个 \(3\)-循环。传递子群 \(G\le S_4\) 若同时含 \(4\)-循环与 \(3\)-循环，则只能是 \(S_4\)。故
\[
\Gal(K_{\mathrm{gr}}/\QQ)\cong S_4,
\qquad
K_{\mathrm{gr}}^{A_4}=\QQ(\sqrt{2777}).
\]
\par
对五次多项式 \(Q_{\mathrm{RH}}(u)\)，模 \(5\) 约化后仍不可约，因此分裂域伽罗瓦群在 \(S_5\) 中传递并含一个 \(5\)-循环。其判别式为
\[
\Disc(Q_{\mathrm{RH}})
=
2^6\cdot 3^2\cdot 7\cdot 743\cdot 1693^2,
\]
非平方，故群不包含于 \(A_5\)。又模 \(17\) 的分解型为 \([1,1,3]\)，故 Frobenius 元给出一个 \(3\)-循环。按五次传递子群分类，含 \(5\)-循环与 \(3\)-循环的传递子群只能为 \(A_5\) 或 \(S_5\)；判别式非平方排除 \(A_5\)，从而
\[
\Gal(K_{\mathrm{RH}}/\QQ)\cong S_5,
\qquad
K_{\mathrm{RH}}^{A_5}=\QQ(\sqrt{5201}).
\]
\par
下面证明独立性。先看 \(K_\Delta\) 与 \(K_{\mathrm{gr}}\)。二者交域是 \(\QQ\) 上的 Galois 扩张，其 Galois 群必须同时是 \(S_3\) 与 \(S_4\) 的商群，因而只能为 \(1\)、\(C_2\) 或 \(S_3\)。若为 \(C_2\)，则二者的唯一二次子域必须相同，但
\[
\QQ(\sqrt{-110})\neq \QQ(\sqrt{2777}),
\]
矛盾。若为 \(S_3\)，则 \(K_\Delta\subseteq K_{\mathrm{gr}}\)，从而 \(K_\Delta\) 必等于 \(K_{\mathrm{gr}}^{V_4}\) 这一 \(S_3\)-商域；然而其唯一二次子域应为
\[
K_{\mathrm{gr}}^{A_4}=\QQ(\sqrt{2777}),
\]
仍与 \(\QQ(\sqrt{-110})\) 矛盾。故
\[
K_\Delta\cap K_{\mathrm{gr}}=\QQ
\]
并得到
\[
\Gal(K_\Delta K_{\mathrm{gr}}/\QQ)\cong S_3\times S_4.
\]
\par
再看 \(K_{\mathrm{RH}}\) 与前两者。首先，\(S_5\) 的非平凡真正规商只有 \(C_2\)，故
\[
K_\Delta\cap K_{\mathrm{RH}}=\QQ,
\qquad
K_{\mathrm{gr}}\cap K_{\mathrm{RH}}=\QQ
\]
只需排除二次子域相同的可能，而
\[
\QQ(\sqrt{5201})
\neq
\QQ(\sqrt{-110}),
\qquad
\QQ(\sqrt{5201})
\neq
\QQ(\sqrt{2777}).
\]
接着考虑与合成域 \(K_\Delta K_{\mathrm{gr}}\) 的交。由于
\[
\bigl(S_3\times S_4\bigr)^{\mathrm{ab}}\cong C_2\times C_2,
\]
其全部二次子域恰由三条非零特征给出，即
\[
\QQ(\sqrt{-110}),\qquad
\QQ(\sqrt{2777}),\qquad
\QQ(\sqrt{-110\cdot 2777}).
\]
这些二次域都不同于 \(\QQ(\sqrt{5201})\)。而 \(S_5\) 侧任何非平凡 Galois 交域只能是其唯一二次子域 \(\QQ(\sqrt{5201})\)，故
\[
K_{\mathrm{RH}}\cap K_\Delta K_{\mathrm{gr}}=\QQ.
\]
于是
\[
\Gal(K_\Delta K_{\mathrm{gr}}K_{\mathrm{RH}}/\QQ)
\cong
\Gal(K_\Delta/\QQ)\times\Gal(K_{\mathrm{gr}}/\QQ)\times\Gal(K_{\mathrm{RH}}/\QQ)
\cong
S_3\times S_4\times S_5.
\]
可复现核验脚本：\texttt{scripts/exp\_xi\_branch\_rh\_zero\_temp\_galois\_chebotarev\_audit.py}。证毕。
\end{proof}

\begin{corollary}[三重 Chebotarev 独立律与显式联合密度]\label{cor:xi-branch-rh-zero-temp-chebotarev-product-densities}
\leanverified{paper\_xi\_branch\_rh\_zero\_temp\_chebotarev\_product\_densities}
沿用定理 \ref{thm:xi-branch-rh-zero-temp-arithmetic-independence} 的记号。对一切不分歧于
\[
\Disc(\Delta)\Disc(R_{\mathrm{gr}})\Disc(Q_{\mathrm{RH}})
\]
的素数 \(p\)，三条多项式
\[
\Delta(u),\qquad R_{\mathrm{gr}}(y),\qquad Q_{\mathrm{RH}}(u)
\]
在 \(\FF_p\) 上的分解型彼此独立；更精确地，对任意共轭类
\[
C_3\subset S_3,\qquad C_4\subset S_4,\qquad C_5\subset S_5,
\]
满足
\[
\delta\Bigl(\mathrm{Frob}_p\in C_3\times C_4\times C_5\Bigr)
=
\frac{|C_3|}{|S_3|}\cdot \frac{|C_4|}{|S_4|}\cdot \frac{|C_5|}{|S_5|}.
\]
特别地：
\begin{enumerate}
  \item 三者同时不可约的素数密度为
  \[
  \boxed{\;
  \frac{2}{6}\cdot\frac{6}{24}\cdot\frac{24}{120}
  =
  \frac{1}{60}.\;}
  \]
  \item 三者同时完全分裂的素数密度为
  \[
  \boxed{\;
  \frac{1}{6}\cdot\frac{1}{24}\cdot\frac{1}{120}
  =
  \frac{1}{17280}.\;}
  \]
  \item \(\Delta\) 不可约、\(R_{\mathrm{gr}}\) 具有 \((3+1)\) 分解、且 \(Q_{\mathrm{RH}}\) 不可约的联合密度为
  \[
  \boxed{\;
  \frac{2}{6}\cdot\frac{8}{24}\cdot\frac{24}{120}
  =
  \frac{1}{45},\;}
  \]
  因为 \(S_4\) 中 \((3+1)\) 型共轭类大小为 \(8\)。
\end{enumerate}
\end{corollary}
\begin{proof}
定理 \ref{thm:xi-branch-rh-zero-temp-arithmetic-independence} 已给出
\[
\Gal(K_\Delta K_{\mathrm{gr}}K_{\mathrm{RH}}/\QQ)\cong S_3\times S_4\times S_5.
\]
Chebotarev 密度定理说明：对不分歧素数，Frobenius 共轭类在该直积群中按计数测度等分布。故任意联合分解型事件的密度都等于三个因子中对应共轭类比例的乘积。
\par
另一方面，多项式在 \(\FF_p\) 上的分解型与 Frobenius 元的循环型一致：三次不可约对应 \(S_3\) 中的 \(3\)-循环，四次不可约对应 \(S_4\) 中的 \(4\)-循环，五次不可约对应 \(S_5\) 中的 \(5\)-循环，而四次的 \((3+1)\) 分解对应 \(S_4\) 中大小为 \(8\) 的 \((3+1)\) 共轭类。将这些类大小代入上式，即得到三条显式密度。证毕。
\end{proof}

\paragraph{原子长度二残差、根单位滤波与 Lee--Yang 符号坍缩}
下列条目把定理 \ref{thm:xi-real-input-40-atomic-witt-even-odd-surgery}、
\ref{thm:xi-atomic-witt-mod-class-localization}、定理
\ref{thm:xi-dirichlet-mode-forces-congruence-residual}、推论
\ref{cor:fold-gauge-anomaly-P10-leyang-chebotarev-product-law} 与结论
\ref{con:xi-terminal-zm-leyang-cubic-modp-splitting-density} 重新压缩为同一条
“原子剥离--同余注入--core 门槛--算术筛选”链；并且该链与
\S\ref{subsubsec:zeta-online-kernel-length-modq-chebotarev} 的 root-of-unity 长度滤波接口完全对齐。

\begin{theorem}[原子 Witt 因子的周期--primitive 精确剥离]\label{thm:xi-atomic-witt-cycle-primitive-exact-peeling}
\leanverified{paper\_xi\_atomic\_witt\_cycle\_primitive\_exact\_peeling}
沿用附录 \ref{app:real-input-40-zeta-u} 的记号。设
\[
\zeta(z,u)=\Delta(z,u)^{-1},
\qquad
\Delta(z,u)=(1-uz^2)\,F(z,u),
\qquad
\zeta_{\mathrm{core}}(z,u):=F(z,u)^{-1}.
\]
定义周期对数系数
\[
\log\zeta(z,u)=\sum_{n\ge 1}\frac{a_n(u)}{n}z^n,
\qquad
\log\zeta_{\mathrm{core}}(z,u)=\sum_{n\ge 1}\frac{a_n^{\mathrm{core}}(u)}{n}z^n,
\]
并定义 primitive 生成函数
\[
P(z,u)=\sum_{n\ge 1}p_n(u)z^n,
\qquad
P_{\mathrm{core}}(z,u)=\sum_{n\ge 1}p_n^{\mathrm{core}}(u)z^n.
\]
则对一切 \(n\ge 1\) 都有严格恒等式
\[
\boxed{\;
a_n(u)=a_n^{\mathrm{core}}(u)+2u^{n/2}\mathbf 1_{\{2\mid n\}}\;}
\]
以及
\[
\boxed{\;
p_n(u)=p_n^{\mathrm{core}}(u)+u\,\mathbf 1_{\{n=2\}}.\;}
\]
因此，原子因子在 primitive 层只占据长度 \(2\) 的单一坐标，而在周期/ghost 层则恰好沿全部偶长度产生几何级数修正。
\end{theorem}
\begin{proof}
由
\[
\zeta(z,u)=(1-uz^2)^{-1}\zeta_{\mathrm{core}}(z,u)
\]
取对数得
\[
\log\zeta(z,u)
=
-\log(1-uz^2)+\log\zeta_{\mathrm{core}}(z,u)
=
\sum_{k\ge 1}\frac{u^k}{k}z^{2k}
+\log\zeta_{\mathrm{core}}(z,u).
\]
比较 \(z^n/n\) 的系数，便得到
\[
a_n(u)-a_n^{\mathrm{core}}(u)=
\begin{cases}
2u^{n/2},&2\mid n,\\
0,&2\nmid n.
\end{cases}
\]
另一方面，命题 \ref{prop:atomic-2-orbit-split-logM} 已给出精确 primitive 剥离
\[
P(z,u)=uz^2+P_{\mathrm{core}}(z,u),
\]
逐项比较系数即得
\[
p_n(u)=p_n^{\mathrm{core}}(u)+u\,\mathbf 1_{\{n=2\}}.
\]
证毕。
\end{proof}

\begin{corollary}[无权核相对 core 恰多出唯一一个 primitive 二步类]\label{cor:xi-atomic-witt-unweighted-unique-two-step-class}
沿用定理 \ref{thm:xi-atomic-witt-cycle-primitive-exact-peeling} 的记号，并取 \(u=1\)。则
\[
p_n(1)=p_n^{\mathrm{core}}(1)+\mathbf 1_{\{n=2\}},
\qquad
a_n(1)=a_n^{\mathrm{core}}(1)+2\,\mathbf 1_{\{2\mid n\}}.
\]
因此完整系统相对于 core 恰多出唯一一个 primitive 长度 \(2\) 轨道类；其在 periodic 层上的全部额外贡献恰为偶长度上的二步原子反复：
\[
a_{2m}(1)-a_{2m}^{\mathrm{core}}(1)=2,
\qquad
a_{2m+1}(1)-a_{2m+1}^{\mathrm{core}}(1)=0
\qquad (m\ge 1).
\]
\end{corollary}
\begin{proof}
由定理 \ref{thm:xi-atomic-witt-cycle-primitive-exact-peeling} 直接在 \(u=1\) 代入即可。证毕。
\leanverified{paper\_xi\_atomic\_witt\_unweighted\_unique\_two\_step\_class}
\end{proof}

把同一原子剥离沿能量权变量 \(u\) 施以根单位 Fourier 反演，还可得到模 \(m\) 能量 residue 侧的完全局域化修正。

\leanverified{paper\_xi\_atomic\_witt\_energy\_residue\_l2\_law}
\begin{theorem}[原子 Witt 因子的能量 residue 单类注入与 \texorpdfstring{$L^2$}{L2} 原子律]\label{thm:xi-atomic-witt-energy-residue-l2-law}
固定整数 \(m\ge 2\)，记 \(\omega_m:=e^{2\pi i/m}\)。对 \(r\in \ZZ/m\ZZ\) 与 \(n\ge 1\)，定义 periodic 与 primitive 的能量 residue 计数
\[
A_{m,r}(n):=\#\{\gamma:\ |\gamma|=n,\ E(\gamma)\equiv r\pmod m\},
\]
\[
P_{m,r}(n):=\#\{\mathfrak p:\ |\mathfrak p|=n,\ E(\mathfrak p)\equiv r\pmod m\}.
\]
并以离散 Fourier 反演定义 core residue 坐标
\[
A_{m,r}^{\mathrm{core}}(n)
:=
\frac1m\sum_{j=0}^{m-1}a_n^{\mathrm{core}}(\omega_m^j)\,\omega_m^{-jr},
\]
\[
P_{m,r}^{\mathrm{core}}(n)
:=
\frac1m\sum_{j=0}^{m-1}p_n^{\mathrm{core}}(\omega_m^j)\,\omega_m^{-jr}.
\]
则对一切 \(n\ge 1\) 与 \(r\in \ZZ/m\ZZ\) 都有严格恒等式
\[
\boxed{\;
A_{m,r}(n)
=
A_{m,r}^{\mathrm{core}}(n)
+2\,\mathbf 1_{\{2\mid n\}}\mathbf 1_{\{r\equiv n/2\;(\mathrm{mod}\,m)\}}.\;}
\]
\[
\boxed{\;
P_{m,r}(n)
=
P_{m,r}^{\mathrm{core}}(n)
+\mathbf 1_{\{n=2\}}\mathbf 1_{\{r\equiv 1\;(\mathrm{mod}\,m)\}}.\;}
\]
并且 twisted 系数差满足精确平方平均律
\[
\boxed{\;
\frac1m\sum_{j=0}^{m-1}
\bigl|a_n(\omega_m^j)-a_n^{\mathrm{core}}(\omega_m^j)\bigr|^2
=
4\,\mathbf 1_{\{2\mid n\}}.\;}
\]
\[
\boxed{\;
\frac1m\sum_{j=0}^{m-1}
\bigl|p_n(\omega_m^j)-p_n^{\mathrm{core}}(\omega_m^j)\bigr|^2
=
\mathbf 1_{\{n=2\}}.\;}
\]
因此，长度 \(2\) 的原子因子在能量 residue 侧并不产生扩散型修正，而是对每个偶长度 \(n\) 只向单一 residue 类 \(r\equiv n/2\pmod m\) 注入恰好两个 periodic 原子；primitive 层则仅在 \((n,r)=(2,1)\) 处留下单点质量。
\end{theorem}
\begin{proof}
由定理 \ref{thm:xi-atomic-witt-cycle-primitive-exact-peeling}，对每个 \(0\le j\le m-1\) 都有
\[
a_n(\omega_m^j)-a_n^{\mathrm{core}}(\omega_m^j)
=
2\,\omega_m^{j n/2}\mathbf 1_{\{2\mid n\}},
\]
\[
p_n(\omega_m^j)-p_n^{\mathrm{core}}(\omega_m^j)
=
\omega_m^j\,\mathbf 1_{\{n=2\}}.
\]
因此
\[
\begin{aligned}
A_{m,r}(n)-A_{m,r}^{\mathrm{core}}(n)
&=
\frac1m\sum_{j=0}^{m-1}
\bigl(a_n(\omega_m^j)-a_n^{\mathrm{core}}(\omega_m^j)\bigr)\omega_m^{-jr} \\
&=
\frac{2\,\mathbf 1_{\{2\mid n\}}}{m}
\sum_{j=0}^{m-1}\omega_m^{j(n/2-r)}.
\end{aligned}
\]
单位根正交性给出
\[
\sum_{j=0}^{m-1}\omega_m^{j(n/2-r)}
=
\begin{cases}
m,& r\equiv n/2\pmod m,\\
0,& r\not\equiv n/2\pmod m,
\end{cases}
\]
从而得到 periodic residue 修正公式。

primitive 层同理：
\[
\begin{aligned}
P_{m,r}(n)-P_{m,r}^{\mathrm{core}}(n)
&=
\frac1m\sum_{j=0}^{m-1}
\bigl(p_n(\omega_m^j)-p_n^{\mathrm{core}}(\omega_m^j)\bigr)\omega_m^{-jr} \\
&=
\frac{\mathbf 1_{\{n=2\}}}{m}
\sum_{j=0}^{m-1}\omega_m^{j(1-r)},
\end{aligned}
\]
故仅当 \(n=2\) 且 \(r\equiv 1\pmod m\) 时取值为 \(1\)。

最后，由 \(|\omega_m^{j n/2}|=1\) 与 \(|\omega_m^j|=1\)，直接得到
\[
\frac1m\sum_{j=0}^{m-1}
\bigl|a_n(\omega_m^j)-a_n^{\mathrm{core}}(\omega_m^j)\bigr|^2
=
\frac1m\sum_{j=0}^{m-1}4\,\mathbf 1_{\{2\mid n\}}
=
4\,\mathbf 1_{\{2\mid n\}},
\]
\[
\frac1m\sum_{j=0}^{m-1}
\bigl|p_n(\omega_m^j)-p_n^{\mathrm{core}}(\omega_m^j)\bigr|^2
=
\frac1m\sum_{j=0}^{m-1}\mathbf 1_{\{n=2\}}
=
\mathbf 1_{\{n=2\}}.
\]
证毕。
\end{proof}

\begin{theorem}[长度同余类滤波中的单类注入律]\label{thm:xi-root-unity-filter-single-class-injection}
固定模数 \(q\ge 2\)，记 \(\omega_q:=e^{2\pi i/q}\)，并取剩余类 \(a\in\ZZ/q\ZZ\)。
定义 primitive 长度滤波
\[
P_{a,q}(z,u):=\frac1q\sum_{j=0}^{q-1}\omega_q^{-aj}P(\omega_q^j z,u),
\qquad
P_{a,q}^{\mathrm{core}}(z,u):=\frac1q\sum_{j=0}^{q-1}\omega_q^{-aj}P_{\mathrm{core}}(\omega_q^j z,u).
\]
则有严格恒等式
\[
\boxed{\;
P_{a,q}(z,u)=P_{a,q}^{\mathrm{core}}(z,u)+u z^2\,\mathbf 1_{\{a\equiv 2\;(\mathrm{mod}\,q)\}}.\;}
\]
进一步，定义周期对数滤波
\[
A_{a,q}(z,u):=\frac1q\sum_{j=0}^{q-1}\omega_q^{-aj}\log\zeta(\omega_q^j z,u),
\]
\[
A_{a,q}^{\mathrm{core}}(z,u):=\frac1q\sum_{j=0}^{q-1}\omega_q^{-aj}\log\zeta_{\mathrm{core}}(\omega_q^j z,u).
\]
则有闭式修正
\[
\boxed{\;
A_{a,q}(z,u)
=
A_{a,q}^{\mathrm{core}}(z,u)
+\sum_{\substack{k\ge 1\\ 2k\equiv a\;(\mathrm{mod}\,q)}}
\frac{u^k z^{2k}}{k}.\;}
\]
因此，原子长度 \(2\) 因子对 root-of-unity 长度筛的作用不是扩散到全部同余类，而是严格局域于单一剩余类 \(2\bmod q\)。
\end{theorem}
\leanverified{paper\_xi\_root\_unity\_filter\_single\_class\_injection}
\begin{proof}
由定理 \ref{thm:xi-atomic-witt-cycle-primitive-exact-peeling}，
\[
P(z,u)=P_{\mathrm{core}}(z,u)+u z^2.
\]
对其施加 root-of-unity 投影，得到
\[
\begin{aligned}
P_{a,q}(z,u)-P_{a,q}^{\mathrm{core}}(z,u)
&=
\frac1q\sum_{j=0}^{q-1}\omega_q^{-aj}u(\omega_q^j z)^2\\
&=
u z^2\cdot \frac1q\sum_{j=0}^{q-1}\omega_q^{j(2-a)}\\
&=
u z^2\,\mathbf 1_{\{a\equiv 2\;(\mathrm{mod}\,q)\}},
\end{aligned}
\]
从而得到第一式。这正是定理 \ref{thm:xi-atomic-witt-mod-class-localization} 的 Fourier 版本。

再由
\[
\log\zeta(z,u)
=
\log\zeta_{\mathrm{core}}(z,u)
-\log(1-uz^2)
=
\log\zeta_{\mathrm{core}}(z,u)
+\sum_{k\ge 1}\frac{u^k z^{2k}}{k},
\]
同样施加 root-of-unity 投影，便有
\[
\begin{aligned}
A_{a,q}(z,u)-A_{a,q}^{\mathrm{core}}(z,u)
&=
\frac1q\sum_{j=0}^{q-1}\omega_q^{-aj}
\sum_{k\ge 1}\frac{u^k(\omega_q^j z)^{2k}}{k}\\
&=
\sum_{k\ge 1}\frac{u^k z^{2k}}{k}
\left(
\frac1q\sum_{j=0}^{q-1}\omega_q^{j(2k-a)}
\right)\\
&=
\sum_{\substack{k\ge 1\\ 2k\equiv a\;(\mathrm{mod}\,q)}}\frac{u^k z^{2k}}{k}.
\end{aligned}
\]
证毕。
\end{proof}

\begin{theorem}[平方根门槛失败的 core 不可约性]\label{thm:xi-square-root-threshold-core-irreducibility}
沿用命题 \ref{prop:real-input-40-output-dirichlet-nonrh} 的记号。对任意 \(m\ge 2\) 与 \(1\le j\le m-1\)，令
\[
\omega_m:=e^{2\pi i/m},
\qquad
\rho_{m,j}:=\frac{1}{\min\{|z|:\ \Delta(z,\omega_m^j)=0\}}.
\]
若 \(\rho_{m,j}>1\)，则必有
\[
\boxed{\;
\rho_{m,j}
=
\frac{1}{\min\{|z|:\ F(z,\omega_m^j)=0\}}.\;}
\]
特别地，任一满足
\[
\theta_{m,j}:=\frac{\log\rho_{m,j}}{\log\lambda}>\frac12
\]
的非平凡扭曲模态，都完全由 core 因子 \(F(z,\omega_m^j)\) 决定；长度 \(2\) 原子因子 \((1-uz^2)\) 不可能制造这类超平方根门槛现象。
\end{theorem}
\begin{proof}
对 \(|u|=1\) 且 \(u=\omega_m^j\)，原子因子满足
\[
1-uz^2=0
\iff
z^2=u^{-1},
\]
故其全部零点都位于单位圆上，即 \(|z|=1\)。
另一方面，
\[
\Delta(z,\omega_m^j)=(1-\omega_m^j z^2)F(z,\omega_m^j).
\]
若 \(\rho_{m,j}>1\)，则按定义
\[
\min\{|z|:\ \Delta(z,\omega_m^j)=0\}<1.
\]
由于原子因子不可能在 \(|z|<1\) 上产生零点，故达到该最小模的零点只能来自 \(F(z,\omega_m^j)\)。
于是
\[
\min\{|z|:\ \Delta(z,\omega_m^j)=0\}
=
\min\{|z|:\ F(z,\omega_m^j)=0\},
\]
对两边取倒数即得结论。证毕。
\end{proof}

\begin{theorem}[Lee--Yang 三次域符号条件下的根存在事件坍缩]\label{thm:xi-leyang-sign-conditional-root-collapse}
令
\[
g(y):=256y^3+411y^2+165y+32,
\]
并记 \(P_{10}(u)\) 为推论 \ref{cor:fold-gauge-anomaly-P10-modp-splitting-densities} 中的十次分歧多项式。
对除去有限坏素后的素数 \(p\)，定义事件
\[
A:=\{\text{\(P_{10}\bmod p\) 含一次因子}\},
\]
\[
B_{\mathrm{split}}:=\{\text{\(g\bmod p\) 全分裂}\},
\qquad
B_{\mathrm{root}}:=\{\text{\(g\bmod p\) 至少有一根}\}.
\]
并定义 Lee--Yang 二次特征
\[
\chi_{\mathrm{LY}}(p):=\mathrm{sgn}\!\bigl(\sigma_3(p)\bigr)
=
\left(\frac{-111}{p}\right)\in\{\pm1\},
\]
其中 \(\sigma_3(p)\in S_3\) 为 \(g\) 的分裂域 Frobenius 元。
则条件密度满足
\[
\boxed{\;
\delta\!\bigl(B_{\mathrm{root}}\mid \chi_{\mathrm{LY}}=+1\bigr)=\frac13,
\qquad
\delta\!\bigl(B_{\mathrm{root}}\mid \chi_{\mathrm{LY}}=-1\bigr)=1,\;}
\]
\[
\boxed{\;
\delta\!\bigl(B_{\mathrm{split}}\mid \chi_{\mathrm{LY}}=+1\bigr)=\frac13,
\qquad
\delta\!\bigl(B_{\mathrm{split}}\mid \chi_{\mathrm{LY}}=-1\bigr)=0.\;}
\]
进一步，由 \(A\) 只依赖 \(S_{10}\) 分量，而
\(\chi_{\mathrm{LY}},B_{\mathrm{split}},B_{\mathrm{root}}\) 只依赖 \(S_3\) 分量，得到
\[
\boxed{\;
\delta\!\bigl(A\cap B_{\mathrm{root}}\mid \chi_{\mathrm{LY}}=+1\bigr)
=
\frac{28319}{134400},
\qquad
\delta\!\bigl(A\cap B_{\mathrm{root}}\mid \chi_{\mathrm{LY}}=-1\bigr)
=
\frac{28319}{44800},\;}
\]
\[
\boxed{\;
\delta\!\bigl(A\cap B_{\mathrm{split}}\mid \chi_{\mathrm{LY}}=+1\bigr)
=
\frac{28319}{134400},
\qquad
\delta\!\bigl(A\cap B_{\mathrm{split}}\mid \chi_{\mathrm{LY}}=-1\bigr)
=
0.\;}
\]
\end{theorem}
\begin{proof}
由结论 \ref{con:xi-terminal-zm-leyang-cubic-modp-splitting-density}，\(g\bmod p\) 的三种分解型与 \(S_3\) 的三类共轭类对应如下：
恒等类对应全分裂 \((1)(1)(1)\)，换位类对应 \((1)(2)\)，三循环类对应不可约 \((3)\)。
因此
\[
B_{\mathrm{split}} \Longleftrightarrow \sigma_3(p)\in \{e\},
\qquad
B_{\mathrm{root}} \Longleftrightarrow \sigma_3(p)\in \{e\}\cup \{\text{换位类}\}.
\]
另一方面，
\[
\chi_{\mathrm{LY}}(p)=+1
\Longleftrightarrow
\sigma_3(p)\in \{e\}\cup\{\text{三循环类}\},
\]
\[
\chi_{\mathrm{LY}}(p)=-1
\Longleftrightarrow
\sigma_3(p)\in \{\text{换位类}\}.
\]
在 \(S_3\) 中，这三类大小分别为 \(1,3,2\)。
故在 \(\chi_{\mathrm{LY}}=+1\) 的条件下，只剩下大小为 \(1\) 与 \(2\) 的两类，其中只有恒等类产生根，于是
\[
\delta\!\bigl(B_{\mathrm{root}}\mid \chi_{\mathrm{LY}}=+1\bigr)
=
\delta\!\bigl(B_{\mathrm{split}}\mid \chi_{\mathrm{LY}}=+1\bigr)
=
\frac{1}{1+2}
=
\frac13.
\]
在 \(\chi_{\mathrm{LY}}=-1\) 的条件下，只剩下换位类，故
\[
\delta\!\bigl(B_{\mathrm{root}}\mid \chi_{\mathrm{LY}}=-1\bigr)=1,
\qquad
\delta\!\bigl(B_{\mathrm{split}}\mid \chi_{\mathrm{LY}}=-1\bigr)=0.
\]

再由命题 \ref{prop:fold-gauge-anomaly-P10-leyang-linear-disjointness} 与推论
\ref{cor:fold-gauge-anomaly-P10-leyang-chebotarev-product-law}，
\(P_{10}\) 的分解型事件与 \(g\) 的 \(S_3\)-分解型事件严格独立；
并由推论 \ref{cor:fold-gauge-anomaly-P10-modp-splitting-densities}，
\[
\delta(A)=\frac{28319}{44800}.
\]
故
\[
\delta\!\bigl(A\cap B_{\mathrm{root}}\mid \chi_{\mathrm{LY}}=\pm1\bigr)
=
\delta(A)\,
\delta\!\bigl(B_{\mathrm{root}}\mid \chi_{\mathrm{LY}}=\pm1\bigr),
\]
\[
\delta\!\bigl(A\cap B_{\mathrm{split}}\mid \chi_{\mathrm{LY}}=\pm1\bigr)
=
\delta(A)\,
\delta\!\bigl(B_{\mathrm{split}}\mid \chi_{\mathrm{LY}}=\pm1\bigr).
\]
代入前述条件密度即得四条显式公式。证毕。
\end{proof}

\begin{corollary}[原子--core--算术三分律]\label{cor:xi-atomic-core-arithmetic-trichotomy}
在真实输入 \(40\) 状态核的一参数输出位势链与 \(P_{10}\times g\) 的算术筛选链上，可将可分辨结构规范分解为
\[
\boxed{\;
\text{原子层}\ \oplus\ \text{core 谱层}\ \oplus\ \text{算术筛选层}\;}
\]
更精确地：
\begin{enumerate}
  \item \textbf{原子层：}
        由定理 \ref{thm:xi-atomic-witt-cycle-primitive-exact-peeling}、
        \ref{thm:xi-atomic-witt-energy-residue-l2-law} 与
        \ref{thm:xi-root-unity-filter-single-class-injection}，
        原子因子 \((1-uz^2)^{-1}\) 在 primitive 层只占据长度 \(2\) 的单一坐标，
        在能量 residue 侧只向 \(r\equiv n/2\pmod m\) 注入单类 periodic 修正，
        并且在长度同余筛中只向剩余类 \(2\bmod q\) 注入局域残差；
  \item \textbf{core 谱层：}
        由定理 \ref{thm:xi-square-root-threshold-core-irreducibility}，
        一切满足 \(\theta_{m,j}>1/2\) 的非平凡扭曲模态都由 core 因子 \(F(z,\omega_m^j)\) 产生，
        因而平方根门槛的失效是纯粹的 core 谱现象；
  \item \textbf{算术筛选层：}
        由定理 \ref{thm:xi-root-unity-filter-single-class-injection} 与
        \ref{thm:xi-leyang-sign-conditional-root-collapse}，
        长度同余类的 root-of-unity 残差与
        \(\chi_{\mathrm{LY}}(p)=\left(\frac{-111}{p}\right)\) 的 Chebotarev 条件化共同构成一条显式筛选链，
        其模素端表现为“\(\chi_{\mathrm{LY}}=-1\) 时必有根、\(\chi_{\mathrm{LY}}=+1\) 时仅以密度 \(1/3\) 留根”的坍缩律。
\end{enumerate}
因此，长度 \(2\) 原子只负责能量 residue 与长度同余类上的局域残差；超越平方根门槛的全部指数复杂性由 core 谱承担；而把这些谱残差投影到算术根存在事件上的，则是由 \(\left(\frac{-111}{p}\right)\) 控制的 Chebotarev--Legendre 筛选。
\end{corollary}

\paragraph{冻结质量反演、容量矩的 Mellin--Stieltjes 精确式与 core 扭曲门槛}
沿着固定冻结相中的 escort/pressure 链、连续容量曲线的 Mellin--Stieltjes 对偶，
以及真实输入 \(40\) 状态核的原子--core 剥离链，还可进一步抽出下列七条互相咬合的接口。
前五条把冻结质量、压力斜率、离散凸结构、整数阈值刚性与容量矩写成可逆账本；后两条则说明非平凡 twisted 平方根门槛完全由 core 因子承担。

\begin{theorem}[冻结区中 escort 质量对最大纤维层的有限层指数集中]\label{thm:xi-time-part57b-escort-maxfiber-layer-concentration}
沿用定理 \ref{thm:xi-terminal-fixed-freezing-tv-collapse} 与定理
\ref{thm:fold-pressure-freezing-threshold} 的记号。对固定实数 \(a>0\)，定义 escort 测度
\[
\nu_{m,a}(x):=\frac{d_m(x)^a}{S_a(m)},
\qquad
S_a(m):=\sum_{x\in X_m}d_m(x)^a.
\]
再记
\[
M_m:=\max_{x\in X_m}d_m(x),\qquad
K_m:=\#\{x\in X_m:\ d_m(x)=M_m\},
\]
\[
M_m^{(2)}:=\max\{d_m(x):\ d_m(x)<M_m\},
\qquad
\mathcal M_m:=\{x\in X_m:\ d_m(x)=M_m\},
\]
其中若次大纤维不存在则约定 \(M_m^{(2)}:=1\)。则对每个 \(m\) 都有有限层界
\[
\boxed{\;
1-\nu_{m,a}(\mathcal M_m)
\le
\frac{|X_m|-K_m}{K_m}\left(\frac{M_m^{(2)}}{M_m}\right)^a
\le
\frac{|X_m|}{K_m}\left(\frac{M_m^{(2)}}{M_m}\right)^a.\;}
\]
若进一步满足
\[
a>a_{\mathrm c}:=\frac{\log\varphi-g_*}{\alpha_*-\alpha_*^{(2)}},
\]
则 escort 测度在 \(\mathcal M_m\) 上指数集中，并且
\[
\limsup_{m\to\infty}\frac1m
\log\bigl(1-\nu_{m,a}(\mathcal M_m)\bigr)
\le
\log\varphi+a\alpha_*^{(2)}-g_*-a\alpha_*<0.
\]
\end{theorem}
\begin{proof}
写
\[
S_a(m)=K_mM_m^a+R_a(m),
\qquad
R_a(m):=\sum_{d_m(x)<M_m}d_m(x)^a.
\]
于是
\[
1-\nu_{m,a}(\mathcal M_m)=\frac{R_a(m)}{S_a(m)}
\le
\frac{R_a(m)}{K_mM_m^a}.
\]
对补集 \(X_m\setminus\mathcal M_m\) 上的每个 \(x\)，都有 \(d_m(x)\le M_m^{(2)}\)，故
\[
R_a(m)\le (|X_m|-K_m)(M_m^{(2)})^a
\le |X_m|(M_m^{(2)})^a,
\]
从而得到两条有限层不等式。

若 \(a>a_{\mathrm c}\)，则由
\[
\frac1m\log|X_m|\to\log\varphi,\qquad
\frac1m\log K_m\to g_*,
\qquad
\frac1m\log M_m\to\alpha_*,
\]
\[
\limsup_{m\to\infty}\frac1m\log M_m^{(2)}=\alpha_*^{(2)},
\]
可得
\[
\limsup_{m\to\infty}\frac1m
\log\bigl(1-\nu_{m,a}(\mathcal M_m)\bigr)
\le
\log\varphi+a\alpha_*^{(2)}-g_*-a\alpha_*.
\]
由 \(a>a_{\mathrm c}\) 的定义，右端严格为负，故得到指数集中。证毕。
\end{proof}
\leanverified{paper\_xi\_time\_part57b\_escort\_maxfiber\_layer\_concentration}

\begin{theorem}[压力谱对 \texorpdfstring{$\alpha_*$}{alpha*} 的高阶斜率恢复与冻结截距恢复]\label{thm:xi-time-part57b-pressure-spectrum-slope-intercept-recovery}
\leanverified{paper\_xi\_time\_part57b\_pressure\_spectrum\_slope\_intercept\_recovery}
设压力函数
\[
P_a:=\lim_{m\to\infty}\frac1m\log S_a(m)
\qquad(a\in\ZZ_{\ge 1})
\]
存在，并满足定理 \ref{thm:fold-pressure-groundstate-bounds} 的双边估计
\[
a\alpha_*+g_*\le P_a\le a\alpha_*+\log\varphi.
\]
则：
\begin{enumerate}
  \item 斜率参数 \(\alpha_*\) 可由压力谱唯一恢复：
  \[
  \boxed{\;
  \alpha_*=\lim_{a\to\infty}\frac{P_a}{a}.\;}
  \]
  \item 若进一步存在某个整数 \(a_0\) 使冻结等式
  \[
  P_a=a\alpha_*+g_*
  \qquad(\forall a\ge a_0)
  \]
  成立，则对每个 \(a\ge a_0\) 都有
  \[
  \boxed{\;
  g_*=P_a-a\alpha_*.\;}
  \]
\end{enumerate}
换言之，高阶压力谱首先唯一锁定最大纤维指数 \(\alpha_*\)；而一旦进入冻结区，截距 \(g_*\) 亦随之由同一压力谱唯一锁定。
\end{theorem}
\begin{proof}
由
\[
a\alpha_*+g_*\le P_a\le a\alpha_*+\log\varphi
\]
两边同时除以 \(a\)，得到
\[
\alpha_*+\frac{g_*}{a}\le \frac{P_a}{a}\le \alpha_*+\frac{\log\varphi}{a}.
\]
令 \(a\to\infty\)，左右两端都收敛到 \(\alpha_*\)，故
\[
\lim_{a\to\infty}\frac{P_a}{a}=\alpha_*.
\]
这证明了第一条。

若对所有 \(a\ge a_0\) 都有冻结等式 \(P_a=a\alpha_*+g_*\)，则直接移项即得
\[
g_*=P_a-a\alpha_*.
\]
证毕。
\end{proof}

\begin{theorem}[纤维压力谱的离散凸性与基态超额的单调下降]\label{thm:xi-time-part57b-pressure-spectrum-discrete-convexity-monotone-excess}
\leanverified{paper\_xi\_time\_part57b\_pressure\_spectrum\_discrete\_convexity\_monotone\_excess}
沿用定理 \ref{thm:xi-time-part57b-pressure-spectrum-slope-intercept-recovery} 的记号。对每个整数 \(a\ge 1\)，记
\[
S_a(m):=\sum_{x\in X_m}d_m(x)^a,
\qquad
P_a:=\lim_{m\to\infty}\frac1m\log S_a(m),
\]
并定义基态超额
\[
E_a:=P_a-a\alpha_*.
\]
则成立：
\begin{enumerate}
  \item 对每个整数 \(a\ge 1\)，压力谱满足离散凸性
  \[
  \boxed{\;2P_{a+1}\le P_a+P_{a+2}.\;}
  \]
  \item 序列 \((E_a)_{a\ge 1}\) 单调不增，即
  \[
  \boxed{\;E_{a+1}\le E_a\qquad(a\ge 1).\;}
  \]
  \item 对每个整数 \(a\ge 1\)，都有
  \[
  \boxed{\;g_*\le E_a\le \log\varphi.\;}
  \]
\end{enumerate}
\end{theorem}
\begin{proof}
对固定 \(m\)，记
\[
\mathbf d^{(m)}:=\bigl(d_m(x)\bigr)_{x\in X_m}.
\]
由幂和的对数凸性（等价于 Hölder 不等式），对每个整数 \(a\ge 1\) 有
\[
S_{a+1}(m)^2\le S_a(m)\,S_{a+2}(m).
\]
两边取对数、除以 \(m\)，再令 \(m\to\infty\)，便得到
\[
2P_{a+1}\le P_a+P_{a+2}.
\]

另一方面，由 \(d_m(x)\le M_m\) 对一切 \(x\in X_m\) 成立，
\[
S_{a+1}(m)=\sum_{x\in X_m}d_m(x)^{a+1}
\le
M_m\sum_{x\in X_m}d_m(x)^a
=
M_m\,S_a(m).
\]
故
\[
\frac1m\log S_{a+1}(m)-\frac{a+1}{m}\log M_m
\le
\frac1m\log S_a(m)-\frac{a}{m}\log M_m.
\]
令 \(m\to\infty\)，利用
\[
\lim_{m\to\infty}\frac1m\log M_m=\alpha_*,
\]
便得到
\[
E_{a+1}\le E_a.
\]

最后，由定理 \ref{thm:fold-pressure-groundstate-bounds}，
\[
a\alpha_*+g_*\le P_a\le a\alpha_*+\log\varphi.
\]
两边同时减去 \(a\alpha_*\)，即得
\[
g_*\le E_a\le \log\varphi.
\]
证毕。
\end{proof}

\begin{corollary}[冻结相整数起始点的单阈值刚性]\label{cor:xi-time-part57b-freezing-first-contact-threshold}
\leanverified{paper\_xi\_time\_part57b\_freezing\_first\_contact\_threshold}
若严格基态间隙 \(\alpha_*^{(2)}<\alpha_*\) 成立，并记
\[
a_{\mathrm c}:=\frac{\log\varphi-g_*}{\alpha_*-\alpha_*^{(2)}}<\infty,
\qquad
a_{\mathrm{fr}}:=\min\{a\in\ZZ_{\ge 1}:\ E_a=g_*\},
\]
则 \(a_{\mathrm{fr}}\) 存在，且满足
\[
\boxed{\;a_{\mathrm{fr}}\le \lfloor a_{\mathrm c}\rfloor+1.\;}
\]
并且对每个整数 \(a\ge a_{\mathrm{fr}}\)，都有
\[
\boxed{\;E_a=g_*.\;}
\]
换言之，冻结相的整数起始点是“第一次触底后永不抬升”的单阈值事件。
\end{corollary}
\begin{proof}
由定理 \ref{thm:fold-pressure-freezing-threshold}，对每个实数 \(a>a_{\mathrm c}\) 都有
\[
P_a=a\alpha_*+g_*,
\]
故对整数
\[
a_0:=\lfloor a_{\mathrm c}\rfloor+1
\]
必有 \(a_0>a_{\mathrm c}\)，从而
\[
E_{a_0}=P_{a_0}-a_0\alpha_*=g_*.
\]
因此集合 \(\{a\in\ZZ_{\ge 1}:E_a=g_*\}\) 非空，\(a_{\mathrm{fr}}\) 良定义，且
\[
a_{\mathrm{fr}}\le a_0=\lfloor a_{\mathrm c}\rfloor+1.
\]

另一方面，由定理 \ref{thm:xi-time-part57b-pressure-spectrum-discrete-convexity-monotone-excess}，
\[
E_{a+1}\le E_a
\qquad\text{且}\qquad
E_a\ge g_*.
\]
于是，一旦某个整数 \(a_1\) 满足 \(E_{a_1}=g_*\)，便对所有 \(a\ge a_1\) 都只能有
\[
g_*\le E_a\le E_{a_1}=g_*,
\]
故
\[
E_a=g_*
\qquad(a\ge a_1).
\]
取 \(a_1=a_{\mathrm{fr}}\) 即得结论。证毕。
\end{proof}

\begin{corollary}[冻结相整数矩阶尾部的精确仿射半群律]\label{cor:xi-time-part57b-pressure-affine-semigroup-law}
\leanverified{paper\_xi\_time\_part57b\_pressure\_affine\_semigroup\_law}
沿用推论 \ref{cor:xi-time-part57b-freezing-first-contact-threshold} 的记号。若整数
\[
a,b>a_{\mathrm c},
\]
则有严格恒等式
\[
\boxed{\;P_a=a\alpha_*+g_*.\;}
\]
并且
\[
\boxed{\;P_{a+1}-P_a=\alpha_*.\;}
\]
以及
\[
\boxed{\;P_{a+b}=P_a+P_b-g_*.\;}
\]
换言之，一旦进入冻结相，整数矩阶上的压力谱尾部满足精确的仿射半群律。
\end{corollary}
\begin{proof}
由定理 \ref{thm:fold-pressure-freezing-threshold}，对任意实数 \(t>a_{\mathrm c}\) 都有
\[
P_t=t\alpha_*+g_*.
\]
特别地，对任意整数 \(a,b>a_{\mathrm c}\)，立得
\[
P_a=a\alpha_*+g_*,
\qquad
P_b=b\alpha_*+g_*,
\qquad
P_{a+b}=(a+b)\alpha_*+g_*.
\]
于是
\[
P_{a+1}-P_a
=
\bigl((a+1)\alpha_*+g_*\bigr)-\bigl(a\alpha_*+g_*\bigr)
=
\alpha_*,
\]
并且
\[
P_{a+b}
=
(a+b)\alpha_*+g_*
=
\bigl(a\alpha_*+g_*\bigr)+\bigl(b\alpha_*+g_*\bigr)-g_*
=
P_a+P_b-g_*.
\]
证毕。
\end{proof}

\begin{theorem}[连续容量曲线与全部幂和矩之间的精确 Mellin--Stieltjes 反演]\label{thm:xi-time-part57b-capacity-moment-exact-mellin-stieltjes}
沿用定义 \ref{def:fold-capacity-tail}、命题 \ref{prop:fold-capacity-tail-equivalence} 与命题
\ref{prop:fold-moment-mellin-stieltjes} 的记号。对固定分辨率 \(m\)，记
\[
N_m(T):=\#\{x\in X_m:\ d_m(x)\ge T\},
\qquad
\mathcal C_m^{\mathrm{cont}}(T):=\sum_{x\in X_m}\min(d_m(x),T).
\]
则对任意实数 \(a>0\) 都有精确公式
\[
\boxed{\;
S_a(m):=\sum_{x\in X_m}d_m(x)^a
=
a\int_0^\infty T^{a-1}N_m(T)\,\dd T.\;}
\]
并且对任意 \(a>1\) 还有
\[
\boxed{\;
S_a(m)
=
a(a-1)\int_0^\infty
T^{a-2}\bigl(2^m-\mathcal C_m^{\mathrm{cont}}(T)\bigr)\,\dd T.\;}
\]
因此，连续容量曲线 \(\mathcal C_m^{\mathrm{cont}}\) 不仅控制可逆容量门槛，而且已经携带了全部实阶幂和矩 \(\{S_a(m)\}_{a>1}\) 的信息。
\end{theorem}
\begin{proof}
对任意整数 \(d\ge 1\) 与任意 \(a>0\)，都有单点恒等式
\[
d^a=a\int_0^d T^{a-1}\,\dd T.
\]
对 \(d=d_m(x)\) 在 \(x\in X_m\) 上求和，并交换有限和与积分，便得到
\[
S_a(m)=a\int_0^\infty T^{a-1}\#\{x\in X_m:\ d_m(x)\ge T\}\,\dd T
=
a\int_0^\infty T^{a-1}N_m(T)\,\dd T.
\]
这正是命题 \ref{prop:fold-moment-mellin-stieltjes} 的实阶写法。

再对 \(a>1\)，有另一单点恒等式
\[
d^a=a(a-1)\int_0^\infty T^{a-2}(d-T)_+\,\dd T.
\]
对 \(x\in X_m\) 求和得
\[
S_a(m)=
a(a-1)\int_0^\infty
T^{a-2}\sum_{x\in X_m}(d_m(x)-T)_+\,\dd T.
\]
而
\[
\sum_{x\in X_m}(d_m(x)-T)_+
=
\sum_{x\in X_m}\bigl(d_m(x)-\min(d_m(x),T)\bigr)
\]
\[
=
\sum_{x\in X_m}d_m(x)-\mathcal C_m^{\mathrm{cont}}(T)
=
2^m-\mathcal C_m^{\mathrm{cont}}(T),
\]
其中最后一步利用
\(
\sum_{x\in X_m}d_m(x)=2^m
\)。
代回即得第二式。证毕。
\end{proof}

\begin{theorem}[原子 Euler 因子不改变任何非平凡 twisted 指数]\label{thm:xi-time-part57b-atomic-factor-zero-twisted-exponent}
沿用定理 \ref{thm:xi-atomic-witt-cycle-primitive-exact-peeling} 的记号。对任意 \(m\ge 2\) 与 \(1\le j\le m-1\)，令
\[
u=\omega_m^j,\qquad
\rho_{m,j}^{\mathrm{full}}
:=
\frac{1}{\min\{|z|:\ \Delta(z,u)=0\}},
\qquad
\rho_{m,j}^{\mathrm{core}}
:=
\frac{1}{\min\{|z|:\ F(z,u)=0\}}.
\]
若记 \(\lambda:=\rho(M(1))=\varphi^2\)，并定义
\[
\theta_{m,j}^{\mathrm{full}}
:=
\frac{\log \rho_{m,j}^{\mathrm{full}}}{\log\lambda},
\qquad
\theta_{m,j}^{\mathrm{core}}
:=
\frac{\log \rho_{m,j}^{\mathrm{core}}}{\log\lambda},
\]
则成立精确关系
\[
\boxed{\;
\rho_{m,j}^{\mathrm{full}}=\max\{1,\rho_{m,j}^{\mathrm{core}}\},\;}
\]
\[
\boxed{\;
\theta_{m,j}^{\mathrm{full}}=\max\{0,\theta_{m,j}^{\mathrm{core}}\}.\;}
\]
特别地，原子 Euler 因子 \((1-uz^2)^{-1}\) 至多把 twisted 指数抬到 \(\theta=0\)；任何真正的正指数 twisted 增长都只能来自 core 因子 \(F(z,u)\)。
\end{theorem}
\begin{proof}
由分解
\[
\Delta(z,u)=(1-uz^2)\,F(z,u),
\]
full 分母的零点由两部分组成。原子因子 \(1-uz^2\) 的零点满足
\[
z^2=u^{-1},
\]
而 \(|u|=1\)，故其全部零点都位于单位圆上，即其贡献的最小模恰为 \(1\)。
另一方面，core 因子 \(F(z,u)\) 的最小零点模长按定义为
\[
\frac{1}{\rho_{m,j}^{\mathrm{core}}}.
\]
因此 full 分母全部零点中的最小模满足
\[
\min\{|z|:\ \Delta(z,u)=0\}
=
\min\left\{1,\frac{1}{\rho_{m,j}^{\mathrm{core}}}\right\}.
\]
取倒数即得
\[
\rho_{m,j}^{\mathrm{full}}=\max\{1,\rho_{m,j}^{\mathrm{core}}\}.
\]
再除以 \(\log\lambda\) 取对数，便得到
\[
\theta_{m,j}^{\mathrm{full}}=\max\{0,\theta_{m,j}^{\mathrm{core}}\}.
\]
证毕。
\end{proof}

\begin{corollary}[twisted 平方根门槛在 core 与 full 之间完全等价]\label{cor:xi-time-part57b-core-full-square-root-barrier-equivalence}
沿用定理 \ref{thm:xi-time-part57b-atomic-factor-zero-twisted-exponent} 的记号，并定义
\[
\theta_m^{\mathrm{full}}
:=
\max_{1\le j\le m-1}\theta_{m,j}^{\mathrm{full}},
\qquad
\theta_m^{\mathrm{core}}
:=
\max_{1\le j\le m-1}\theta_{m,j}^{\mathrm{core}}.
\]
则有等价关系
\[
\boxed{\;
\theta_m^{\mathrm{full}}>\frac12
\quad\Longleftrightarrow\quad
\theta_m^{\mathrm{core}}>\frac12.\;}
\]
因此，文稿中 Dirichlet/primitive 扭曲误差不具备 RH 形平方根指数的现象，在剥离原子因子 \((1-uz^2)^{-1}\) 后完全保留；相反，将该门槛失效归因于原子 Euler 因子的解释并不成立。
\end{corollary}
\begin{proof}
由定理 \ref{thm:xi-time-part57b-atomic-factor-zero-twisted-exponent}，
\[
\theta_{m,j}^{\mathrm{full}}=\max\{0,\theta_{m,j}^{\mathrm{core}}\}
\qquad(1\le j\le m-1).
\]
由于阈值 \(1/2\) 严格大于 \(0\)，故对每个 \(j\) 都有
\[
\theta_{m,j}^{\mathrm{full}}>\frac12
\quad\Longleftrightarrow\quad
\theta_{m,j}^{\mathrm{core}}>\frac12.
\]
对 \(j\) 取最大值即得结论。证毕。
\end{proof}

\begin{remark}[有限层审计接口]\label{rem:xi-time-part57b-freezing-capacity-core-twist-audit}
脚本 \texttt{scripts/exp\_xi\_freezing\_capacity\_core\_twist\_audit.py}
对三类有限层数据作逐项核验：其一是最大纤维层上的 escort 质量与有限层集中上界，
其二是尾函数/连续容量曲线与实阶矩的 Mellin--Stieltjes 两种积分恒等式，
其三是真实输入 \(40\) 状态核在若干模数下的 full/core twisted 半径与指数一致性。
\end{remark}

\endinput

\paragraph{冻结微正则分解与原子 residue 的单维扰动}
在冻结 escort 态向最大纤维层的指数集中、以及单原子 Euler 因子在 periodic/core residue 侧的精确剥离都已建立之后，还可把这两条链进一步压缩为一组完全闭式的终端结论：冻结相中的 escort 态相对于基态微正则分布具有精确二分解，而原子 residue 修正在有限 Fourier 侧则形成单维等幅相位向量。

\begin{conclusion}[冻结相 escort 态相对于基态微正则分布的精确二分解与 \texorpdfstring{$f$}{f}-散度闭式]\label{con:xi-time-part57ba-freezing-escort-groundstate-fdiv-closed}
沿用定理 \ref{thm:xi-fixed-freezing-escort-renyi-spectrum-collapse} 与推论 \ref{cor:xi-fixed-freezing-escort-groundstate-tv-identity} 的记号，固定 \(a>a_{\mathrm c}\)。记
\[
\sigma_m:=\mathrm{Unif}(A_m^{\max})
=\frac{1}{K_m}\ind_{A_m^{\max}}(\cdot),
\qquad
p_{m,a}:=\pi_{m,a}(A_m^{\max}).
\]
则存在概率测度 \(\eta_{m,a}\)；当 \(p_{m,a}<1\) 时其支撑于 \(X_m\setminus A_m^{\max}\)，并且
\[
\pi_{m,a}=p_{m,a}\sigma_m+(1-p_{m,a})\eta_{m,a},
\qquad
\pi_{m,a}(\,\cdot\mid A_m^{\max})=\sigma_m.
\]
进一步，对任意满足 \(f(1)=0\) 且 \(f(0)<\infty\) 的 Csisz\'ar--Morimoto 散度
\[
D_f(\sigma_m\|\pi_{m,a})
:=
\sum_{x\in X_m}\pi_{m,a}(x)\,
f\!\left(\frac{\sigma_m(x)}{\pi_{m,a}(x)}\right),
\]
都有精确闭式
\[
\boxed{\;
D_f(\sigma_m\|\pi_{m,a})
=
p_{m,a}\,f\!\left(\frac{1}{p_{m,a}}\right)
+(1-p_{m,a})f(0).\;}
\]
特别地，
\[
\boxed{\;
\|\pi_{m,a}-\sigma_m\|_{\mathrm{TV}}=1-p_{m,a},\;}
\]
\[
\boxed{\;
D_{\mathrm{KL}}(\sigma_m\|\pi_{m,a})=-\log p_{m,a},\;}
\]
\[
\boxed{\;
\chi^2(\sigma_m\|\pi_{m,a})=\frac{1}{p_{m,a}}-1.\;}
\]
并且对任意固定的 \(f\)，上述散度满足
\[
D_f(\sigma_m\|\pi_{m,a})
=
O\!\bigl(\exp(-m\Delta(a)+o(m))\bigr).
\]
\end{conclusion}
\begin{proof}
由于 \(d_m(x)=M_m\) 在 \(A_m^{\max}\) 上恒定，对任意 \(x\in A_m^{\max}\) 都有
\[
\pi_{m,a}(x)=\frac{M_m^a}{S_a(m)}=\frac{p_{m,a}}{K_m}.
\]
因此 \(\pi_{m,a}\) 在 \(A_m^{\max}\) 上的条件分布恰为 \(\sigma_m\)。若 \(p_{m,a}<1\)，取 \(\eta_{m,a}\) 为 \(\pi_{m,a}\) 在补集 \(X_m\setminus A_m^{\max}\) 上的条件分布；若 \(p_{m,a}=1\)，由于第二项系数为零，任取 \(X_m\) 上概率测度 \(\eta_{m,a}\) 即可。于是得到所述精确分解。

再计算 \(f\)-散度。对 \(x\in A_m^{\max}\)，有
\[
\frac{\sigma_m(x)}{\pi_{m,a}(x)}
=
\frac{1/K_m}{p_{m,a}/K_m}
=
\frac{1}{p_{m,a}},
\]
而对 \(x\notin A_m^{\max}\)，有 \(\sigma_m(x)=0\)。故
\[
\begin{aligned}
D_f(\sigma_m\|\pi_{m,a})
&=
\sum_{x\in A_m^{\max}}\pi_{m,a}(x)\,
f\!\left(\frac{1}{p_{m,a}}\right)
+
\sum_{x\notin A_m^{\max}}\pi_{m,a}(x)\,f(0)\\
&=
p_{m,a}\,f\!\left(\frac{1}{p_{m,a}}\right)
+(1-p_{m,a})f(0).
\end{aligned}
\]
这给出一般闭式。

对 \(f(t)=t\log t\)，得到
\[
D_{\mathrm{KL}}(\sigma_m\|\pi_{m,a})
=
p_{m,a}\cdot \frac{1}{p_{m,a}}\log\frac{1}{p_{m,a}}
=
-\log p_{m,a}.
\]
对 \(f(t)=(t-1)^2\)，得到
\[
\chi^2(\sigma_m\|\pi_{m,a})
=
p_{m,a}\left(\frac{1}{p_{m,a}}-1\right)^2+(1-p_{m,a})
=
\frac{1}{p_{m,a}}-1.
\]
总变差公式则由支撑分解直接给出：
\[
\|\pi_{m,a}-\sigma_m\|_{\mathrm{TV}}
=
\frac12\sum_{x\in A_m^{\max}}\frac{1-p_{m,a}}{K_m}
+\frac12\sum_{x\notin A_m^{\max}}\pi_{m,a}(x)
=
1-p_{m,a}.
\]

最后，由定理 \ref{thm:xi-fixed-freezing-escort-renyi-spectrum-collapse} 的第四条，
\[
1-p_{m,a}\le \exp\!\bigl(-m\Delta(a)+o(m)\bigr).
\]
当 \(m\) 充分大时有 \(1/p_{m,a}\in[1,2]\)。由于固定的凸函数 \(f\) 在紧区间 \([1,2]\) 上 Lipschitz，故
\[
f\!\left(\frac{1}{p_{m,a}}\right)=O(1-p_{m,a}).
\]
代回一般闭式即得
\[
D_f(\sigma_m\|\pi_{m,a})
=
O(1-p_{m,a})
=
O\!\bigl(\exp(-m\Delta(a)+o(m))\bigr).
\]
证毕。
\end{proof}

\begin{conclusion}[冻结相 \texorpdfstring{$\min$}{min}-entropy 的精确有限尺度修正]\label{con:xi-time-part57ba-freezing-minentropy-finite-size}
沿用结论 \ref{con:xi-time-part57ba-freezing-escort-groundstate-fdiv-closed} 的记号。则 escort 分布的 \(\min\)-entropy 满足精确恒等式
\[
\boxed{\;
H_\infty(\pi_{m,a})
:=
-\log\max_{x\in X_m}\pi_{m,a}(x)
=
\log K_m-\log p_{m,a}.\;}
\]
从而
\[
\boxed{\;
H_\infty(\pi_{m,a})
=
\log K_m
+O\!\bigl(\exp(-m\Delta(a)+o(m))\bigr).\;}
\]
特别地，若 \(g_*=\lim_{m\to\infty}m^{-1}\log K_m\)，则
\[
\boxed{\;
H_\infty(\pi_{m,a})-g_*m
=
\bigl(\log K_m-g_*m\bigr)
+O\!\bigl(\exp(-m\Delta(a)+o(m))\bigr).\;}
\]
\end{conclusion}
\begin{proof}
由上一结论，对 \(x\in A_m^{\max}\) 有
\[
\pi_{m,a}(x)=\frac{p_{m,a}}{K_m},
\]
而对 \(x\notin A_m^{\max}\) 有 \(\pi_{m,a}(x)\le p_{m,a}/K_m\)，因为这些点的权重对应严格小于 \(M_m\) 的纤维。故最大概率恰在 \(A_m^{\max}\) 上达到，并满足
\[
\max_{x\in X_m}\pi_{m,a}(x)=\frac{p_{m,a}}{K_m}.
\]
取负对数即得第一式。

再由
\[
1-p_{m,a}\le \exp\!\bigl(-m\Delta(a)+o(m)\bigr)
\]
与标准估计
\[
-\log p_{m,a}=O(1-p_{m,a}),
\]
便得到
\[
H_\infty(\pi_{m,a})
=
\log K_m+O\!\bigl(\exp(-m\Delta(a)+o(m))\bigr).
\]
最后一式只是把主项 \(g_*m\) 从 \(\log K_m\) 中分离出来。证毕。
\end{proof}

\begin{conclusion}[原子 residue 修正的 delta--flat 极值结构]\label{con:xi-time-part57ba-atomic-residue-delta-flat-extremizer}
沿用定理 \ref{thm:xi-atomic-witt-cycle-primitive-exact-peeling} 与定理 \ref{thm:xi-atomic-witt-energy-residue-l2-law} 的记号。固定偶整数 \(n\) 与模数 \(m\ge 2\)，定义原子 residue 修正
\[
A_{m,r}^{\mathrm{cycle}}(n)
:=
A_{m,r}^{\mathrm{full}}(n)-A_{m,r}^{\mathrm{core}}(n),
\]
以及其 twisted Fourier 向量
\[
a_n^{\mathrm{cycle}}(\omega_m^j)
:=
a_n(\omega_m^j)-a_n^{\mathrm{core}}(\omega_m^j).
\]
则有显式公式
\[
\boxed{\;
A_{m,r}^{\mathrm{cycle}}(n)
=
2\,\ind_{\{r\equiv n/2\;(\mathrm{mod}\,m)\}},\;}
\]
\[
\boxed{\;
a_n^{\mathrm{cycle}}(\omega_m^j)=2\,\omega_m^{j n/2}\qquad(j\in\ZZ/m\ZZ).\;}
\]
因此：
\begin{enumerate}
  \item residue 侧支撑恰为单点，即
  \[
  \boxed{\;\#\supp A_{m,\bullet}^{\mathrm{cycle}}(n)=1.\;}
  \]
  \item twisted Fourier 侧全部模长相同，即
  \[
  \boxed{\;\abs{a_n^{\mathrm{cycle}}(\omega_m^j)}=2\qquad(\forall j\in\ZZ/m\ZZ).\;}
  \]
  \item Parseval 能量满足
  \[
  \boxed{\;
  \sum_{r\in\ZZ/m\ZZ}\abs{A_{m,r}^{\mathrm{cycle}}(n)}^2=4,\qquad
  \frac1m\sum_{j=0}^{m-1}\abs{a_n^{\mathrm{cycle}}(\omega_m^j)}^2=4.\;}
  \]
  \item 对任意 \(1\le p\le \infty\)，都有
  \[
  \boxed{\;\|A_{m,\bullet}^{\mathrm{cycle}}(n)\|_{\ell^p}=2,\;}
  \]
  以及
  \[
  \boxed{\;
  \|a_{m,\bullet}^{\mathrm{cycle}}(n)\|_{\ell^p}
  =
  2m^{1/p}\quad(1\le p<\infty),\qquad
  \|a_{m,\bullet}^{\mathrm{cycle}}(n)\|_{\ell^\infty}=2,\;}
  \]
  其中
  \(
  a_{m,\bullet}^{\mathrm{cycle}}(n):=
  (a_n^{\mathrm{cycle}}(\omega_m^j))_{j\in\ZZ/m\ZZ}.
  \)
\end{enumerate}
换言之，原子 residue 修正在 residue 侧呈现单点支撑，而在 twisted Fourier 侧呈现等幅相位向量；二者由有限循环群上的精确 Fourier 对偶所连接。
\end{conclusion}
\begin{proof}
上述两条显式公式正是定理 \ref{thm:xi-atomic-witt-energy-residue-l2-law} 与定理 \ref{thm:xi-atomic-witt-cycle-primitive-exact-peeling} 的直接重述。由 \(A_{m,r}^{\mathrm{cycle}}(n)\) 只在唯一的剩余类 \(r\equiv n/2\pmod m\) 上取值 \(2\)，立得单点支撑与
\[
\|A_{m,\bullet}^{\mathrm{cycle}}(n)\|_{\ell^p}=2
\qquad(1\le p\le \infty).
\]
另一方面，对每个 \(j\) 都有 \(\abs{\omega_m^{j n/2}}=1\)，故
\[
\abs{a_n^{\mathrm{cycle}}(\omega_m^j)}=2
\qquad(\forall j\in\ZZ/m\ZZ),
\]
从而
\[
\frac1m\sum_{j=0}^{m-1}\abs{a_n^{\mathrm{cycle}}(\omega_m^j)}^2
=
\frac1m\sum_{j=0}^{m-1}4
=
4
\]
以及
\[
\|a_{m,\bullet}^{\mathrm{cycle}}(n)\|_{\ell^p}
=
\left(\sum_{j=0}^{m-1}2^p\right)^{1/p}
=
2m^{1/p}
\qquad(1\le p<\infty),
\]
\[
\|a_{m,\bullet}^{\mathrm{cycle}}(n)\|_{\ell^\infty}=2.
\]
最后，residue 侧平方和只有唯一非零坐标贡献，故
\[
\sum_{r\in\ZZ/m\ZZ}\abs{A_{m,r}^{\mathrm{cycle}}(n)}^2=4.
\]
证毕。
\end{proof}

\begin{conclusion}[full/core residue 差的单维秩一扰动与统一投影界]\label{con:xi-time-part57ba-core-full-residue-rankone-perturbation}
沿用结论 \ref{con:xi-time-part57ba-atomic-residue-delta-flat-extremizer} 的记号，并定义差向量
\[
\Delta A_{m,r}(n)
:=
A_{m,r}^{\mathrm{full}}(n)-A_{m,r}^{\mathrm{core}}(n)
=
A_{m,r}^{\mathrm{cycle}}(n).
\]
若记 \(r_0\equiv n/2\pmod m\)，则 residue 侧有
\[
\boxed{\;
\Delta A_{m,\bullet}(n)=2e_{r_0},\;}
\]
其中 \(e_{r_0}\) 为 \(\CC^{\ZZ/m\ZZ}\) 的标准基向量；而 twisted Fourier 侧则有
\[
\boxed{\;
\bigl(a_n(\omega_m^j)-a_n^{\mathrm{core}}(\omega_m^j)\bigr)_{j\in\ZZ/m\ZZ}
=
2(\omega_m^{j n/2})_{j\in\ZZ/m\ZZ}.\;}
\]
因此 full/core residue 差始终局域在一维坐标子空间中。特别地，对任意满足
\[
Q=Q^*=Q^2,
\qquad
Q\mathbf 1=0
\]
的 Hermitian 投影 \(Q\)，都有统一界
\[
\boxed{\;
\|Q\Delta A_{m,\bullet}(n)\|_2\le 2.\;}
\]
故任意由 residue 向量经过固定正交投影所得到的去均值统计量，其 full/core 差都至多表现为一个单维扰动，并且该扰动的 \(\ell^2\) 幅值与模数 \(m\) 无关。
\end{conclusion}
\begin{proof}
由上一结论，
\[
\Delta A_{m,r}(n)=2\,\ind_{\{r=r_0\}},
\]
故 \(\Delta A_{m,\bullet}(n)=2e_{r_0}\)，这说明 residue 差确实落在一维坐标子空间中。其 twisted Fourier 图像同样由上一结论的显式公式直接给出：
\[
a_n(\omega_m^j)-a_n^{\mathrm{core}}(\omega_m^j)=2\omega_m^{j n/2}.
\]

再由 \(Q=Q^*=Q^2\) 可知 \(Q\) 是 \(\ell^2\) 上的正交投影，从而是压缩算子：
\[
\|Qv\|_2\le \|v\|_2
\qquad(\forall v\in\CC^{\ZZ/m\ZZ}).
\]
取 \(v=\Delta A_{m,\bullet}(n)\)，并利用
\[
\|\Delta A_{m,\bullet}(n)\|_2=\|2e_{r_0}\|_2=2,
\]
即可得到
\[
\|Q\Delta A_{m,\bullet}(n)\|_2\le 2.
\]
证毕。
\end{proof}

\endinput

\paragraph{冻结 escort 的基态等权塌缩与实阶 Rényi 熵率刚性}
在冻结相的总变差恒等式、基态微正则二分解与实阶 Rényi 压力差公式都已建立之后，还可把基态层上的等权塌缩与全部信息型熵率的刚性收束为同一条直接接口。

\begin{corollary}[冻结 escort 的基态等权塌缩与实阶 Rényi 熵率刚性]\label{cor:xi-time-part57bb-freezing-groundstate-equidistribution-renyi-rigidity}
沿用定理 \ref{thm:xi-fixed-freezing-escort-renyi-spectrum-collapse}、推论 \ref{cor:xi-fixed-freezing-escort-groundstate-tv-identity} 与定理 \ref{thm:xi-fixed-freezing-escort-bounded-observable-collapse} 的记号。设严格基态间隙
\[
\alpha_*^{(2)}<\alpha_*,
\]
并固定
\[
a>a_{\mathrm c}.
\]
记
\[
u_m^{\max}:=\mathrm{Unif}(A_m^{\max}),
\qquad
\Delta(a):=a(\alpha_*-\alpha_*^{(2)})-(\log\varphi-g_*)>0.
\]
则成立：
\begin{enumerate}
  \item 全变差满足精确恒等式
  \[
  \|\pi_{m,a}-u_m^{\max}\|_{\mathrm{TV}}
  =
  1-\pi_{m,a}(A_m^{\max})
  \le
  \exp\!\bigl(-m\Delta(a)+o(m)\bigr).
  \]
  因而对任意有界函数 \(f_m:X_m\to\mathbb R\)，都有
  \[
  \left|
  \int f_m\,\dd\pi_{m,a}
  -
  \frac1{K_m}\sum_{x\in A_m^{\max}}f_m(x)
  \right|
  \le
  2\|f_m\|_\infty
  \exp\!\bigl(-m\Delta(a)+o(m)\bigr).
  \]
  \item 对任意实数 \(s\in(1,\infty)\)，都有
  \[
  \lim_{m\to\infty}\frac1m H_s(\pi_{m,a})=g_*.
  \]
  若进一步取 \(s=\infty\)，则同样有
  \[
  \lim_{m\to\infty}\frac1m H_\infty(\pi_{m,a})=g_*.
  \]
\end{enumerate}
换言之，冻结区内的 escort 态在分布层面指数塌缩到最大纤维层上的等权微正则态，而在信息层面则把全部实阶 Rényi 熵率统一压缩到同一基态复杂度 \(g_*\)。
\end{corollary}
\begin{proof}
第一条中的总变差恒等式与指数界正是推论 \ref{cor:xi-fixed-freezing-escort-groundstate-tv-identity} 的内容；对有界测试函数的估计则由定理 \ref{thm:xi-fixed-freezing-escort-bounded-observable-collapse} 立即得到。

对第二条，任取 \(s\in(1,\infty)\)。由于 \(a>a_{\mathrm c}\) 且 \(s>1\)，自动有
\[
as>a>a_{\mathrm c}.
\]
于是定理 \ref{thm:xi-fixed-freezing-escort-renyi-spectrum-collapse} 的第二条给出
\[
\lim_{m\to\infty}\frac1m H_s(\pi_{m,a})=g_*.
\]
而 \(s=\infty\) 的情形正是同一定理第三条的内容。证毕。
\end{proof}

\begin{corollary}[冻结 escort 在最大纤维子集上的微正则计数公式]\label{cor:xi-time-part57bb-freezing-groundstate-subset-equidistribution}
沿用推论 \ref{cor:xi-time-part57bb-freezing-groundstate-equidistribution-renyi-rigidity} 的记号。对任意子集
\[
B_m\subseteq A_m^{\max},
\]
都有精确公式
\[
\pi_{m,a}(B_m)
=
\pi_{m,a}(A_m^{\max})\frac{|B_m|}{K_m}.
\]
从而
\[
\left|
\pi_{m,a}(B_m)-\frac{|B_m|}{K_m}
\right|
=
\bigl(1-\pi_{m,a}(A_m^{\max})\bigr)\frac{|B_m|}{K_m}
\le
\exp\!\bigl(-m\Delta(a)+o(m)\bigr).
\]
因此，冻结区中最大纤维内的任意子事件都以指数速度微正则化。
\end{corollary}
\begin{proof}
由推论 \ref{cor:xi-time-part57bb-freezing-groundstate-equidistribution-renyi-rigidity} 的第一条，\(\pi_{m,a}\) 在 \(A_m^{\max}\) 上的条件分布恰为均匀分布 \(u_m^{\max}\)。故对任意 \(x\in A_m^{\max}\) 都有
\[
\pi_{m,a}(x)=\frac{\pi_{m,a}(A_m^{\max})}{K_m}.
\]
于是
\[
\pi_{m,a}(B_m)
=
\sum_{x\in B_m}\pi_{m,a}(x)
=
\pi_{m,a}(A_m^{\max})\frac{|B_m|}{K_m}.
\]
两边减去 \(|B_m|/K_m\)，便得到
\[
\left|
\pi_{m,a}(B_m)-\frac{|B_m|}{K_m}
\right|
=
\bigl(1-\pi_{m,a}(A_m^{\max})\bigr)\frac{|B_m|}{K_m}.
\]
最后再用同一推论中的指数估计
\[
1-\pi_{m,a}(A_m^{\max})\le \exp\!\bigl(-m\Delta(a)+o(m)\bigr)
\]
即得结论。证毕。
\leanverified{paper\_xi\_time\_part57bb\_freezing\_groundstate\_subset\_equidistribution}
\end{proof}

\endinput

\paragraph{冻结相的均匀简并化、局部精确反演率与双温度分离}
在冻结 escort 对最大纤维层的指数集中、整数矩阶尾部的精确仿射律以及全微态精确反演阈值已经建立之后，还可把冻结端面的配分函数、局部精确编码率与全局位预算阈值进一步压缩为同一条接口。

\begin{theorem}[冻结相的均匀简并化]\label{thm:xi-time-part57bc-freezing-uniform-degeneracy}
沿用定理 \ref{thm:xi-fixed-freezing-escort-renyi-spectrum-collapse} 与推论 \ref{cor:xi-time-part57bb-freezing-groundstate-equidistribution-renyi-rigidity} 的记号。设
\[
\alpha_*^{(2)}<\alpha_*,
\qquad
a>a_{\mathrm c}:=\frac{\log\varphi-g_*}{\alpha_*-\alpha_*^{(2)}},
\]
并记
\[
A_m^{\max}:=\{x\in X_m:\ d_m(x)=M_m\},
\qquad
K_m:=|A_m^{\max}|,
\qquad
u_m^{\max}:=\mathrm{Unif}(A_m^{\max}),
\]
\[
\pi_{m,a}(x):=\frac{d_m(x)^a}{S_a(m)},
\qquad
\Delta(a):=a(\alpha_*-\alpha_*^{(2)})-(\log\varphi-g_*)>0.
\]
则有
\[
S_a(m)=K_mM_m^a\Bigl(1+O\!\bigl(e^{-m\Delta(a)+o(m)}\bigr)\Bigr),
\]
并且
\[
\|\pi_{m,a}-u_m^{\max}\|_{\mathrm{TV}}
=
O\!\bigl(e^{-m\Delta(a)+o(m)}\bigr).
\]
换言之，冻结区内的 \(a\)-阶配分函数已经彻底由最大纤维层支配，而 escort 态则以同一指数尺度塌缩到该层上的均匀分布。
\end{theorem}
\begin{proof}
记
\[
R_a(m):=\sum_{d_m(x)<M_m}d_m(x)^a.
\]
则
\[
S_a(m)=K_mM_m^a+R_a(m).
\]
另一方面，
\[
\pi_{m,a}(A_m^{\max})
=
\frac{K_mM_m^a}{S_a(m)}
=
\frac{1}{1+\varepsilon_m},
\qquad
\varepsilon_m:=\frac{R_a(m)}{K_mM_m^a}.
\]
由定理 \ref{thm:xi-time-part57b-escort-maxfiber-layer-concentration} 以及推论 \ref{cor:xi-time-part57bb-freezing-groundstate-equidistribution-renyi-rigidity}，
\[
1-\pi_{m,a}(A_m^{\max})
=
\frac{\varepsilon_m}{1+\varepsilon_m}
\le
\exp\!\bigl(-m\Delta(a)+o(m)\bigr).
\]
于是对充分大的 \(m\)，右端小于 \(1/2\)，从而
\[
\varepsilon_m
\le
2\exp\!\bigl(-m\Delta(a)+o(m)\bigr)
=
O\!\bigl(e^{-m\Delta(a)+o(m)}\bigr).
\]
代回
\[
S_a(m)=K_mM_m^a(1+\varepsilon_m)
\]
便得第一式。

第二式则由同一推论的总变差恒等式
\[
\|\pi_{m,a}-u_m^{\max}\|_{\mathrm{TV}}
=
1-\pi_{m,a}(A_m^{\max})
\]
与上式立得。证毕。
\end{proof}

\begin{theorem}[冻结支撑上的精确局部反演率]\label{thm:xi-time-part57bc-frozen-support-exact-local-inversion-rate}
\leanverified{paper\_xi\_time\_part57bc\_frozen\_support\_exact\_local\_inversion\_rate}
沿用定理 \ref{thm:xi-full-microstate-exact-inversion-bitrate-threshold} 的记号，并定义最大纤维支撑
\[
\Omega_m^{\max}
:=
\bigcup_{x\in A_m^{\max}}\Fold_m^{-1}(x).
\]
记 \(B_m^{\max}\) 为下述最小比特数：给定读出 \(x\in A_m^{\max}\) 后，在 \(\Omega_m^{\max}\) 上实现精确微态反演所需的最小侧信息预算。则对每个 \(m\) 都有严格恒等式
\[
B_m^{\max}
=
\left\lceil \log_2 M_m\right\rceil.
\]
特别地，
\[
\lim_{m\to\infty}\frac{B_m^{\max}}{m}
=
\frac{\alpha_*}{\log 2}.
\]
因此，在冻结支撑上，局部精确可逆率恰由最大纤维指数 \(\alpha_*\) 单独锁定。
\end{theorem}
\begin{proof}
固定 \(x\in A_m^{\max}\)。由定义，
\[
\bigl|\Fold_m^{-1}(x)\bigr|=M_m.
\]
若只允许 \(B\) 比特侧信息，则在该纤维上至多能区分 \(2^B\) 个微态。要对整条最大纤维精确反演，必要条件为
\[
2^B\ge M_m,
\]
故
\[
B_m^{\max}\ge \left\lceil \log_2 M_m\right\rceil.
\]

反之，对每个 \(x\in A_m^{\max}\) 任选双射
\[
\iota_x:\Fold_m^{-1}(x)\xrightarrow{\sim}\{0,1,\dots,M_m-1\}.
\]
把 \(\iota_x(\omega)\) 编成长度 \(\lceil\log_2M_m\rceil\) 的二进制字并作为侧信息，即可在固定 \(x\) 的条件下精确恢复任意
\(\omega\in\Fold_m^{-1}(x)\)。
故下界可达，从而
\[
B_m^{\max}
=
\left\lceil \log_2 M_m\right\rceil.
\]
最后再由定理 \ref{thm:xi-full-microstate-exact-inversion-bitrate-threshold} 中
\[
\lim_{m\to\infty}\frac{1}{m}\log M_m=\alpha_*
\]
即可得到
\[
\lim_{m\to\infty}\frac{B_m^{\max}}{m}
=
\frac{\alpha_*}{\log 2}.
\]
证毕。
\end{proof}

\begin{corollary}[两种信息温度的严格分离]\label{cor:xi-time-part57bc-two-information-temperature-separation}
定义
\[
\rho_{\mathrm{global}}:=\log_2\frac{2}{\varphi},
\qquad
\rho_{\mathrm{local}}:=\frac{\alpha_*}{\log 2}.
\]
则：
\begin{enumerate}
  \item 在均匀微态输入的全局反演问题中，若一列比特预算 \(B_m\) 满足
  \[
  \mathrm{Succ}_m^{\mathrm{bin}}(B_m)\ge 1-\varepsilon
  \qquad
  (0<\varepsilon<1\ \text{固定}),
  \]
  则必有
  \[
  \liminf_{m\to\infty}\frac{B_m}{m}\ge \rho_{\mathrm{global}}.
  \]
  \item 在冻结相渐近全部质量所支撑的最大纤维族 \(\Omega_m^{\max}\) 上，精确局部反演的最小比特率恰为
  \[
  \lim_{m\to\infty}\frac{B_m^{\max}}{m}
  =
  \rho_{\mathrm{local}}.
  \]
\end{enumerate}
因此，该体系同时携带一个外层像空间亏损率 \(\rho_{\mathrm{global}}\) 与一个内层纤维定位率 \(\rho_{\mathrm{local}}\)；二者源于完全不同的机制，不能混同。
\end{corollary}
\begin{proof}
第一条正是定理 \ref{thm:xi-time-part68-entropy-deficit-triple-identification} 的后半部分。第二条则是定理 \ref{thm:xi-time-part57bc-frozen-support-exact-local-inversion-rate} 的直接重述。证毕。
\end{proof}

\begin{corollary}[冻结斜率即编码斜率]\label{cor:xi-time-part57bc-freezing-slope-coding-rate}
记
\[
\beta_{\max}
:=
\lim_{m\to\infty}\frac{B_m^{\max}}{m}
=
\frac{\alpha_*}{\log 2}.
\]
则对任意整数 \(a>a_{\mathrm c}\)，都有
\[
P_a
=
g_*+a(\log 2)\beta_{\max}.
\]
特别地，冻结相中的高阶自由能直线，其斜率恰等于最大纤维精确反演率的自然对数版本。
\end{corollary}
\begin{proof}
由定理 \ref{thm:xi-time-part57bc-frozen-support-exact-local-inversion-rate}，
\[
(\log 2)\beta_{\max}=\alpha_*.
\]
再由推论 \ref{cor:xi-time-part57b-pressure-affine-semigroup-law}，
\[
P_a=a\alpha_*+g_*
\qquad
(a>a_{\mathrm c},\ a\in\ZZ),
\]
代入即得
\[
P_a=g_*+a(\log 2)\beta_{\max}.
\]
证毕。
\end{proof}

\endinput

\paragraph{碰撞势的局部对偶、素数模阈值与 dyadic 三谱迹}
碰撞位势 \(P_2(\theta)=\log\rho(e^\theta)\) 已在附录中被压缩为一条三次代数压力链：它一方面给出 \(\theta=0\) 处全部低阶累积量的 \(\sqrt5\)-闭式，另一方面在单位根扭曲与 \(u=-1\) 奇偶通道上保留足够低阶的谱数据，使局部 Legendre 对偶、介观模数阈值与 dyadic 周期迹振荡可以继续闭合为一组不再依赖外加参数化的终端结论。

\begin{theorem}[碰撞速率函数在平衡点的四阶 Cramér--Edgeworth 局部几何]\label{thm:xi-real-input-40-collision-rate-cramer-edgeworth}
沿用命题 \ref{prop:real-input-40-collision-pressure}、推论 \ref{cor:real-input-40-collision-cumulants}、\ref{cor:real-input-40-collision-cumulants-higher} 与注 \ref{rem:real-input-40-collision-ldp-param} 的记号。按注 \ref{rem:real-input-40-collision-ldp-param} 的归一化，定义
\[
I_2(\alpha):=\sup_{\theta\in\RR}\bigl(\alpha\theta-(P_2(\theta)-P_2(0))\bigr),
\qquad
\alpha_*:=P_2'(0)=\frac{3-\sqrt5}{10}.
\]
再记
\[
\kappa_2:=P_2''(0)=\frac{6\sqrt5-5}{125},\qquad
\kappa_3:=P_2^{(3)}(0)=\frac{9+10\sqrt5}{625},\qquad
\kappa_4:=P_2^{(4)}(0)=\frac{7+24\sqrt5}{3125}.
\]
则在 \(\alpha_*\) 的某个邻域内，若写
\[
\delta:=\alpha-\alpha_*,
\]
便有解析展开
\[
\boxed{\;
I_2(\alpha_*+\delta)
=
\frac{\delta^2}{2\kappa_2}
-\frac{\kappa_3}{6\kappa_2^3}\,\delta^3
+\frac{3\kappa_3^2-\kappa_2\kappa_4}{24\kappa_2^5}\,\delta^4
+O(\delta^5).\;}
\]
并且前三个非平凡局部系数可完全显式写成
\[
\boxed{\;
\frac{1}{2\kappa_2}
=
\frac{125}{62}+\frac{75\sqrt5}{31}
\approx 7.4259709133059428\;}
\]
\[
\boxed{\;
-\frac{\kappa_3}{6\kappa_2^3}
=
-\frac{849375}{59582}-\frac{525250\sqrt5}{89373}
\approx -27.3970573347852767\;}
\]
\[
\boxed{\;
\frac{3\kappa_3^2-\kappa_2\kappa_4}{24\kappa_2^5}
=
\frac{32390871875}{343549812}
+\frac{4803058125\sqrt5}{114516604}
\approx 188.068114201582058.\;}
\]
特别地，由于 \(\kappa_3>0\)，局部奇次项严格不对称：
\[
\boxed{\;
I_2(\alpha_*+\delta)-I_2(\alpha_*-\delta)
=
-\frac{\kappa_3}{3\kappa_2^3}\,\delta^3+O(\delta^5)
<0
\qquad(\delta>0\ \text{充分小}).\;}
\]
因此在均值点附近，右侧偏差严格比左侧偏差“更便宜”。
\end{theorem}
\begin{proof}
由推论 \ref{cor:real-input-40-collision-cumulants} 与推论 \ref{cor:real-input-40-collision-cumulants-higher}，
\[
P_2(\theta)-P_2(0)
=
\alpha_*\theta+\frac{\kappa_2}{2}\theta^2+\frac{\kappa_3}{6}\theta^3+\frac{\kappa_4}{24}\theta^4+O(\theta^5),
\]
故
\[
\alpha(\theta):=P_2'(\theta)
=
\alpha_*+\kappa_2\theta+\frac{\kappa_3}{2}\theta^2+\frac{\kappa_4}{6}\theta^3+O(\theta^4).
\]
由 \(\kappa_2>0\) 可知 \(\theta=0\) 处存在局部反函数 \(\theta=\theta(\delta)\)。逐次比较幂次，得到
\[
\theta(\delta)
=
\frac{\delta}{\kappa_2}
-\frac{\kappa_3}{2\kappa_2^3}\,\delta^2
+\frac{3\kappa_3^2-\kappa_2\kappa_4}{6\kappa_2^5}\,\delta^3
+O(\delta^4).
\]
另一方面，归一化 Legendre 对偶满足 \(I_2(\alpha_*)=0\) 且
\[
I_2'(\alpha)=\theta(\alpha)
\]
在 \(\alpha_*\) 邻域内成立，因此
\[
I_2(\alpha_*+\delta)=\int_0^\delta \theta(s)\,\dd s.
\]
逐项积分便得盒装展开式。再把 \(\kappa_2,\kappa_3,\kappa_4\) 的闭式代入并化简，即得三条显式系数公式。最后以 \(\delta\) 与 \(-\delta\) 相减并保留奇次项，即得第二条盒装式；由于 \(\kappa_3=(9+10\sqrt5)/625>0\)，右侧严格小于左侧。脚本 \texttt{scripts/exp\_xi\_real\_input\_40\_collision\_ldp\_median\_twist\_audit.py} 对反函数系数与速率函数系数给出符号核验。证毕。
\end{proof}

\begin{corollary}[碰撞计数的 Edgeworth 连续中位点具有显式常数位移]\label{cor:xi-real-input-40-collision-edgeworth-continuous-median-shift}
沿用定理 \ref{thm:xi-real-input-40-collision-rate-cramer-edgeworth} 的记号。令 \(S_n\) 为长度 \(n\) 上碰撞观测 \(\varphi_2\) 的总和，并把 \(\mathfrak m_n\) 定义为下列一阶 Edgeworth 半质量方程在 \(n\alpha_*\) 附近的唯一实解：
\[
\Phi\!\left(\frac{\mathfrak m_n-n\alpha_*}{\sqrt{n\kappa_2}}\right)
+\frac{\kappa_3}{6\kappa_2^{3/2}\sqrt n}
\left(
1-\left(\frac{\mathfrak m_n-n\alpha_*}{\sqrt{n\kappa_2}}\right)^2
\right)
\phi\!\left(\frac{\mathfrak m_n-n\alpha_*}{\sqrt{n\kappa_2}}\right)
=
\frac12,
\]
其中 \(\Phi,\phi\) 分别为标准正态分布函数与密度。则
\[
\boxed{\;
\mathfrak m_n
=
n\alpha_*-\frac{\kappa_3}{6\kappa_2}+o(1)
=
n\frac{3-\sqrt5}{10}
-\left(\frac{23}{310}+\frac{52\sqrt5}{2325}\right)+o(1).\;}
\]
特别地，连续中心位置始终位于 \(n\alpha_*\) 的左侧一个非零常数量级处。另一方面，由于 \(\alpha_*=(3-\sqrt5)/10\notin\QQ\)，任何整数中位数序列都不可能满足
\[
m_n=n\alpha_*+C+o(1)
\]
对某个固定常数 \(C\) 成立；因此上式刻画的是 Edgeworth 中心位置本身，而离散中位数的精确描述必须再附加独立的 lattice 修正。
\end{corollary}
\begin{proof}
对于有限状态解析加权矩阵的一维加性泛函，附录 \ref{app:pressure-analytic} 的标准解析扰动法逐字适用；结合推论 \ref{cor:real-input-40-collision-cumulants} 与推论 \ref{cor:real-input-40-collision-cumulants-higher}，碰撞总数的标准化变量
\[
Z_n:=\frac{S_n-n\alpha_*}{\sqrt{n\kappa_2}}
\]
满足一阶 Edgeworth 近似
\[
\PP(Z_n\le x)
=
\Phi(x)
+\frac{\kappa_3}{6\kappa_2^{3/2}\sqrt n}(1-x^2)\phi(x)
+O(n^{-1})
\]
在中心窗口内一致成立。令
\[
x_n:=\frac{\mathfrak m_n-n\alpha_*}{\sqrt{n\kappa_2}}.
\]
由中心极限定理知 \(x_n\to0\)。将 \(x=x_n\) 代入上式，并用
\[
\Phi(x_n)-\frac12=\phi(0)x_n+O(x_n^3),
\qquad
(1-x_n^2)\phi(x_n)=\phi(0)+O(x_n^2),
\]
得到
\[
0=\phi(0)x_n+\frac{\kappa_3}{6\kappa_2^{3/2}\sqrt n}\phi(0)+O(n^{-1}),
\]
从而
\[
x_n=-\frac{\kappa_3}{6\kappa_2^{3/2}\sqrt n}+O(n^{-1}).
\]
乘回 \(\sqrt{n\kappa_2}\) 即得
\[
\mathfrak m_n-n\alpha_*=-\frac{\kappa_3}{6\kappa_2}+o(1).
\]
最后，将 \(\kappa_2,\kappa_3\) 的闭式直接代入并化简，便得到
\[
-\frac{\kappa_3}{6\kappa_2}
=
-\frac{23}{310}-\frac{52\sqrt5}{2325}.
\]
又因 \(\alpha_*\notin\QQ\)，序列 \(n\alpha_*\) 的小数部分不可能收敛，故整数值中位数不可能收敛到某个固定常数平移的形式；若要对其作精确刻画，必须额外引入 lattice 修正。证毕。
\end{proof}

\begin{theorem}[模素数扭曲的精确介观窗口：\texorpdfstring{$p^2$}{p2} 尺度阈值律]\label{thm:xi-real-input-40-collision-prime-twist-p2-threshold}
沿用命题 \ref{prop:real-input-40-collision-pressure} 与推论 \ref{cor:real-input-40-collision-twist-fourth} 的记号。对每个奇素数 \(p\)，记
\[
\rho_p:=\rho\!\bigl(W(e^{2\pi i/p})\bigr),
\qquad
\lambda:=\rho(W(1))=\varphi^2,
\]
并定义常数
\[
c_2:=\frac{2\pi^2}{125}(6\sqrt5-5),
\qquad
c_4:=\frac{2\pi^4}{9375}(7+24\sqrt5).
\]
则当 \(p\to\infty\) 时有精确渐近
\[
\boxed{\;
\log\frac{\rho_p}{\lambda}
=
-\frac{c_2}{p^2}+\frac{c_4}{p^4}+O(p^{-6}).\;}
\]
因此，对任意沿奇素数指标取的正整数列 \(\{n_p\}\)：
\begin{enumerate}
  \item 若 \(n_p/p^2\to 0\)，则
  \[
  \left(\frac{\rho_p}{\lambda}\right)^{n_p}\to 1;
  \]
  \item 若 \(n_p/p^2\to \tau\in(0,\infty)\)，则
  \[
  \left(\frac{\rho_p}{\lambda}\right)^{n_p}\to e^{-c_2\tau};
  \]
  \item 若 \(n_p/p^2\to\infty\)，则
  \[
  \left(\frac{\rho_p}{\lambda}\right)^{n_p}\to 0.
  \]
\end{enumerate}
换言之，模 \(p\) 的非平凡扭曲通道在长度尺度 \(n\asymp p^2\) 上发生唯一的介观过渡：低于该尺度时几乎不可见；高于该尺度时则被指数压制。
\end{theorem}
\begin{proof}
推论 \ref{cor:real-input-40-collision-twist-fourth} 在 \(t_p=2\pi/p\) 处给出
\[
\log\frac{\rho_p}{\lambda}
=
-\frac{P_2''(0)}{2}\left(\frac{2\pi}{p}\right)^2
+\frac{P_2^{(4)}(0)}{24}\left(\frac{2\pi}{p}\right)^4
+O(p^{-6}),
\]
将推论 \ref{cor:real-input-40-collision-cumulants} 与推论 \ref{cor:real-input-40-collision-cumulants-higher} 中的闭式代入，即得第一条盒装式。等价地，
\[
\log\frac{\rho_p}{\lambda}
=
-\frac{c_2}{p^2}\left(1-\frac{c_4}{c_2p^2}+O(p^{-4})\right)
=
-\frac{c_2+o(1)}{p^2}.
\]
于是
\[
n_p\log\frac{\rho_p}{\lambda}
=
-\bigl(c_2+o(1)\bigr)\frac{n_p}{p^2},
\]
三种极限情形立刻分别给出 \(0\)、\(-c_2\tau\) 与 \(-\infty\)。指数化即得所述结论。
\par
由于周期迹层的模 \(p\) 扭曲通道正是 \(\Tr(W(e^{2\pi i/p})^n)\)，而 primitive 层只是在其上再施加 Möbius/Witt 投影，上述阈值便同时给出两层通道的唯一介观尺度。脚本 \texttt{scripts/exp\_xi\_real\_input\_40\_collision\_ldp\_median\_twist\_audit.py} 对若干奇素数样本的 \(p^{-6}\) 余项进行数值核验。证毕。
\end{proof}

\begin{theorem}[模 \texorpdfstring{$2$}{2} 扭曲周期迹的精确三谱分解]\label{thm:xi-real-input-40-collision-mod2-trace-three-spectrum}
沿用命题 \ref{prop:real-input-40-collision-pressure} 与注 \ref{rem:real-input-40-collision-parity} 的记号。令
\[
F=\begin{pmatrix}1&1\\ 1&0\end{pmatrix},
\qquad
A:=F\otimes F,
\qquad
W(-1):=D(-1)A,
\]
并定义
\[
a_n(-1):=\Tr\!\bigl(W(-1)^n\bigr)\qquad(n\ge 1).
\]
则 \(\det(xI-W(-1))\) 精确分解为
\[
(x+1)\bigl(x^3-2x^2-1\bigr).
\]
记 \(\rho_-\) 为三次方程
\[
x^3-2x^2-1=0
\]
的最大实根，并记其余两根为 \(\xi,\bar\xi\)。则成立：
\begin{enumerate}
  \item
  \[
  |\xi|=|\bar\xi|=\rho_-^{-1/2},
  \qquad
  \xi+\bar\xi=2-\rho_-;
  \]
  \item 存在唯一角 \(\vartheta\in(0,\pi)\) 使
  \[
  \xi=\rho_-^{-1/2}e^{i\vartheta},
  \qquad
  \cos\vartheta=\frac{(2-\rho_-)\sqrt{\rho_-}}{2};
  \]
  \item 对一切 \(n\ge 1\)，有完全闭式
  \[
  \boxed{\;
  a_n(-1)
  =
  \rho_-^n+(-1)^n+2\rho_-^{-n/2}\cos(n\vartheta).\;}
  \]
\end{enumerate}
特别地，
\[
a_n(-1)=\rho_-^n+O(1),
\qquad
\frac1n\log|a_n(-1)|\to \log\rho_-.
\]
因此模 \(2\) 奇偶扭曲下不存在额外隐藏模态：全部振荡都由单一实根 \(\rho_-\) 与单一角频率 \(\vartheta\) 完全控制。
\end{theorem}
\begin{proof}
命题 \ref{prop:real-input-40-collision-pressure} 的特征多项式公式在 \(u=-1\) 处退化为
\[
\det(xI-W(-1))=(x+1)(x^3-2x^2-1).
\]
因此 \(W(-1)\) 的四个特征值恰为
\[
-1,\qquad \rho_-,\qquad \xi,\qquad \bar\xi.
\]
对 monic 三次多项式 \(x^3-2x^2+0x-1\) 应用 Vi\`ete 关系，得到
\[
\rho_-+\xi+\bar\xi=2,
\qquad
\rho_-\xi\bar\xi=1.
\]
由于 \(\xi,\bar\xi\) 为共轭复根，第二式给出
\[
|\xi|^2=\xi\bar\xi=\rho_-^{-1},
\]
从而 \(|\xi|=|\bar\xi|=\rho_-^{-1/2}\)；第一式则给出
\[
2\Re\xi=2-\rho_-.
\]
于是存在唯一 \(\vartheta\in(0,\pi)\) 使
\[
\xi=\rho_-^{-1/2}e^{i\vartheta},
\qquad
\cos\vartheta=\frac{\Re\xi}{|\xi|}
=
\frac{(2-\rho_-)\sqrt{\rho_-}}{2}.
\]
最后，迹等于全部特征值的 \(n\) 次幂之和，故
\[
a_n(-1)
=
(-1)^n+\rho_-^n+\xi^n+\bar\xi^n
=
\rho_-^n+(-1)^n+2\rho_-^{-n/2}\cos(n\vartheta).
\]
将右式除以 \(\rho_-^n\) 得
\[
\rho_-^{-n}a_n(-1)
=
1+(-1)^n\rho_-^{-n}+2\rho_-^{-3n/2}\cos(n\vartheta),
\]
故 \(\rho_-^{-n}a_n(-1)\to1\)，于是既有 \(a_n(-1)=\rho_-^n+O(1)\)，又有
\[
\frac1n\log|a_n(-1)|\to\log\rho_-.
\]
脚本 \texttt{scripts/exp\_xi\_real\_input\_40\_collision\_ldp\_median\_twist\_audit.py} 对特征多项式分解以及前若干 \(n\) 的迹样本进行逐项核验。证毕。
\end{proof}

\begin{theorem}[模 \texorpdfstring{$2$}{2} 奇偶扭曲对应的负频率代数支]\label{thm:xi-real-input-40-collision-mod2-negative-frequency-branch}
沿用定理 \ref{thm:xi-real-input-40-collision-mod2-trace-three-spectrum} 的记号，并令
\[
\alpha_-:=
\frac{(\rho_-^2-2\rho_--1)(\rho_--1)}
{2\rho_-^3-5\rho_-^2+4\rho_-+1}.
\]
则 \(\alpha_-\) 是三次代数数，并满足不可约方程
\[
\boxed{\;
59\alpha^3-59\alpha^2+15\alpha+2=0.\;}
\]
该三次多项式仅有一个实根，且
\[
\alpha_-\approx -0.0947099154877496.
\]
特别地，
\[
\alpha_-\notin \left[0,\frac12\right].
\]
因此，模 \(2\) 奇偶扭曲不落在定理 \ref{thm:xi-time-part9p-positive-branch-global-coordinate} 的实正扭曲分支上，而对应于同一代数参数化中的唯一负频率实支。
\end{theorem}
\begin{proof}
由定理 \ref{thm:xi-real-input-40-collision-mod2-trace-three-spectrum}，
\(\rho_-\) 满足
\[
\rho_-^3-2\rho_-^2-1=0.
\]
另一方面，定理 \ref{thm:xi-time-part9p-collision-frequency-algebraic-ldp} 已给出参数化公式
\[
\alpha(\rho)=
\frac{(\rho^2-2\rho-1)(\rho-1)}
{2\rho^3-5\rho^2+4\rho+1}.
\]
因此 \(\alpha_-=\alpha(\rho_-)\) 满足二元代数系统
\[
\rho^3-2\rho^2-1=0,
\]
\[
\alpha(2\rho^3-5\rho^2+4\rho+1)-(\rho^2-2\rho-1)(\rho-1)=0.
\]
对 \(\rho\) 取 Sylvester 结果式可得
\[
2\bigl(59\alpha^3-59\alpha^2+15\alpha+2\bigr)=0,
\]
故 \(\alpha_-\) 必满足所示三次方程。由有理根判别法，其可能有理根只能取
\[
\pm 1,\ \pm 2,\ \pm \frac{1}{59},\ \pm \frac{2}{59},
\]
直接代入均不为零，因此该三次在 \(\QQ\) 上不可约。

记
\[
f(\alpha):=59\alpha^3-59\alpha^2+15\alpha+2.
\]
其判别式为
\[
\Delta_f=-625931<0,
\]
故 \(f\) 恰有一个实根。又
\[
f(0)=2>0,
\qquad
f\!\left(-\frac{1}{10}\right)=-\frac{149}{1000}<0,
\]
故该唯一实根位于 \((-\tfrac{1}{10},0)\)。由于 \(\alpha_-\) 本身就是该三次的实根，只能等于这一唯一实根，从而
\[
\alpha_-\approx -0.0947099154877496.
\]
最后，由定理 \ref{thm:xi-time-part9p-positive-branch-global-coordinate}，实正扭曲支与 \(\alpha\in(0,\tfrac12)\) 严格一一对应；既然 \(\alpha_-<0\)，该模 \(2\) 分支便不可能属于实正扭曲支。上述结式消元、残差与数值根可由脚本 \texttt{scripts/exp\_xi\_real\_input\_40\_collision\_ldp\_median\_twist\_audit.py} 逐项核验。证毕。
\end{proof}

\begin{theorem}[奇素数单位根扭曲的近主模态刚性]\label{thm:xi-real-input-40-collision-prime-near-perron-rigidity}
沿用定理 \ref{thm:xi-real-input-40-collision-prime-twist-p2-threshold} 的记号。对奇素数 \(p\) 定义
\[
\Theta_p:=\frac{\log \rho_p}{\log\lambda}.
\]
则当 \(p\to\infty\) 时有四阶渐近
\[
\boxed{\;
\Theta_p
=
1-\frac{6\sqrt5-5}{250\log\lambda}\Bigl(\frac{2\pi}{p}\Bigr)^2
+\frac{7+24\sqrt5}{75000\log\lambda}\Bigl(\frac{2\pi}{p}\Bigr)^4
+O(p^{-6}).\;}
\]
特别地，
\[
\Theta_p
=
1-\frac{1.3809571881649568\ldots}{p^2}
+\frac{1.3098894548658713\ldots}{p^4}
+O(p^{-6}),
\]
从而 \(\Theta_p\to 1\)。
\end{theorem}
\begin{proof}
由定理 \ref{thm:xi-real-input-40-collision-prime-twist-p2-threshold}，
\[
\log\frac{\rho_p}{\lambda}
=-\frac{c_2}{p^2}+\frac{c_4}{p^4}+O(p^{-6}),
\]
其中
\[
c_2=\frac{2\pi^2}{125}(6\sqrt5-5),
\qquad
c_4=\frac{2\pi^4}{9375}(7+24\sqrt5).
\]
因此
\[
\log\rho_p
=\log\lambda-\frac{c_2}{p^2}+\frac{c_4}{p^4}+O(p^{-6}).
\]
两边除以 \(\log\lambda\) 即得
\[
\Theta_p
=
1-\frac{c_2}{(\log\lambda)\,p^2}
+\frac{c_4}{(\log\lambda)\,p^4}
+O(p^{-6}),
\]
这正是盒装式。将 \(\lambda=\varphi^2\) 代入并化简，便得到后面的数值系数。最后由主项显然可知 \(\Theta_p\to 1\)。证毕。
\end{proof}

\begin{remark}[平方根门槛失效的强化]\label{rem:xi-real-input-40-collision-prime-near-perron-meaning}
上一定理强于单纯的 \(\Theta_p>\frac12\)。它表明奇素数模上的非平凡单位根扭曲并非只是在若干模数上越过平方根门槛，而是以显式 \(p^{-2}\) 速率贴近未扭曲 Perron 指数 \(1\)；换言之，这族扭曲通道在大素数极限中几乎保留了全部主模态指数。
\end{remark}

\begin{corollary}[素数模上的平方根门槛呈余有限失效]\label{cor:xi-real-input-40-collision-prime-cofinite-sqrt-failure}
沿用定理 \ref{thm:xi-real-input-40-collision-prime-near-perron-rigidity} 与
\ref{thm:xi-real-input-40-collision-mod2-trace-three-spectrum} 的记号。则存在 \(p_0\)，使得对所有奇素数 \(p>p_0\) 都有
\[
\Theta_p>\frac12.
\]
另一方面，模 \(2\) 扭曲指数
\[
\theta_2:=\frac{\log\rho_-}{\log\lambda}
\]
亦满足
\[
\theta_2>\frac12.
\]
因此，在素数模集合中，满足 RH 形平方根门槛的 Dirichlet 扭曲误差至多只剩有限个例外；等价地，平方根门槛在素数模方向呈余有限失效。
\end{corollary}
\begin{proof}
由定理 \ref{thm:xi-real-input-40-collision-prime-near-perron-rigidity}，\(\Theta_p\to 1\)，故对充分大的奇素数 \(p\) 必有 \(\Theta_p>\frac12\)。

对模 \(2\) 的情形，由定理 \ref{thm:xi-real-input-40-collision-mod2-trace-three-spectrum}，\(\rho_-\) 是三次方程
\[
x^3-2x^2-1=0
\]
的最大实根。又该多项式在 \(x=2\) 处取值为 \(-1\)，在 \(x=3\) 处取值为 \(8\)，故
\[
\rho_->2>\varphi=\lambda^{1/2}.
\]
于是
\[
\theta_2=\frac{\log\rho_-}{\log\lambda}>\frac12.
\]
综上，素数模上的平方根门槛失效只可能保留有限个例外。证毕。
\end{proof}

\input{subsubsec__xi-time-protocol-conclusions-part58a-dirichlet-minimal-witness-worst-exponent}
\input{subsubsec__xi-time-protocol-conclusions-part58b-real-input40-congruence-ldp-synthesis}

\input{subsubsec__xi-time-protocol-conclusions-part59-bernoulli-zero-natural-extension-resonance}
\input{subsubsec__xi-time-protocol-conclusions-part59a-resonance-closure-pisot-bernoulli}

\endinput


\endinput

\paragraph{全微态阈值的极端矩恢复、冻结压力分解与零温地基层分裂}
在全微态精确反演的有限层比特阈值、冻结相的双指数分离、原子 Euler 因子的 primitive 剥离以及零温四态地基层闭式已经建立之后，还可进一步抽取出下列六条跨模块后果。

\begin{corollary}[全可逆阈值的极端矩重建]\label{cor:xi-full-inversion-threshold-extreme-moment-recovery}
沿用定理 \ref{thm:xi-full-microstate-exact-inversion-bitrate-threshold} 的记号。对每个固定分辨率 \(m\) 与每个 \(q>0\)，记
\[
S_q(m):=\sum_{x\in X_m}d_m(x)^q.
\]
则有
\[
\boxed{\;
M_m=\lim_{q\to\infty}S_q(m)^{1/q}.\;}
\]
从而
\[
\boxed{\;
B_m^{\mathrm{full}}
=
\left\lceil
\log_2\!\left(\lim_{q\to\infty}S_q(m)^{1/q}\right)
\right\rceil.\;}
\]
因此，全微态精确反演阈值不仅由最大纤维规模 \(M_m\) 唯一确定，也由正阶幂和矩谱 \(\{S_q(m)\}_{q>0}\) 的端点极限唯一恢复。
\end{corollary}
\begin{proof}
对固定 \(m\)，向量
\[
\mathbf d^{(m)}:=\bigl(d_m(x)\bigr)_{x\in X_m}
\]
只有有限个坐标。有限维范数极限定理给出
\[
\lim_{q\to\infty}\left(\sum_{x\in X_m}d_m(x)^q\right)^{1/q}
=
\max_{x\in X_m}d_m(x)
=
M_m.
\]
这便得到第一式。再将其代入定理 \ref{thm:xi-full-microstate-exact-inversion-bitrate-threshold} 的恒等式
\[
B_m^{\mathrm{full}}=\left\lceil \log_2 M_m\right\rceil
\]
即可得到第二式。证毕。
\end{proof}

\begin{theorem}[冻结相压力的地址成本--残余简并分解]\label{thm:xi-fixed-freezing-pressure-address-cost-residual-decomposition}
沿用定理 \ref{thm:xi-full-microstate-exact-inversion-bitrate-threshold}、定理 \ref{thm:fold-pressure-freezing-threshold} 与定理 \ref{thm:xi-fixed-freezing-escort-renyi-spectrum-collapse} 的记号。对任意 \(a>a_{\mathrm c}\)，记
\[
\beta_{\mathrm{inv}}
:=
\lim_{m\to\infty}\frac{B_m^{\mathrm{full}}}{m}
=
\frac{\alpha_*}{\log 2},
\]
并记冻结相的 min-entropy 率为
\[
h_\infty(a)
:=
\lim_{m\to\infty}\frac{1}{m}
\Bigl(-\log\max_{x\in X_m}\pi_{m,a}(x)\Bigr)
=
g_*.
\]
则冻结相压力满足严格分解
\[
\boxed{\;
P_a
=
a(\log 2)\,\beta_{\mathrm{inv}}+h_\infty(a).\;}
\]
换言之，冻结相压力恰分解为“实现全微态精确反演所需的指数地址成本”与“冻结端面的残余简并熵率”两部分。
\end{theorem}
\begin{proof}
由定理 \ref{thm:fold-pressure-freezing-threshold}，对任意 \(a>a_{\mathrm c}\) 都有
\[
P_a=a\alpha_*+g_*.
\]
另一方面，定理 \ref{thm:xi-full-microstate-exact-inversion-bitrate-threshold} 给出
\[
\beta_{\mathrm{inv}}=\frac{\alpha_*}{\log 2},
\]
而定理 \ref{thm:xi-fixed-freezing-escort-renyi-spectrum-collapse} 在 \(q=\infty\) 时给出
\[
h_\infty(a)=g_*.
\]
将这两式代回即可得到
\[
P_a
=
a\alpha_*+g_*
=
a(\log 2)\,\beta_{\mathrm{inv}}+h_\infty(a).
\]
证毕。
\end{proof}
\leanverified{paper\_xi\_fixed\_freezing\_pressure\_address\_cost\_residual\_decomposition}

\begin{corollary}[冻结相的斜率--截距反演公式]\label{cor:xi-fixed-freezing-slope-intercept-inversion}
沿用定理 \ref{thm:xi-fixed-freezing-pressure-address-cost-residual-decomposition} 的记号。则有
\[
\boxed{\;
\lim_{q\to\infty}\frac{P_q}{q}=\alpha_*,
\qquad
\beta_{\mathrm{inv}}
=
\frac{1}{\log 2}\lim_{q\to\infty}\frac{P_q}{q}.\;}
\]
并且对任意 \(a>a_{\mathrm c}\)，成立
\[
\boxed{\;
h_\infty(a)
=
P_a-a\lim_{q\to\infty}\frac{P_q}{q}.\;}
\]
因此，冻结相的截距并非独立参数，而是由压力谱减去其无穷远斜率后精确恢复。
\end{corollary}
\begin{proof}
由定理 \ref{thm:xi-time-part57b-pressure-spectrum-slope-intercept-recovery}，
\[
\lim_{q\to\infty}\frac{P_q}{q}=\alpha_*.
\]
再与
\[
\beta_{\mathrm{inv}}=\frac{\alpha_*}{\log 2}
\]
合并，便得第二式。

另一方面，定理
\ref{thm:xi-fixed-freezing-pressure-address-cost-residual-decomposition}
给出
\[
P_a=a\alpha_*+h_\infty(a).
\]
将 \(\alpha_*\) 以
\[
\lim_{q\to\infty}\frac{P_q}{q}
\]
替换，即得
\[
h_\infty(a)
=
P_a-a\lim_{q\to\infty}\frac{P_q}{q}.
\]
证毕。
\end{proof}

\begin{theorem}[冻结相中高阶幂和比值的指数锁定]\label{thm:xi-fixed-freezing-power-sum-ratio-locking}
在严格基态间隙
\[
\alpha_*^{(2)}<\alpha_*
\]
下，沿用定理 \ref{thm:xi-fixed-freezing-escort-renyi-spectrum-collapse} 的记号。记
\[
M_m:=\max_{x\in X_m}d_m(x),
\qquad
K_m:=\#\{x\in X_m:\ d_m(x)=M_m\},
\]
\[
M_m^{(2)}:=\max\{d_m(x):\ d_m(x)<M_m\},
\qquad
S_t(m):=\sum_{x\in X_m}d_m(x)^t.
\]
若次大纤维不存在，则约定 \(M_m^{(2)}:=1\)。对任意固定整数 \(k\ge 1\) 与任意
\[
a>a_{\mathrm c}:=\frac{\log\varphi-g_*}{\alpha_*-\alpha_*^{(2)}},
\]
都有
\[
\boxed{\;
\frac{S_{a+k}(m)}{S_a(m)}
=
M_m^k\Bigl(1+O\bigl(e^{-m\Delta(a)+o(m)}\bigr)\Bigr),\;}
\]
其中
\[
\Delta(a):=a(\alpha_*-\alpha_*^{(2)})-(\log\varphi-g_*)>0.
\]
\end{theorem}
\begin{proof}
把
\[
R_t(m):=\sum_{x:\,d_m(x)<M_m}d_m(x)^t
\]
记作基态外余项，则对任意 \(t>0\) 都有精确分解
\[
S_t(m)=K_mM_m^t+R_t(m).
\]
并且
\[
0\le R_t(m)\le |X_m|\,(M_m^{(2)})^t.
\]
由
\[
M_m=e^{\alpha_* m+o(m)},
\qquad
K_m=e^{g_* m+o(m)},
\qquad
|X_m|=e^{(\log\varphi)m+o(m)},
\]
以及
\[
\limsup_{m\to\infty}\frac1m\log M_m^{(2)}=\alpha_*^{(2)},
\]
可得
\[
\frac{R_a(m)}{K_mM_m^a}
\le
\exp\!\bigl(-m\Delta(a)+o(m)\bigr).
\]
因此
\[
S_a(m)=K_mM_m^a\Bigl(1+O\bigl(e^{-m\Delta(a)+o(m)}\bigr)\Bigr).
\]
同理，
\[
\frac{R_{a+k}(m)}{K_mM_m^{a+k}}
\le
\exp\!\bigl(-m\Delta(a+k)+o(m)\bigr)
\le
\exp\!\bigl(-m\Delta(a)+o(m)\bigr),
\]
因为
\[
\Delta(a+k)=\Delta(a)+k(\alpha_*-\alpha_*^{(2)})\ge \Delta(a).
\]
故
\[
S_{a+k}(m)=K_mM_m^{a+k}\Bigl(1+O\bigl(e^{-m\Delta(a)+o(m)}\bigr)\Bigr).
\]
两式相除便得到
\[
\frac{S_{a+k}(m)}{S_a(m)}
=
M_m^k\,
\frac{1+O\bigl(e^{-m\Delta(a)+o(m)}\bigr)}
{1+O\bigl(e^{-m\Delta(a)+o(m)}\bigr)}
=
M_m^k\Bigl(1+O\bigl(e^{-m\Delta(a)+o(m)}\bigr)\Bigr).
\]
证毕。
\end{proof}

\begin{corollary}[冻结 escort 测度的正幂矩塌缩]\label{cor:xi-fixed-freezing-escort-positive-moment-collapse}
沿用定理 \ref{thm:xi-fixed-freezing-power-sum-ratio-locking} 的记号。对任意固定整数 \(k\ge 1\)，若
\[
X\sim \pi_{m,a},
\qquad
\pi_{m,a}(x):=\frac{d_m(x)^a}{S_a(m)},
\qquad
a>a_{\mathrm c},
\]
则有
\[
\boxed{\;
\EE_{\pi_{m,a}}\!\left[\left(\frac{d_m(X)}{M_m}\right)^k\right]
=
1+O\bigl(e^{-m\Delta(a)+o(m)}\bigr),\;}
\]
并且
\[
\boxed{\;
\EE_{\pi_{m,a}}\!\left[\left(1-\frac{d_m(X)}{M_m}\right)^k\right]
=
O\bigl(e^{-m\Delta(a)+o(m)}\bigr).\;}
\]
\end{corollary}
\begin{proof}
由定义直接得到
\[
\EE_{\pi_{m,a}}\!\left[\left(\frac{d_m(X)}{M_m}\right)^k\right]
=
\frac{S_{a+k}(m)}{M_m^kS_a(m)}.
\]
代入定理 \ref{thm:xi-fixed-freezing-power-sum-ratio-locking} 即得第一式。

对第二式，作二项展开
\[
\left(1-\frac{d_m(X)}{M_m}\right)^k
=
\sum_{j=0}^{k}(-1)^j\binom{k}{j}\left(\frac{d_m(X)}{M_m}\right)^j.
\]
取 \(\pi_{m,a}\)-期望后，\(j=0\) 项恒等于 \(1\)，而对每个 \(1\le j\le k\) 由第一式有
\[
\EE_{\pi_{m,a}}\!\left[\left(\frac{d_m(X)}{M_m}\right)^j\right]
=
1+O\bigl(e^{-m\Delta(a)+o(m)}\bigr).
\]
于是
\[
\EE_{\pi_{m,a}}\!\left[\left(1-\frac{d_m(X)}{M_m}\right)^k\right]
=
\sum_{j=0}^{k}(-1)^j\binom{k}{j}
+O\bigl(e^{-m\Delta(a)+o(m)}\bigr).
\]
再利用
\[
\sum_{j=0}^{k}(-1)^j\binom{k}{j}=0
\]
便得结论。证毕。
\end{proof}

\begin{theorem}[零温缩放极限的原子--地基精确分裂]\label{thm:xi-zero-temp-atomic-ground-exact-splitting}
沿用附录 \ref{app:real-input-40-zeta-u} 的记号。对真实输入 \(40\) 状态核，
\[
\Delta(z,u)=(1-uz^2)\,F(z,u).
\]
再令
\[
B_\infty=
\begin{pmatrix}
0&1&1&0\\
0&0&1&0\\
0&0&3&1\\
1&0&0&3
\end{pmatrix}.
\]
则对任意固定 \(w\in\CC\)，在零温缩放
\[
z=\sqrt{w/u}
\]
下有
\[
\boxed{\;
\lim_{u\to+\infty}\Delta\!\left(\sqrt{w/u},u\right)
=
(1-w)\det(I-wB_\infty).\;}
\]
等价地，零温窗口内的极限解析对象严格分裂为一个 primitive/Big-Witt 原子因子 \((1-w)\) 与一个四状态地基 SFT 的动力学行列式 \(\det(I-wB_\infty)\)。
\end{theorem}
\begin{proof}
由附录 \ref{app:real-input-40-zeta-u}，
\[
\begin{aligned}
F(z,u)
&=
1-z-6uz^2+2uz^3+(9u^2-u)z^4\\
&\qquad +(u-3u^2)z^5-(u^3+2u^2)z^6\\
&\qquad +(4u^3-3u^2)z^7+(u^3-u^4)z^8.
\end{aligned}
\]
把
\[
z=\sqrt{w/u}
\]
代入，得到
\[
\begin{aligned}
F\!\left(\sqrt{w/u},u\right)
&=
1-u^{-1/2}w^{1/2}-6w+2u^{-1/2}w^{3/2}\\
&\qquad +(9-u^{-1})w^2+(u^{-3/2}-3u^{-1/2})w^{5/2}\\
&\qquad -(1+2u^{-1})w^3+(4u^{-1/2}-3u^{-3/2})w^{7/2}\\
&\qquad +(u^{-1}-1)w^4.
\end{aligned}
\]
令 \(u\to+\infty\)，所有带负半整数次幂的项都消失，于是
\[
\lim_{u\to+\infty}F\!\left(\sqrt{w/u},u\right)
=
1-6w+9w^2-w^3-w^4.
\]
另一方面，由定理 \ref{thm:xi-zero-temperature-fourstate-artin-mazur-zeta-primitive-asymptotic}，
\[
1-6w+9w^2-w^3-w^4=\det(I-wB_\infty).
\]
再由
\[
1-u\left(\sqrt{w/u}\right)^2=1-w,
\]
便得到
\[
\lim_{u\to+\infty}\Delta\!\left(\sqrt{w/u},u\right)
=
(1-w)\det(I-wB_\infty).
\]
证毕。
\end{proof}

\begin{corollary}[零温主极点由地基层而非原子层支配]\label{cor:xi-zero-temp-leading-pole-ground-dominance}
沿用定理 \ref{thm:xi-zero-temp-atomic-ground-exact-splitting} 的记号。记
\[
y_\infty:=\rho(B_\infty)\approx 3.596154676009,
\qquad
b_\infty:=y_\infty^{-1}\approx 0.278074802141.
\]
则极限多项式
\[
(1-w)\det(I-wB_\infty)
\]
在正实轴上的最内层零点为
\[
\boxed{\;w=b_\infty,\;}
\]
而非原子零点 \(w=1\)。因此，零温缩放窗口中的主导自由能尺度完全由四状态地基层的 Perron 根 \(y_\infty\) 决定，而不由原子因子 \((1-w)\) 决定。
\end{corollary}
\begin{proof}
由定理 \ref{thm:xi-zero-temperature-fourstate-artin-mazur-zeta-primitive-asymptotic}，
\[
\det(I-wB_\infty)=1-6w+9w^2-w^3-w^4.
\]
其零点恰为 \(B_\infty\) 各特征值的倒数。由于 \(B_\infty\) 为非负不可约矩阵，其 Perron 根
\[
y_\infty=\rho(B_\infty)>1
\]
给出唯一最小正实倒数
\[
b_\infty=y_\infty^{-1}\in(0,1).
\]
另一方面，原子因子 \((1-w)\) 的正实零点为 \(w=1\)。由于
\[
b_\infty<1,
\]
在全部正实零点中更靠近原点的是 \(w=b_\infty\) 而非 \(w=1\)。故主导尺度来自地基层行列式 \(\det(I-wB_\infty)\)，而非原子因子。证毕。
\end{proof}

\begin{theorem}[零温地基四次数域与 Lee--Yang 三次数域的线性互不相交]\label{thm:xi-zero-temp-quartic-leyang-cubic-linear-disjoint}
记
\[
K_\infty:=\QQ(y_\infty),
\qquad
y_\infty=\rho(B_\infty),
\]
其中 \(y_\infty\) 满足
\[
y^4-6y^3+9y^2-y-1=0.
\]
再令
\[
K_{\mathrm{LY}}:=\QQ(u)=\QQ(b_{\mathrm{LY}}),
\]
其中 \(u\) 与 \(b_{\mathrm{LY}}\) 分别为文中 Lee--Yang 分支系数与振幅平方常数所生成的共同三次数域。则有
\[
\boxed{\;
[K_\infty:\QQ]=4,\qquad [K_{\mathrm{LY}}:\QQ]=3,\qquad
K_\infty\cap K_{\mathrm{LY}}=\QQ.\;}
\]
从而
\[
\boxed{\;
[K_\infty K_{\mathrm{LY}}:\QQ]=12,\;}
\]
并且 \(K_\infty\) 与 \(K_{\mathrm{LY}}\) 在 \(\QQ\) 上线性互不相交。
\end{theorem}
\begin{proof}
由定理 \ref{thm:xi-zero-temperature-fourstate-perron-parry-closed-form}，多项式
\[
y^4-6y^3+9y^2-y-1
\]
在 \(\QQ\) 上不可约，故
\[
[K_\infty:\QQ]=4.
\]
另一方面，由推论 \ref{cor:xi-terminal-zm-leyang-cubic-rational-observable-primitivity}，
\[
K_{\mathrm{LY}}=\QQ(u)=\QQ(b_{\mathrm{LY}})
\]
且
\[
[K_{\mathrm{LY}}:\QQ]=3.
\]

设
\[
E:=K_\infty\cap K_{\mathrm{LY}}.
\]
由塔式公式，
\[
[E:\QQ]\mid [K_\infty:\QQ]=4,
\qquad
[E:\QQ]\mid [K_{\mathrm{LY}}:\QQ]=3.
\]
因此 \([E:\QQ]\) 同时整除 \(4\) 与 \(3\)，只能等于 \(1\)。于是
\[
K_\infty\cap K_{\mathrm{LY}}=\QQ.
\]
再由塔式公式，
\[
[K_\infty K_{\mathrm{LY}}:\QQ]
=
\frac{[K_\infty:\QQ][K_{\mathrm{LY}}:\QQ]}{[K_\infty\cap K_{\mathrm{LY}}:\QQ]}
=
4\cdot 3
=
12.
\]
由于特征 \(0\) 下有限扩张皆可分，乘积次数达到乘积上界即等价于线性互不相交，故结论成立。证毕。
\end{proof}

\endinput


\endinput


\endinput
